% Transformation et recyclage des déchets — SVT 1ère A — Cours signé OnBuch+
% Programme officiel MINESEC. Compiler avec : tectonic NA07-transformation-recyclage-dechets.tex
\newcommand{\DOCMATIERE}{SVT}
\newcommand{\DOCNIVEAU}{1ère A}
\newcommand{\DOCMODULE}{Module 3 : L'Éducation à l'Environnement et au Développement Durable}
\newcommand{\DOCLECON}{NA07}
\newcommand{\DOCTITRE}{Transformation et recyclage des déchets}
% OnBuch+ — préambule « cours signés » ENRICHI (compilé avec tectonic / XeLaTeX)
% Système visuel « Pop » d'OnBuch+ : fond crème (jamais blanc), bordures encre
% épaisses, ombres « dures » décalées, titres Archivo Black en capitales.
% Version MATHÉMATIQUES : graphes (pgfplots/tikz), boîtes pédagogiques
% (définition, propriété, méthode, attention, le savais-tu, exercice résolu,
% exercice, TP, bilan), page de garde et signature OnBuch+ ancrée en bas.
%
% Macros attendues dans main.tex AVANT \input{preamble} :
%   \DOCMATIERE  \DOCNIVEAU  \DOCMODULE  \DOCLECON (numéro)  \DOCTITRE
\documentclass[11pt]{article}

\usepackage{fontspec}
\usepackage[french]{babel}
\usepackage[a4paper,margin=2cm,top=3.4cm,bottom=2.8cm,headheight=1.4cm,footskip=1.3cm]{geometry}
\usepackage{amsmath,amssymb,amsfonts}
\usepackage{mathtools} % \xmapsto, \coloneqq... souvent utilisés par les modèles
\usepackage{cancel} % simplifications barrées (bilans de forces)
% \abs{z} (module, valeur absolue) et \norm{u} (norme), utilisés par les modèles
\providecommand{\abs}{}\let\abs\relax\DeclarePairedDelimiter{\abs}{\lvert}{\rvert}
\providecommand{\norm}{}\let\norm\relax\DeclarePairedDelimiter{\norm}{\lVert}{\rVert}
\usepackage[table]{xcolor} % [table] : fournit aussi \rowcolors (alternance de lignes)
\usepackage{pagecolor}
\usepackage{tikz}
\usetikzlibrary{arrows.meta,positioning,calc,shapes.geometric,shapes.misc,shapes.arrows,decorations.pathmorphing,decorations.pathreplacing,decorations.markings,patterns,shadows,mindmap,trees,fit,backgrounds,matrix,intersections,angles,quotes}
\usepackage{pgfplots}
\usepackage[version=4]{mhchem} % équations nucléaires \ce{^{235}_{92}U}
\usepackage[european,siunitx]{circuitikz} % schémas électriques
\pgfplotsset{compat=1.18}
\usepgfplotslibrary{fillbetween,groupplots} % aires entre courbes (calcul intégral)
\usepackage{enumitem}
\usepackage{fancyhdr}
% [most] charge toutes les bibliothèques tcolorbox (listings, minted, raster,
% documentation...) et gonfle inutilement la mémoire TeX sur un document long
% (~300 boîtes) : on ne charge que ce qui est réellement utilisé.
\usepackage[skins,breakable]{tcolorbox}
\usepackage{graphicx}
\usepackage{colortbl}
\usepackage{booktabs}
\usepackage{tabularx}
\usepackage{diagbox} % case d'en-tête diagonale des tableaux à double entrée
% Colonnes extensibles centrée / gauche / droite (C, L, R), souvent utilisées par les modèles
\newcolumntype{C}{>{\centering\arraybackslash}X}
\newcolumntype{L}{>{\raggedright\arraybackslash}X}
\newcolumntype{R}{>{\raggedleft\arraybackslash}X}
\usepackage{array}
\usepackage{multicol}
\usepackage{siunitx}
% parse-numbers=false : accepte aussi \SI{2,5 \times 10^{-3}}{...}
\sisetup{locale=FR,output-decimal-marker={,},group-minimum-digits=4,parse-numbers=false}
\DeclareSIUnit\ohm{\ensuremath{\symup{\Omega}}} % U+2126 absent de la police
\DeclareSIUnit\division{div} % réglages d'oscilloscope (ms/div, V/div)
\DeclareSIUnit\tr{tr} % tours (vitesse de rotation, ex. \SI{1500}{\tr\per\minute})
\DeclareSIUnit\ch{ch} % cheval-vapeur (puissance moteur, unité usuelle hors SI)
\DeclareSIUnit\franccfa{FCFA} % franc CFA (coûts, contexte camerounais)
\DeclareSIUnit\dioptre{\ensuremath{\delta}} % vergence d'une lentille (dioptrie)
\usepackage{qrcode}
% \degree : commande fréquemment utilisée par les modèles pour un angle ou une
% température en dehors de \SI{}{} (non fournie par les paquets chargés).
\providecommand{\degree}{^{\circ}}
% \qed (fin de démonstration), utilisé par les modèles sans amsthm
\providecommand{\qed}{\hfill$\square$}
% \barre : double barre des valeurs interdites dans un tableau de variations (tkz-tab)
\providecommand{\barre}{\ensuremath{\|}}
% environnement proof (amsthm non chargé) utilisé par les modèles
\ifdefined\proof\else\newenvironment{proof}[1][Démonstration]{\par\noindent\textit{#1.}\ }{\qed\par}\fi
\usepackage{lastpage}
\usepackage{hyperref}
\hypersetup{hidelinks}

% ============================================================
% Palette « Pop » OnBuch+ — mêmes hex que la classe P (plus_ui.dart)
% ============================================================
\definecolor{poporange}{HTML}{F59321}
\definecolor{popdark}{HTML}{E07A0C}
\definecolor{popcream}{HTML}{FFF6E8}
\definecolor{popink}{HTML}{1C1714}
\definecolor{poppurple}{HTML}{7A5AE0}
\definecolor{popgreen}{HTML}{2E9E6A}
\definecolor{poppink}{HTML}{E0457B}
\definecolor{popblue}{HTML}{2F7DE1}
\definecolor{popgold}{HTML}{C98A12}
\definecolor{popmuted}{HTML}{8A7F76}
\definecolor{popline}{HTML}{EADFD2}
\definecolor{popgray}{HTML}{8A7F76}
\colorlet{popgrey}{popgray}
% Alias tolérants (le modèle nomme parfois les couleurs différemment)
\colorlet{popwhite}{white}
\colorlet{popblack}{popink}
\colorlet{popred}{poppink}
\colorlet{popyellow}{popgold}
% Teintes claires pour fonds de tableaux / zones
\colorlet{popblueL}{popblue!10!white}
\colorlet{poporangeL}{poporange!14!white}
\colorlet{popgreenL}{popgreen!12!white}
\colorlet{poppurpleL}{poppurple!12!white}
\colorlet{poppinkL}{poppink!12!white}
\colorlet{popgoldL}{popgold!14!white}

\pagecolor{popcream}

% ============================================================
% Typographie
% ============================================================
\defaultfontfeatures{Ligatures=TeX}
\setmainfont{PlusJakartaSans-Medium.ttf}[Path=../../../fonts/,BoldFont=PlusJakartaSans-Bold.ttf]
% Maths en sans-serif (Fira Math) avec chiffres et lettres droites de Plus
% Jakarta, pour que les formules se fondent dans le texte.
\usepackage[math-style=ISO,bold-style=ISO]{unicode-math}
\setmathfont{FiraMath-Regular.otf}[Path=../../../fonts/]
\setmathfont{PlusJakartaSans-Medium.ttf}[Path=../../../fonts/,range={up/num,up/{Latin,latin},\mathup}]
\usepackage{icomma} % virgule décimale sans espace en mode math
% Caractères mathématiques Unicode absents des polices de texte (Plus Jakarta,
% Archivo Black) : rendus par leur équivalent mathématique, en texte comme en
% formule — sinon ils s'affichent en carré vide (ex. « Arithmétique dans ℤ »).
\usepackage{newunicodechar}
\newunicodechar{ℝ}{\ensuremath{\mathbb{R}}}
\newunicodechar{ℤ}{\ensuremath{\mathbb{Z}}}
\newunicodechar{ℂ}{\ensuremath{\mathbb{C}}}
\newunicodechar{ℕ}{\ensuremath{\mathbb{N}}}
\newunicodechar{ℚ}{\ensuremath{\mathbb{Q}}}
\newunicodechar{𝔻}{\ensuremath{\mathbb{D}}}
\newunicodechar{α}{\ensuremath{\alpha}}
\newunicodechar{β}{\ensuremath{\beta}}
\newunicodechar{γ}{\ensuremath{\gamma}}
\newunicodechar{δ}{\ensuremath{\delta}}
\newunicodechar{ε}{\ensuremath{\varepsilon}}
\newunicodechar{θ}{\ensuremath{\theta}}
\newunicodechar{λ}{\ensuremath{\lambda}}
\newunicodechar{μ}{\ensuremath{\mu}}
\newunicodechar{σ}{\ensuremath{\sigma}}
\newunicodechar{τ}{\ensuremath{\tau}}
\newunicodechar{φ}{\ensuremath{\varphi}}
\newunicodechar{ψ}{\ensuremath{\psi}}
\newunicodechar{ω}{\ensuremath{\omega}}
\newunicodechar{Σ}{\ensuremath{\Sigma}}
% majuscules produites par \MakeUppercase dans les titres (α-aminés -> Α)
\newunicodechar{Α}{\ensuremath{\alpha}}
\newunicodechar{Β}{\ensuremath{\beta}}
\newunicodechar{Γ}{\ensuremath{\Gamma}}
\newunicodechar{Δ}{\ensuremath{\Delta}}
\newunicodechar{Ε}{\ensuremath{\varepsilon}}
\newunicodechar{Θ}{\ensuremath{\Theta}}
\newunicodechar{Λ}{\ensuremath{\Lambda}}
\newunicodechar{Μ}{\ensuremath{\mu}}
\newunicodechar{Τ}{\ensuremath{\tau}}
\newunicodechar{Φ}{\ensuremath{\Phi}}
\newunicodechar{Ψ}{\ensuremath{\Psi}}
\newunicodechar{Ω}{\ensuremath{\Omega}}
\newunicodechar{↦}{\ensuremath{\mapsto}}
\newunicodechar{∈}{\ensuremath{\in}}
\newunicodechar{∉}{\ensuremath{\notin}}
\newunicodechar{∧}{\ensuremath{\wedge}}
\newunicodechar{⇔}{\ensuremath{\Leftrightarrow}}
\newunicodechar{⟺}{\ensuremath{\Longleftrightarrow}}
\newunicodechar{⇒}{\ensuremath{\Rightarrow}}
\newunicodechar{⟹}{\ensuremath{\Longrightarrow}}
\newunicodechar{∘}{\ensuremath{\circ}}
\newunicodechar{∪}{\ensuremath{\cup}}
\newunicodechar{∩}{\ensuremath{\cap}}
\newunicodechar{≡}{\ensuremath{\equiv}}
\newunicodechar{⊂}{\ensuremath{\subset}}
\newunicodechar{⊥}{\ensuremath{\perp}}
\newunicodechar{∃}{\ensuremath{\exists}}
\newunicodechar{∀}{\ensuremath{\forall}}
\newunicodechar{∛}{\ensuremath{\sqrt[3]{\,}}}
\newunicodechar{∜}{\ensuremath{\sqrt[4]{\,}}}
\newunicodechar{⃗}{\ensuremath{{}^{\to}}}
\newunicodechar{ⁿ}{\ifmmode^{n}\else\textsuperscript{\ensuremath{n}}\fi}
\newunicodechar{ˣ}{\ifmmode^{x}\else\textsuperscript{\ensuremath{x}}\fi}
\newunicodechar{ᵇ}{\ifmmode^{b}\else\textsuperscript{\ensuremath{b}}\fi}
\newunicodechar{ʳ}{\ifmmode^{r}\else\textsuperscript{\ensuremath{r}}\fi}
\newunicodechar{ᵘ}{\ifmmode^{u}\else\textsuperscript{\ensuremath{u}}\fi}
\newunicodechar{ʸ}{\ifmmode^{y}\else\textsuperscript{\ensuremath{y}}\fi}
\newunicodechar{ᵃ}{\ifmmode^{a}\else\textsuperscript{\ensuremath{a}}\fi}
\newunicodechar{ᵉ}{\ifmmode^{e}\else\textsuperscript{\ensuremath{e}}\fi}
\newunicodechar{ᵏ}{\ifmmode^{k}\else\textsuperscript{\ensuremath{k}}\fi}
\newunicodechar{ᵖ}{\ifmmode^{p}\else\textsuperscript{\ensuremath{p}}\fi}
\newunicodechar{ᴺ}{\ifmmode^{N}\else\textsuperscript{\ensuremath{N}}\fi}
\newunicodechar{ᶜ}{\ifmmode^{c}\else\textsuperscript{\ensuremath{c}}\fi}
\newunicodechar{ʰ}{\ifmmode^{h}\else\textsuperscript{\ensuremath{h}}\fi}
\newunicodechar{ᵠ}{\ifmmode^{q}\else\textsuperscript{\ensuremath{q}}\fi}
\newunicodechar{ₙ}{\ifmmode_{n}\else\textsubscript{\ensuremath{n}}\fi}
\newunicodechar{ₐ}{\ifmmode_{a}\else\textsubscript{\ensuremath{a}}\fi}
\newunicodechar{ᵢ}{\ifmmode_{i}\else\textsubscript{\ensuremath{i}}\fi}
\newunicodechar{ᵣ}{\ifmmode_{r}\else\textsubscript{\ensuremath{r}}\fi}
\newunicodechar{ₖ}{\ifmmode_{k}\else\textsubscript{\ensuremath{k}}\fi}
\newunicodechar{ᵦ}{\ifmmode_{\beta}\else\textsubscript{\ensuremath{\beta}}\fi}
\newunicodechar{ₓ}{\ifmmode_{x}\else\textsubscript{\ensuremath{x}}\fi}
\newfontfamily\popdisplay{ArchivoBlack-Regular.ttf}[Path=../../../fonts/]
\newfontfamily\poplogo{SpaceGrotesk-Bold.ttf}[Path=../../../fonts/]
\newfontfamily\popsemi{PlusJakartaSans-SemiBold.ttf}[Path=../../../fonts/]
\newfontfamily\popxbold{PlusJakartaSans-ExtraBold.ttf}[Path=../../../fonts/]
\newfontfamily\popxboldls{PlusJakartaSans-ExtraBold.ttf}[Path=../../../fonts/,LetterSpace=3.0]

\setlist[enumerate]{leftmargin=1.7em,itemsep=5pt,topsep=4pt}
\setlist[itemize]{leftmargin=1.7em,itemsep=4pt,topsep=4pt,label={\color{poporange}\textbullet}}
\setlength{\parindent}{0pt}
\setlength{\parskip}{5pt}
\renewcommand{\arraystretch}{1.25}

% pgfplots : style Pop par défaut
\pgfplotsset{
  every axis/.append style={axis line style={popink,line width=1pt},tick style={popink},
    grid style={popline,line width=0.5pt},label style={font=\small},tick label style={font=\footnotesize}},
  popcurve/.style={poporange,line width=1.8pt,no markers,smooth},
}
% Même style disponible dans un \draw TikZ simple (hors pgfplots)
\tikzset{popcurve/.style={poporange,line width=1.8pt}}
\tikzset{popgrid/.style={popline,very thin}}
\colorlet{popcurve}{poporange} % le nom du style est parfois utilisé comme couleur

% ============================================================
% Pill / squiggle
% ============================================================
\newcommand{\poppill}[3]{%
  \tikz[baseline=-0.55ex]\node[draw=popink,line width=1.2pt,rounded corners=2.6pt,%
    fill=#2,inner xsep=6.5pt,inner ysep=3pt]{%
    \popxboldls\fontsize{7}{7}\selectfont\color{#3}\MakeUppercase{#1}};%
}
\newcommand{\popsquiggle}[1]{%
  \tikz[baseline=-0.4ex]{
    \draw[#1,line width=2.6pt,line cap=round]
      plot [smooth] coordinates {(0,0.09) (0.34,0.01) (0.68,0.21) (1.02,0.01) (1.36,0.10)};
  }%
}
% Mot-clé surligné (marqueur orange sous le texte)
\newcommand{\cle}[1]{{\popxbold\color{popdark}#1}}

% ============================================================
% En-tête / pied
% ============================================================
\pagestyle{fancy}
\fancyhf{}
\renewcommand{\headrulewidth}{0pt}
\renewcommand{\footrulewidth}{0pt}
\lhead{\raisebox{-0.15ex}{\poplogo\fontsize{13}{13}\selectfont\color{popink}OnBuch\color{poporange}+}}
\rhead{\poppill{\DOCMATIERE\ \textbullet\ \DOCNIVEAU\ \textbullet\ Leçon \DOCLECON}{white}{popink}}
\lfoot{\popxbold\fontsize{7.6}{7.6}\selectfont\color{popmuted}\MakeUppercase{Cours signé OnBuch+}\ \textbullet\ {\popsemi\DOCTITRE}}
\rfoot{\tikz[baseline=-0.6ex]\node[rounded corners=8.5pt,fill=popink,minimum height=17pt,minimum width=17pt,inner xsep=5pt,inner sep=0]{\popxbold\fontsize{8}{8}\selectfont\color{popcream}\thepage\,/\,\pageref*{LastPage}};}

% ============================================================
% Titres
% ============================================================
\newcommand{\chaptitre}[2]{%
  \begin{tcolorbox}[enhanced,colback=poporange,colframe=popink,boxrule=2.6pt,arc=11pt,
      drop shadow=popink,left=15pt,right=15pt,top=12pt,bottom=11pt,before skip=3pt,after skip=18pt]
    \poppill{#2}{popink}{popcream}\par\vspace{9pt}
    {\raggedright\hyphenpenalty=10000\exhyphenpenalty=10000\popdisplay\fontsize{22}{24}\selectfont\color{popcream}\MakeUppercase{#1}\par}
  \end{tcolorbox}
}
% Section numérotée : pastille orange avec le numéro + titre Archivo
\newcounter{popsec}
\newcommand{\coursec}[1]{%
  \stepcounter{popsec}\setcounter{popsub}{0}%
  \par\vspace{16pt}\noindent
  \begin{minipage}{\linewidth}
  \tikz[baseline=-0.7ex]\node[draw=popink,line width=1.6pt,fill=poporange,rounded corners=4pt,minimum size=22pt,inner sep=0]{\popdisplay\fontsize{12}{12}\selectfont\color{popcream}\thepopsec};%
  \hspace{8pt}{\popdisplay\fontsize{15}{17}\selectfont\color{popink}\MakeUppercase{#1}}\par
  \vspace{4pt}\noindent{\color{poporange}\rule{36pt}{3.6pt}}
  \end{minipage}\par\nopagebreak\vspace{8pt}
}
\newcounter{popsub}
\newcommand{\courssub}[1]{%
  \stepcounter{popsub}%
  \par\vspace{9pt}\noindent{\popxbold\fontsize{11.5}{13}\selectfont\color{popink}{\color{poporange}\thepopsec.\thepopsub}\hspace{6pt}#1}\par\nopagebreak\vspace{4pt}
}

% ============================================================
% Boîtes pédagogiques — même recette : fond blanc, bordure épaisse colorée,
% ombre dure encre, badge plein. Titre optionnel [#1].
% ============================================================
\tcbset{popbase/.style={breakable,enhanced,colback=white,boxrule=2.3pt,arc=9pt,
  shadow={3pt}{-3pt}{0pt}{popink},left=14pt,right=14pt,top=11pt,bottom=11pt,before skip=11pt,after skip=15pt,
  fontupper=\popsemi\fontsize{10.6}{14.5}\selectfont\color{popink},
  fontlower=\popsemi\fontsize{10.6}{14.5}\selectfont\color{popink}}}
% Coche dessinée (le glyphe ✓ n'existe pas dans les polices du document)
\newcommand{\popcheck}[1]{\tikz[baseline=-0.5ex]\draw[#1,line width=1.6pt,line cap=round,line join=round](-0.13,0.0)--(-0.04,-0.09)--(0.13,0.1);}
\providecommand{\coche}{\popcheck{popgreen}}
\providecommand{\resultat}[1]{\par\smallskip\noindent\fcolorbox{popgreen}{popgreenL}{\parbox{\dimexpr\linewidth-2\fboxsep-2\fboxrule\relax}{\textbf{Résultat.} #1}}\par\smallskip}
\newcommand{\popboxhead}[3]{\poppill{#1}{#2}{white}\ifx\relax#3\relax\else\hspace{6pt}{\popxbold\fontsize{10.6}{12}\selectfont\color{popink}#3}\fi\par\vspace{7pt}}

\newtcolorbox{aretenir}[1][]{popbase,colframe=popgreen,before upper={\popboxhead{À retenir}{popgreen}{#1}}}
\newtcolorbox{exemplebox}[1][]{popbase,colframe=poporange,before upper={\popboxhead{Exemple}{poporange}{#1}}}
\newtcolorbox{definition}[1][]{popbase,colframe=popblue,before upper={\popboxhead{Définition}{popblue}{#1}}}
\newtcolorbox{propriete}[1][]{popbase,colframe=poppurple,before upper={\popboxhead{Propriété}{poppurple}{#1}}}
\newtcolorbox{methode}[1][]{popbase,colframe=popgold,before upper={\popboxhead{Méthode}{popgold}{#1}}}
\newtcolorbox{attention}[1][]{popbase,colframe=poppink,before upper={\popboxhead{Attention}{poppink}{#1}}}
\newtcolorbox{experience}[1][]{popbase,colframe=popblue,colback=popblueL,before upper={\popboxhead{Expérience}{popblue}{#1}}}
% « Le savais-tu ? » : carte violette pleine (contexte, culture, Cameroun)
\newtcolorbox{savaistu}[1][]{popbase,colback=poppurple,colframe=popink,
  fontupper=\popsemi\fontsize{10.4}{14.2}\selectfont\color{white},
  before upper={\poppill{Le savais-tu\,?}{white}{poppurple}\ifx\relax#1\relax\else\hspace{6pt}{\popxbold\color{white}#1}\fi\par\vspace{7pt}}}
% Exercice résolu : énoncé en haut, \tcblower puis la solution sur fond crème
\newcounter{popexo}\newcounter{popexr}
\newtcolorbox{exoresolu}[1][]{popbase,colframe=poporange,colbacklower=popcream,
  segmentation style={popink,line width=1.2pt,dashed},
  before upper={\stepcounter{popexr}\popboxhead{Exercice résolu \thepopexr}{poporange}{#1}},
  before lower={\poppill{Solution}{popink}{popcream}\par\vspace{6pt}}}
% Exercice (non résolu en place) — la correction est dans la partie Corrigés
\newtcolorbox{exercice}[1][]{popbase,colframe=popink,
  before upper={\stepcounter{popexo}\popboxhead{Exercice \thepopexo}{popink}{#1}}}
% Corrigé
\newtcolorbox{corrige}[1][]{popbase,colframe=popgreen,colback=popgreenL,
  before upper={\popboxhead{Corrigé}{popgreen}{#1}}}
% Objectifs / prérequis / bilan
\newtcolorbox{objectifs}{popbase,colframe=popink,colback=poporangeL,
  before upper={\popboxhead{Objectifs de la leçon}{popink}{}}}
\newtcolorbox{prerequis}{popbase,colframe=popink,colback=popblueL,
  before upper={\popboxhead{Prérequis}{popblue}{}}}
\newtcolorbox{bilan}{popbase,colframe=popink,colback=white,boxrule=2.8pt,
  before upper={\popboxhead{Fiche bilan}{poporange}{L'essentiel à connaître pour l'examen}}}
% Figure encadrée avec légende
\newtcolorbox{popfigure}[1][]{enhanced,breakable=false,colback=white,colframe=popline,boxrule=1.4pt,arc=8pt,
  left=10pt,right=10pt,top=10pt,bottom=8pt,before skip=10pt,after skip=12pt,halign=center,
  lower separated=false,halign lower=center,
  fontlower=\popsemi\fontsize{9}{11.5}\selectfont\color{popmuted},
  #1}
\newcommand{\legende}[1]{\par\vspace{6pt}{\popsemi\fontsize{9}{11.5}\selectfont\color{popmuted}#1}}

% ============================================================
% Signature OnBuch+
% ============================================================
\newcommand{\poplogoinline}[1]{{\poplogo\fontsize{#1}{#1}\selectfont\color{popink}OnBuch\color{poporange}+}}
\newcommand{\signatureonbuch}{%
  \begin{tcolorbox}[enhanced,colback=popink,colframe=popink,boxrule=2.2pt,arc=10pt,
      drop shadow=poporange,left=14pt,right=14pt,top=10pt,bottom=10pt,before skip=0pt,after skip=0pt]
    \tikz[baseline=-0.7ex]\node[circle,fill=popgreen,draw=popcream,line width=1.3pt,minimum size=22pt,inner sep=0]{\popcheck{white}};%
    \hspace{9pt}\begin{minipage}[c]{0.62\linewidth}
      {\popxbold\fontsize{12}{13}\selectfont\color{popcream}Signé\ \textbullet\ {\poplogo OnBuch\color{poporange}+}}\par\vspace{2pt}
      {\raggedright\popsemi\fontsize{8}{10.5}\selectfont\color{popline}\MakeUppercase{Équipe pédagogique OnBuch+}\\\MakeUppercase{Conforme au programme officiel MINESEC}\par}
    \end{minipage}\hfill
    {\poplogo\fontsize{22}{22}\selectfont\color{popcream}OnBuch\color{poporange}+}
  \end{tcolorbox}%
}

% ============================================================
% Page de garde
% #1 titre · #2 matière · #3 niveau · #4 module · #5 n° leçon · #6 durée ·
% #7 description / objectifs (texte ou itemize)
% ============================================================
\newcommand{\pagegarde}[7]{%
  \thispagestyle{empty}
  \newgeometry{a4paper,margin=1.7cm,top=1.6cm,bottom=1.4cm}%
  \begin{tikzpicture}[remember picture,overlay]
    \fill[poporange!18] ([xshift=-3.2cm,yshift=-3.6cm]current page.north east) circle (4.6cm);
    \fill[poppurple!14] ([xshift=2.4cm,yshift=7.6cm]current page.south west) circle (3.8cm);
  \end{tikzpicture}%
  {\poplogo\fontsize{30}{30}\selectfont\color{popink}OnBuch\color{poporange}+}\hfill
  \raisebox{0.6ex}{\poppill{Cours signé \textbullet\ Premium}{popink}{popcream}}\par
  \vspace{1pt}\popsquiggle{poppurple}\par
  \vspace{8pt}
  \begin{tcolorbox}[enhanced,colback=poporange,colframe=popink,boxrule=2.8pt,arc=13pt,
      drop shadow=popink,left=18pt,right=18pt,top=12pt,bottom=13pt,after skip=10pt]
    \poppill{#2 \textbullet\ #3}{popink}{popcream}\hspace{5pt}\poppill{#4}{white}{popink}\par\vspace{8pt}
    {\popdisplay\fontsize{14}{14}\selectfont\color{popink}LEÇON #5}\par\vspace{4pt}
    {\raggedright\hyphenpenalty=10000\exhyphenpenalty=10000\popdisplay\fontsize{25}{27}\selectfont\color{popcream}\MakeUppercase{#1}\par}
    \vspace{8pt}
    {\popxbold\fontsize{9.5}{9.5}\selectfont\color{popink}\MakeUppercase{Durée conseillée : #6}}
  \end{tcolorbox}
  \begin{tcolorbox}[enhanced,colback=white,colframe=popink,boxrule=2.2pt,arc=10pt,
      drop shadow=popink,left=16pt,right=16pt,top=10pt,bottom=10pt,after skip=8pt]
    \poppill{Dans cette leçon}{poppurple}{white}\par\vspace{5pt}
    \popsemi\fontsize{9.3}{12.6}\selectfont\color{popink}#7
  \end{tcolorbox}
  \noindent\begin{minipage}[t]{0.487\linewidth}
    \begin{tcolorbox}[enhanced,colback=popgreen,colframe=popink,boxrule=2.2pt,arc=10pt,
        drop shadow=popink,left=11pt,right=11pt,top=10pt,bottom=10pt,height=3.1cm,valign=center]
      \tcbox[colback=white,colframe=popink,boxrule=1.3pt,arc=5pt,boxsep=0pt,nobeforeafter,
        left=3pt,right=3pt,top=3pt,bottom=3pt,valign=center]{\qrcode[height=1.9cm]{https://play.google.com/store/apps/details?id=com.luvvix.onbuch&hl=fr}}%
      \hspace{8pt}\begin{minipage}[c]{0.5\linewidth}
        {\popdisplay\fontsize{11}{12.5}\selectfont\color{white}\MakeUppercase{Télécharge OnBuch}}\par\vspace{4pt}
        {\raggedright\popsemi\fontsize{8.4}{11}\selectfont\color{white}Gratuit \textbullet\ Google Play\\Tuteur IA, annales, concours, résultats\par}
      \end{minipage}
    \end{tcolorbox}
  \end{minipage}\hfill
  \begin{minipage}[t]{0.487\linewidth}
    \begin{tcolorbox}[enhanced,colback=poppurple,colframe=popink,boxrule=2.2pt,arc=10pt,
        drop shadow=popink,left=11pt,right=11pt,top=10pt,bottom=10pt,height=3.1cm,valign=center]
      \tikz[baseline=-0.7ex]\node[circle,fill=white,draw=popink,line width=1.3pt,minimum size=1.6cm,inner sep=0]{\popdisplay\fontsize{24}{24}\selectfont\color{poppurple}+};%
      \hspace{8pt}\begin{minipage}[c]{0.6\linewidth}
        {\popdisplay\fontsize{11}{12.5}\selectfont\color{white}\MakeUppercase{Passe à OnBuch+}}\par\vspace{4pt}
        {\raggedright\popsemi\fontsize{8.4}{11}\selectfont\color{white}Tous les cours signés, Tuteur IA illimité, fiches PDF — onglet OnBuch+ de l'app\par}
      \end{minipage}
    \end{tcolorbox}
  \end{minipage}
  \vspace{8pt}
  \popcontenu
  \vfill
  \signatureonbuch
  \restoregeometry
  \newpage
}

% Rangée « Ce cours contient » (page de garde)
\newcommand{\poptile}[3]{% #1 couleur #2 gros texte #3 légende
  \begin{tcolorbox}[enhanced,colback=#1,colframe=popink,boxrule=2pt,arc=9pt,shadow={2.5pt}{-2.5pt}{0pt}{popink},
      left=6pt,right=6pt,top=9pt,bottom=9pt,halign=center,valign=center,height=1.75cm,width=\linewidth,nobeforeafter]
    {\popdisplay\fontsize{12.5}{13}\selectfont\color{white}\MakeUppercase{#2}}\par\vspace{3pt}
    {\centering\popsemi\fontsize{7.8}{9.5}\selectfont\color{white}#3\par}
  \end{tcolorbox}}
\newcommand{\popcontenu}{%
  \par\noindent{\popxboldls\fontsize{7.5}{7.5}\selectfont\color{popmuted}CE COURS SIGNÉ CONTIENT}\par\vspace{7pt}
  \noindent\begin{minipage}[t]{0.235\linewidth}\poptile{popblue}{Cours}{complet, illustré, pas à pas}\end{minipage}\hfill
  \begin{minipage}[t]{0.235\linewidth}\poptile{poporange}{Méthodes}{et exercices résolus}\end{minipage}\hfill
  \begin{minipage}[t]{0.235\linewidth}\poptile{poppink}{Exercices}{type examen + corrigés}\end{minipage}\hfill
  \begin{minipage}[t]{0.235\linewidth}\poptile{popgreen}{Fiche bilan}{l'essentiel à retenir}\end{minipage}\par}

% Sommaire
\newtcolorbox{sommaire}{enhanced,colback=white,colframe=popink,boxrule=2.2pt,arc=10pt,
  drop shadow=popink,left=18pt,right=18pt,top=13pt,bottom=13pt,before skip=6pt,after skip=20pt,
  before upper={\poppill{Sommaire}{poppurple}{white}\par\vspace{9pt}\popsemi\fontsize{10.6}{14.5}\selectfont\color{popink}}}

% Signature de fin de cours — poussée tout en bas de la dernière page
\newcommand{\findecours}{%
  \par\vfill\nopagebreak
  \begin{center}\popsquiggle{poporange}\end{center}\vspace{4pt}
  \signatureonbuch
}
\AtBeginDocument{\def\llbracket{\mathopen{[\mkern-3.5mu[}}\def\rrbracket{\mathclose{]\mkern-3.5mu]}}} % intervalles d'entiers (glyphe ⟦ absent de la police)
% Glyphes absents de Fira Math : redéfinis à partir de glyphes disponibles
\AtBeginDocument{%
  \def\setminus{\mathbin{\backslash}}\let\smallsetminus\setminus
  \def\vdots{\mathord{\vbox{\baselineskip3.2pt\lineskiplimit0pt\kern4pt\hbox{.}\hbox{.}\hbox{.}}}}%
  \def\ddots{\mathinner{\mkern1mu\raise6.5pt\hbox{.}\mkern2mu\raise3.6pt\hbox{.}\mkern2mu\raise0.7pt\hbox{.}\mkern1mu}}%
  \def\triangle{\mathord{\scalebox{0.9}{$\symup{\Delta}$}}}%
  \def\nabla{\mathord{\raisebox{\depth}{\scalebox{1}[-1]{$\symup{\Delta}$}}}}%
  \def\checkmark{\popcheck{popdark}}%
  \def\star{\mathbin{\ast}}\def\bigstar{\mathord{\ast}}%
  \def\diamond{\mathbin{\scalebox{0.6}{$\Diamond$}}}%
  \def\bigcup{\mathop{\mathchoice{\raisebox{-0.3em}{\scalebox{1.7}{$\cup$}}}{\scalebox{1.25}{$\cup$}}{\cup}{\cup}}\displaylimits}%
  \def\bigcap{\mathop{\mathchoice{\raisebox{-0.3em}{\scalebox{1.7}{$\cap$}}}{\scalebox{1.25}{$\cap$}}{\cap}{\cap}}\displaylimits}%
  \def\lesseqqgtr{\mathrel{\lessgtr}}\def\bigoplus{\mathop{\oplus}\displaylimits}%
}
\providecommand{\ssi}{si et seulement si} % abréviation fréquente
\AtBeginDocument{\providecommand{\wideparen}{\overparen}} % arc de cercle

%% FIGFIX-BEGIN v3 — correctifs globaux des figures (docs/onbuch-generation/tools/figfix)
\usepackage{adjustbox}\usepackage{etoolbox}
% toute figure TikZ/pgfplots plus large que la ligne est réduite à la largeur disponible
\BeforeBeginEnvironment{tikzpicture}{\begin{adjustbox}{max width=\linewidth}}
\AfterEndEnvironment{tikzpicture}{\end{adjustbox}}
% légendes pgfplots lisibles par-dessus les courbes
\pgfplotsset{every axis/.append style={legend style={fill=white,fill opacity=0.92,text opacity=1,draw=black!15,inner sep=3pt}}}
% étiquettes d'arêtes (syntaxe edge["..."]) avec fond blanc
\tikzset{every edge quotes/.append style={fill=white,inner sep=1.2pt}}
% bulles de cartes mentales : texte à la ligne dans la bulle, bulle un peu plus grande
\tikzset{every concept/.append style={text width=2.0cm,align=center,minimum size=2.3cm,inner sep=2pt}}
%% FIGFIX-END
\begin{document}

\pagegarde{Transformation et recyclage des déchets}{SVT}{1ère A}{Module 3 : L'Éducation à l'Environnement et au Développement Durable}{NA07}{4 h}{Cette leçon présente les différents types de déchets, leurs impacts sur l'environnement et la santé, ainsi que les principes du tri, de la transformation et du recyclage. Elle s'appuie sur des exemples camerounais de valorisation (compostage, biogaz, recyclage du plastique) pour former des citoyens acteurs d'une gestion durable des déchets.
\vspace{4pt}
\begin{multicols}{2}\small
\begin{itemize}
\item À la fin de cette leçon, je sais définir ce qu'est un déchet et citer les principales sources de déchets domestiques et urbains.
\item À la fin de cette leçon, je sais classer les déchets en catégories (organiques, plastiques, dangereux, inertes) selon leur nature.
\item À la fin de cette leçon, je sais expliquer les impacts environnementaux et sanitaires des déchets mal gérés.
\item À la fin de cette leçon, je sais décrire les principes du tri des déchets et la règle des 3R (Réduire, Réutiliser, Recycler).
\item À la fin de cette leçon, je sais proposer un mode de transformation ou de recyclage adapté à un déchet donné.
\item À la fin de cette leçon, je sais expliquer le fonctionnement du compostage et de la production de biogaz.
\end{itemize}\end{multicols}}

\begin{sommaire}
\begin{multicols}{2}
\begin{enumerate}
\item Les déchets : définition, sources et typologie
\item Impacts des déchets sur l'environnement et la santé
\item Le tri des déchets et la règle des 3R
\item Transformation et recyclage : principes et filières
\item Valorisation des déchets organiques : compostage et biogaz
\item Agir pour une gestion responsable des déchets
\item Activité d'intégration
\item Méthodes et exercices résolus
\item Exercices
\item Corrigés des exercices
\item Fiche bilan
\end{enumerate}
\end{multicols}
\end{sommaire}

\begin{objectifs}
\begin{itemize}
\item À la fin de cette leçon, je sais définir ce qu'est un déchet et citer les principales sources de déchets domestiques et urbains.
\item À la fin de cette leçon, je sais classer les déchets en catégories (organiques, plastiques, dangereux, inertes) selon leur nature.
\item À la fin de cette leçon, je sais expliquer les impacts environnementaux et sanitaires des déchets mal gérés.
\item À la fin de cette leçon, je sais décrire les principes du tri des déchets et la règle des 3R (Réduire, Réutiliser, Recycler).
\item À la fin de cette leçon, je sais proposer un mode de transformation ou de recyclage adapté à un déchet donné.
\item À la fin de cette leçon, je sais expliquer le fonctionnement du compostage et de la production de biogaz.
\item À la fin de cette leçon, je sais citer des exemples locaux de valorisation des déchets au Cameroun.
\item À la fin de cette leçon, je sais concevoir et mener une action de sensibilisation à la gestion responsable des déchets dans mon environnement.
\end{itemize}
\end{objectifs}

\begin{prerequis}
\begin{itemize}
\item Notions d'écosystème et de chaîne alimentaire (classes de collège)
\item Notion de pollution de l'eau, de l'air et du sol (Seconde / début du Module 3)
\item Notion de développement durable et ses trois piliers (début du Module 3)
\item Connaissance du quotidien : gestion des ordures dans son quartier ou son village
\item Notions élémentaires de fermentation et de décomposition par les micro-organismes
\end{itemize}
\end{prerequis}

\newpage

\coursec{Les déchets : définition, sources et typologie}

\textbf{Situation d'accroche.} Chaque jour, la ville de Douala produit plus de \num{3 000} tonnes de déchets, dont une grande partie finit dans les rues, les caniveaux et les lagunes, provoquant inondations et maladies. Pourtant, à Bafoussam, une coopérative de femmes transforme les épluchures de manioc en compost vendu aux maraîchers, et à Yaoundé, des jeunes fabriquent des pavés à partir de sachets plastiques fondus. Comment ces déchets qui nous encombrent peuvent-ils devenir une richesse ? Quels déchets peut-on trier, transformer ou recycler dans notre quartier ? Cette leçon nous donne les clés pour passer du problème à la solution. Mais avant de transformer, il faut d'abord comprendre ce qu'est un déchet, d'où il vient et comment on le classe.

\courssub{Qu'est-ce qu'un déchet ?}

\textbf{Observation.} À la fin de la journée au marché Mokolo à Yaoundé, on ramasse des épluchures de banane plantain, des feuilles de chou abîmées, des sachets plastiques déchirés, des boîtes de conserve vides et des piles usagées. Tous ces objets ont un point commun : leurs propriétaires s'en sont débarrassés parce qu'ils ne leur étaient plus utiles.

\begin{definition}[Déchet]
Un \cle{déchet} est toute substance ou tout objet dont le détenteur se défait ou dont il a l'intention de se défaire, parce qu'il ne lui est plus utile ou qu'il est devenu encombrant.
\end{definition}

\begin{attention}[Un déchet n'est pas forcément inutile]
Un objet peut être un déchet pour une personne et une ressource pour une autre : les épluchures de manioc sont un déchet pour la cuisinière, mais une matière première pour le composteur. C'est toute la logique de la valorisation que nous étudierons dans cette leçon.
\end{attention}

\courssub{Les sources de déchets}

Les déchets proviennent de toutes les activités humaines. On distingue les principales sources suivantes :

\begin{itemize}
\item les \cle{ménages} : déchets domestiques (restes de repas, emballages, vieux vêtements, appareils usagés) ;
\item les \cle{marchés} : épluchures, feuilles et fruits abîmés, sachets, cartons ;
\item les \cle{industries} : chutes de matériaux, emballages, produits chimiques, eaux usées ;
\item les \cle{hôpitaux et centres de santé} : seringues usagées, gants, médicaments périmés, pansements souillés ;
\item les \cle{écoles et administrations} : papiers, cartouches d'encre, restes de cantine ;
\item les \cle{chantiers de construction} : gravats, sable, débris de briques et de ciment ;
\item l'\cle{agriculture et l'élevage} : fumier, pailles, résidus de récolte, emballages de pesticides.
\end{itemize}

\begin{exemplebox}[Ordre de grandeur au Cameroun]
En milieu urbain camerounais, chaque habitant produit en moyenne \SI{0,5}{kg} à \SI{1}{kg} de déchets par jour. Pour une ville de \num{100 000} habitants, cela représente jusqu'à \SI{100}{t} par jour, soit plus de \SI{36 000}{t} par an !
\end{exemplebox}

\courssub{Classification des déchets selon leur nature}

Pour gérer les déchets, il faut d'abord les classer. La classification la plus utile repose sur la \emph{nature} du déchet, c'est-à-dire la matière dont il est fait.

\begin{definition}[Déchet biodégradable]
Un déchet \cle{biodégradable} est un déchet qui peut être décomposé par des micro-organismes (bactéries, champignons) en substances simples (eau, gaz carbonique, matière organique minéralisée) qui retournent dans les cycles naturels.
\end{definition}

\begin{definition}[Déchet dangereux]
Un déchet \cle{dangereux} est un déchet qui présente un risque pour la santé humaine ou pour l'environnement, parce qu'il est toxique, inflammable, corrosif, infectieux ou radioactif.
\end{definition}

On distingue ainsi les grandes catégories suivantes :

\begin{center}
\begin{tabularx}{\linewidth}{|l|X|X|}
\hline
\rowcolor{popblueL} \textbf{Catégorie} & \textbf{Exemples} & \textbf{Caractéristique} \\
\hline
Déchets organiques & restes de nourriture, épluchures, fumier, feuilles mortes & biodégradables, fermentescibles \\
\hline
Déchets plastiques & sachets, bouteilles, emballages & non biodégradables, très persistants \\
\hline
Déchets métalliques & boîtes de conserve, ferraille, canettes & recyclables, oxydables \\
\hline
Verre & bouteilles, bocaux, vitres cassées & non biodégradable, recyclable à l'infini \\
\hline
Papier et carton & cahiers, journaux, cartons d'emballage & biodégradables, recyclables \\
\hline
Déchets inertes & gravats, sable, briques cassées & ne se décomposent pas, chimiquement stables \\
\hline
Déchets dangereux & piles, batteries, médicaments périmés, huiles usées, déchets hospitaliers & toxiques ou infectieux \\
\hline
\end{tabularx}
\end{center}

\textbf{Zoom sur trois catégories essentielles.}

\emph{Les déchets organiques} constituent la part la plus importante des déchets ménagers au Cameroun. Ils sont dégradés par les micro-organismes du sol : c'est le principe naturel de la décomposition, que l'on peut exploiter dans le compostage.

\emph{Les déchets plastiques} posent un problème majeur : fabriqués à partir du pétrole, ils ne sont pas reconnus comme nourriture par les micro-organismes. Un sac plastique persiste dans la nature pendant \num{100} à \num{1 000} ans. Ils bouchent les caniveaux (cause d'inondations à Douala et Yaoundé), étouffent les animaux et polluent les sols.

\emph{Les déchets dangereux} exigent une prudence particulière : une pile jetée au sol libère des métaux lourds (mercure, plomb, cadmium) qui contaminent le sol et les nappes phréatiques ; les seringues usagées peuvent transmettre des maladies comme le VIH ou les hépatites.

\begin{popfigure}
\begin{center}
\begin{tikzpicture}[scale=1.1]
% Camembert : organiques 60%, plastiques 10%, papier 8%, métaux/verre 7%, autres 15%
\fill[popgreen!70] (0,0) -- (90:2.6) arc (90:-126:2.6) -- cycle;
\fill[poporange!75] (0,0) -- (-126:2.6) arc (-126:-162:2.6) -- cycle;
\fill[popgold!80] (0,0) -- (-162:2.6) arc (-162:-190.8:2.6) -- cycle;
\fill[popblue!60] (0,0) -- (-190.8:2.6) arc (-190.8:-216:2.6) -- cycle;
\fill[popmuted!50] (0,0) -- (-216:2.6) arc (-216:-270:2.6) -- cycle;
\draw[popdark] (0,0) circle (2.6);
% Étiquettes
\node at (-18:1.7) {\textbf{60\,\%}};
\node[align=center] at (-18:3.5) {\textbf{Déchets organiques}\\ (épluchures, restes)};
\node at (-144:1.7) {\textbf{10\,\%}};
\node[align=center] at (-144:3.6) {\textbf{Plastiques}};
\node at (-176:1.7) {\textbf{8\,\%}};
\node[align=center] at (-186:3.7) {\textbf{Papier/carton}};
\node at (-203:1.6) {\textbf{7\,\%}};
\node[align=center] at (-215:3.7) {\textbf{Métaux/verre}};
\node at (-243:1.7) {\textbf{15\,\%}};
\node[align=center] at (-243:3.7) {\textbf{Autres}\\ (textiles, inertes...)};
\end{tikzpicture}
\end{center}
\legende{Composition moyenne des déchets ménagers urbains au Cameroun. Les déchets organiques dominent largement (environ 60\,\%), ce qui montre l'énorme potentiel du compostage dans nos villes.}
\end{popfigure}

\courssub{Biodégradabilité et durées de décomposition}

\textbf{Démarche d'investigation.} \emph{Observation :} une pelure de banane jetée derrière la maison disparaît en quelques semaines, tandis qu'un sachet plastique jeté au même endroit est toujours visible des mois plus tard. \emph{Hypothèse :} tous les déchets ne se décomposent pas à la même vitesse. \emph{Vérification :} on enterre dans un même sol une pelure de banane, un morceau de papier, un sachet plastique et un fragment de verre, et on les exhume chaque mois. \emph{Résultat :} la pelure a disparu au bout de quelques semaines, le papier en quelques mois, le plastique et le verre sont intacts. \emph{Conclusion :} la durée de décomposition dépend de la nature chimique du déchet : seuls les déchets organiques et le papier sont rapidement biodégradables.

\begin{popfigure}
\begin{center}
\begin{tikzpicture}[scale=0.95]
% Échelle logarithmique simplifiée : 0 = semaines, 1 = mois, 2 = années, 3 = siècles, 4 = millénaires
\draw[->, thick, popdark] (0,0) -- (11,0) node[below left=2pt and -30pt] {\footnotesize Durée de décomposition (échelle simplifiée)};
\foreach \x/\lab in {0.5/{quelques semaines}, 3/{quelques mois}, 5.5/{quelques années}, 8/{siècles}, 10.3/{millénaires}}
  \draw[popdark] (\x,0) -- (\x,-0.15) node[below, font=\footnotesize, align=center] {\lab};
% Barres
\fill[popgreen!70] (0,4.6) rectangle (0.9,5.2); \node[left, font=\small] at (0,4.9) {Pelure de banane};
\fill[popgold!75] (0,3.7) rectangle (3.0,4.3); \node[left, font=\small] at (0,4.0) {Journal, papier};
\fill[popgreen!50] (0,2.8) rectangle (5.2,3.4); \node[left, font=\small] at (0,3.1) {Mégot, coton};
\fill[poporange!75] (0,1.9) rectangle (8.2,2.5); \node[left, font=\small] at (0,2.2) {Sac plastique (100--1 000 ans)};
\fill[popblue!60] (0,1.0) rectangle (10.3,1.6); \node[left, font=\small] at (0,1.3) {Verre (plusieurs milliers d'années)};
\end{tikzpicture}
\end{center}
\legende{Durées de décomposition comparées de quelques déchets courants (barres horizontales, échelle logarithmique simplifiée). Plus le déchet est persistant, plus son abandon dans la nature est grave.}
\end{popfigure}

\begin{exemplebox}[Un exemple qui fait réfléchir]
Un sac plastique mis au rebut aujourd'hui peut encore exister dans la nature en l'an 2500. Autrement dit, le sachet qui a servi cinq minutes à emballer un beignet polluera le sol plus longtemps que toute l'histoire écrite du Cameroun !
\end{exemplebox}

\begin{methode}[Reconnaître la catégorie d'un déchet]
\begin{enumerate}
\item \textbf{Identifier la matière} : de quoi le déchet est-il fait (matière végétale ou animale, plastique, métal, verre, papier) ?
\item \textbf{Identifier l'origine} : d'où vient-il (cuisine, chantier, hôpital, atelier) ?
\item \textbf{Se poser la question de la biodégradabilité} : les micro-organismes peuvent-ils le décomposer en un temps raisonnable ?
\item \textbf{Se poser la question du danger} : contient-il des substances toxiques, infectieuses ou inflammables (piles, huiles, seringues, produits chimiques) ?
\item \textbf{Conclure} : ranger le déchet dans la bonne catégorie pour choisir ensuite le bon mode de traitement.
\end{enumerate}
\emph{Exemple :} une batterie de moto usagée est faite de métal et d'acide (étape 1), vient d'un atelier (étape 2), n'est pas biodégradable (étape 3) et contient du plomb et de l'acide toxiques (étape 4) : c'est un \textbf{déchet dangereux}, à ne jamais jeter dans la nature ni au feu.
\end{methode}

\begin{attention}[Brûler ne fait pas disparaître !]
Erreur fréquente : croire que tout déchet brûlé disparaît sans conséquence. En réalité, la combustion des plastiques et des déchets mélangés à l'air libre dégage des \cle{fumées toxiques} contenant notamment des \cle{dioxines}, substances cancérigènes qui polluent l'air, retombent sur les sols et contaminent les cultures. Brûler un déchet, c'est transformer une pollution du sol en pollution de l'air. Seules des installations spécialisées, à très haute température (supérieure à \SI{850}{\celsius}), peuvent incinérer certains déchets sans danger.
\end{attention}

\begin{exoresolu}[Classer les déchets d'une cour de récréation]
Énoncé. À la fin de la récréation, on ramasse dans la cour du lycée : des pelures d'orange, un sachet de biscuits vide, une pile usagée tombée d'un téléphone, une bouteille en verre cassée et une feuille de papier chiffonnée. Classe chaque déchet dans sa catégorie et indique s'il est biodégradable.
\tcblower
Solution détaillée.
\begin{itemize}
\item \textbf{Pelures d'orange} : matière végétale $\rightarrow$ déchet \emph{organique}, biodégradable (quelques semaines). Destination possible : compost.
\item \textbf{Sachet de biscuits} : matière plastique $\rightarrow$ déchet \emph{plastique}, non biodégradable (centaines d'années). Destination : recyclage ou enfouissement contrôlé.
\item \textbf{Pile usagée} : contient des métaux lourds toxiques $\rightarrow$ déchet \emph{dangereux}, non biodégradable. Destination : collecte spécifique, jamais à la poubelle ordinaire ni au feu.
\item \textbf{Bouteille en verre cassée} : déchet \emph{verre}, non biodégradable (milliers d'années) mais recyclable à l'infini. Attention aux coupures lors de la manipulation.
\item \textbf{Feuille de papier} : déchet \emph{papier}, biodégradable (quelques mois) et recyclable.
\end{itemize}
\end{exoresolu}

\begin{savaistu}[Le défi des déchets au Cameroun]
Le Cameroun produit plusieurs millions de tonnes de déchets par an, et dans les grandes villes, moins de la moitié de ces déchets est réellement collectée par les services d'hygiène (HYSACAM et sous-traitants). Le reste s'accumule dans les quartiers, les caniveaux et les cours d'eau. C'est pourquoi le tri à la source, le compostage et le recyclage, que nous étudierons dans les sections suivantes, sont des solutions essentielles : ils réduisent la quantité de déchets à collecter tout en créant des revenus et des emplois.
\end{savaistu}

\begin{aretenir}[L'essentiel de la section]
\begin{itemize}
\item Un \cle{déchet} est toute substance ou objet dont le détenteur se défait ou a l'intention de se défaire.
\item Les sources de déchets sont multiples : ménages, marchés, industries, hôpitaux, écoles, chantiers, agriculture.
\item On classe les déchets par nature : \cle{organiques} (biodégradables), \cle{plastiques}, \cle{métaux}, \cle{verre}, \cle{papier/carton}, \cle{inertes} et \cle{dangereux}.
\item La \cle{biodégradabilité} varie énormément : de quelques semaines (pelure de banane) à plusieurs milliers d'années (verre).
\item Les déchets dangereux (piles, médicaments, huiles, déchets hospitaliers) exigent une collecte spécifique ; brûler les déchets à l'air libre produit des fumées toxiques (dioxines).
\end{itemize}
\end{aretenir}

\textbf{Transition.} Maintenant que nous savons identifier et classer les déchets, nous allons voir comment les trier efficacement et appliquer la règle des 3R (Réduire, Réutiliser, Recycler) pour limiter leur impact sur notre environnement.

\coursec{Impacts des déchets sur l'environnement et la santé}

Dans la section précédente, nous avons appris à reconnaître les différents types de déchets — organiques, plastiques, dangereux — et à identifier leurs sources. Mais une question fondamentale demeure : pourquoi est-il si grave de les abandonner dans la nature ? Pourquoi un sac plastique jeté dans un caniveau ou une pile usagée enfouie dans le jardin ne sont-ils pas de simples « saletés » sans conséquence ? Cette section répond à cette question en montrant, mécanisme par mécanisme, comment un déchet mal géré devient un \cle{polluant} qui atteint les sols, l'eau, l'air, puis finalement l'être humain lui-même.

\courssub{Situation-problème : le puits du quartier}

\begin{exemplebox}[Une observation qui interroge]
Dans un quartier de Yaoundé, une famille puise son eau dans un puits situé à une trentaine de mètres d'une décharge sauvage où l'on jette de tout : ordures ménagères, piles usagées, vieilles batteries de téléphones, restes de peinture. Après quelques années, l'eau du puits prend un goût désagréable et des cas de diarrhées se multiplient dans le voisinage. Comment des déchets déposés \emph{sur le sol} ont-ils pu contaminer une eau puisée \emph{sous le sol}, à plusieurs mètres de profondeur ?
\end{exemplebox}

C'est toute la problématique de cette section : les déchets ne restent jamais là où on les jette. Ils se déplacent, se transforment, et empruntent des « routes » invisibles — l'eau qui s'infiltre, l'air qui circule, la chaîne alimentaire — pour revenir vers nous.

\courssub{Pollution des sols : une contamination lente mais durable}

\textbf{Observation et hypothèse.} Sur une ancienne décharge sauvage, on remarque que la végétation pousse mal et que les cultures voisines donnent des rendements faibles. Hypothèse : les déchets ont libéré dans le sol des substances toxiques.

\textbf{Mécanismes.} Deux processus principaux expliquent cette pollution :

\begin{itemize}
\item L'\cle{accumulation des plastiques} : un sac plastique met de 100 à 1000 ans à se dégrader complètement. Pendant ce temps, il étouffe le sol (imperméabilisation), gêne l'enracinement des plantes et se fragmente sous l'effet des UV et de l'érosion mécanique en \cle{microplastiques} (particules $<$ 5 mm) qui persistent dans l'environnement et entrent dans la chaîne alimentaire.
\item La \cle{lixiviation des métaux lourds} : les piles, batteries, ampoules et appareils électroniques contiennent des métaux toxiques (plomb Pb, mercure Hg, cadmium Cd). Sous l'action des eaux de pluie qui traversent la décharge, ces métaux sont solubilisés et entraînés vers le bas : c'est le phénomène de \cle{lixiviation}.
\end{itemize}

\begin{definition}[Lixiviat]
Le \cle{lixiviat}, appelé aussi « jus de décharge », est le liquide chargé de substances polluantes (métaux lourds, composés organiques toxiques, sels minéraux, microbes pathogènes) qui s'écoule d'une décharge lorsque l'eau de pluie la traverse. Non traité, il s'infiltre dans le sol et peut atteindre les \cle{nappes phréatiques}, c'est-à-dire les réserves d'eau souterraines exploitées pour la consommation humaine.
\end{definition}

\begin{savaistu}[Le pouvoir polluant d'une simple pile]
Une seule pile alcaline jetée au sol peut polluer jusqu'à \SI{1}{m^3} de terre pendant plusieurs dizaines d'années, à cause du mercure, du plomb ou du cadmium qu'elle contient. C'est pourquoi les piles usagées ne doivent jamais rejoindre la poubelle ordinaire ni être jetées dans la nature : elles doivent être collectées séparément dans des bacs spécifiques (souvent disponibles dans les grandes surfaces ou les mairies).
\end{savaistu}

\begin{attention}[Erreur fréquente : confusion lessivage / lixiviation]
On parle de \emph{lessivage} pour l'entraînement des éléments nutritifs (nitrates, potassium) du sol vers la nappe par l'eau de pluie. Pour les déchets, le terme exact est \emph{lixiviation} : c'est la solubilisation et l'entraînement de polluants \emph{contenus dans les déchets eux-mêmes} par l'eau de percolation. Ne confonds pas les deux termes.
\end{attention}

\begin{popfigure}
\begin{tikzpicture}[font=\small, node distance=0.4cm]
% Chaîne horizontale
\node[draw=poporange, fill=poporangeL, rounded corners, minimum width=2.5cm, minimum height=1.4cm, align=center] (dech) {\textbf{Décharge}\\ \textbf{sauvage}\\ (piles, batteries,\\ déchets électroniques)};
\node[draw=popblue, fill=popblueL, rounded corners, minimum width=2.5cm, minimum height=1.4cm, align=center, right=1.2cm of dech] (lix) {\textbf{Lixiviat}\\ (jus de décharge\\ chargé en métaux\\ lourds, sels, microbes)};
\node[draw=poppurple, fill=poppurpleL, rounded corners, minimum width=2.5cm, minimum height=1.4cm, align=center, right=1.2cm of lix] (nappe) {\textbf{Nappe}\\ \textbf{phréatique}\\ (eau souterraine\\ contaminée)};
\node[draw=popgreen, fill=popgreenL, rounded corners, minimum width=2.5cm, minimum height=1.4cm, align=center, right=1.2cm of nappe] (puits) {\textbf{Puits / Forage}\\ (captage d'eau\\ de boisson)};
\node[draw=poppink, fill=poppinkL, rounded corners, minimum width=2.5cm, minimum height=1.4cm, align=center, right=1.2cm of puits] (homme) {\textbf{Homme}\\ (intoxications\\ chroniques, maladies)};
\draw[-{Stealth}, very thick, popdark] (dech) -- (lix) node[midway, above, font=\scriptsize]{pluie + infiltration};
\draw[-{Stealth}, very thick, popdark] (lix) -- (nappe) node[midway, above, font=\scriptsize]{percolation};
\draw[-{Stealth}, very thick, popdark] (nappe) -- (puits) node[midway, above, font=\scriptsize]{pompage};
\draw[-{Stealth}, very thick, popdark] (puits) -- (homme) node[midway, above, font=\scriptsize]{consommation};
\end{tikzpicture}
\legende{Chaîne de contamination : un déchet dangereux jeté en décharge sauvage peut, par la lixiviation, atteindre la nappe phréatique, puis le point de captage, puis l'être humain.}
\end{popfigure}

\courssub{Pollution de l'eau : rivières, caniveaux et eau de boisson}

Dans beaucoup de villes camerounaises, les rivières (Mfoundi, Wouri, Dibamba...) et les caniveaux servent malheureusement de poubelles à ciel ouvert. Les conséquences sont directes et mesurables :

\begin{itemize}
\item \textbf{Eutrophisation et asphyxie du milieu aquatique} : Les déchets organiques en décomposition sont dégradés par les micro-organismes aérobies qui consomment l'\cle{oxygène dissous} dans l'eau. La concentration en $\ce{O2}$ chute (demande biochimique en oxygène DBO5 élevée), entraînant la mort des poissons et de la faune benthique par asphyxie. La rivière devient « morte » (eau noire, odeurs de $\ce{H2S}$).
\item \textbf{Contamination chimique} : Les substances toxiques (métaux lourds, hydrocarbures, résidus de pesticides, médicaments) dissoutes dans le lixiviat ou jetées directement contaminent l'eau utilisée pour la boisson, la cuisine ou l'irrigation.
\item \textbf{Risque microbiologique} : Les eaux souillées par des matières fécales (couches, absence de latrines) et les ordures deviennent des milieux de survie pour les bactéries (\emph{Vibrio cholerae}, \emph{Salmonella typhi}), virus et parasites responsables de maladies hydriques.
\end{itemize}

\begin{exemplebox}[La mort d'une rivière urbaine]
Le fleuve Mfoundi à Yaoundé, ou la rivière Dibamba à Douala, illustrent ce phénomène : la charge organique quotidienne dépasse la capacité d'auto-épuration du cours d'eau. La décomposition anaérobie s'installe, l'oxygène disparaît, les poissons meurent, et l'eau devient impropre à tout usage. Les pêcheurs perdent leur ressource, les riverains une source d'eau, et l'écosystème s'effondre.
\end{exemplebox}

\courssub{Pollution de l'air : combustion sauvage et gaz de décharge}

\textbf{La combustion sauvage (brûlage à l'air libre).} Pratique encore courante dans les cours et aux abords des routes, elle dégage des fumées très toxiques :
\begin{itemize}
\item Particules fines (\cle{PM2,5}, \cle{PM10}) pénétrant profond dans les poumons.
\item Monoxyde de carbone (\ce{CO}), asphyxiant.
\item \cle{Dioxines} et furanes : polluants organiques persistants (POP), très toxiques, cancérigènes, formés lors de la combustion \emph{incomplète} de plastiques contenant du chlore (PVC) ou en présence de métaux catalyseurs (cuivre).
\end{itemize}
Ces fumées irritent les voies respiratoires, aggravent l'asthme et la BPCO, et augmentent le risque de cancers du poumon, surtout chez les enfants et les personnes âgées exposées quotidiennement.

\textbf{Les gaz de décharge (biogaz non maîtrisé).} Même sans combustion, une décharge pollue l'air : les déchets organiques enfouis se décomposent \emph{en l'absence d'oxygène} (décomposition \cle{anaérobie} / méthanisation sauvage) et dégagent un mélange gazeux : \cle{méthane} (\ce{CH4}, 50-60\%), dioxyde de carbone (\ce{CO2}, 30-40\%), et traces de sulfure d'hydrogène (\ce{H2S}, odeur d'œuf pourri) et de composés organiques volatils.

\begin{definition}[Gaz à effet de serre (GES)]
Un \cle{gaz à effet de serre} est un gaz atmosphérique qui absorbe le rayonnement infrarouge émis par la Terre et piège une partie de la chaleur, contribuant ainsi au réchauffement climatique. Les principaux d'origine anthropique sont le dioxyde de carbone (\ce{CO2}), le méthane (\ce{CH4}) et le protoxyde d'azote (\ce{N2O}).
\end{definition}

\begin{propriete}[Pouvoir de réchauffement global (PRG) du méthane — résultat admis]
Selon les rapports du GIEC (AR5/AR6), le méthane a un pouvoir de réchauffement global (PRG) environ \cle{28 à 30 fois supérieur} à celui du dioxyde de carbone sur un horizon de 100 ans (valeur souvent arrondie à 25 dans les manuels scolaires actuels, référence AR4). Ce résultat est \emph{admis} : il n'est pas à démontrer au Probatoire. Il montre que les décharges sauvages, en émettant du $\ce{CH4}$ non capté, contribuent significativement au \cle{changement climatique}.
\end{propriete}

\courssub{Impacts sanitaires : quand les déchets rendent malade}

Les déchets mal gérés sont des alliés directs des maladies, au premier rang desquelles figurent les premières causes de morbidité au Cameroun :

\begin{itemize}
\item \textbf{Gîtes larvaires artificiels} : les emballages abandonnés (boîtes de conserve, pneus usagés, sachets plastiques, noix de coco) retiennent l'eau de pluie et deviennent des lieux de ponte privilégiés pour \emph{Anopheles} (vecteur du \cle{paludisme}, première cause de consultation dans les formations sanitaires camerounaises) et \emph{Aedes} (vecteur de la dengue, du chikungunya, du Zika).
\item \textbf{Eaux stagnantes souillées / contamination fécale} : elles favorisent la propagation du \cle{choléra} (\emph{Vibrio cholerae}) et de la \cle{fièvre typhoïde} (\emph{Salmonella Typhi}), maladies hydriques à déclaration obligatoire, dont les épidémies récurrentes au Cameroun sont liées à l'insalubrité.
\item \textbf{Blessures et infections} : les objets tranchants (verre brisé, lames de rasoir, canettes, aiguilles) provoquent des coupures qui peuvent s'infecter (tétanos, septicémie), notamment chez les enfants qui jouent pieds nus et chez les récupérateurs de déchets (souvent sans EPI).
\item \textbf{Intoxications chimiques aiguës ou chroniques} : par contact cutané, inhalation ou ingestion accidentelle de produits chimiques (médicaments périmés, pesticides domestiques, métaux lourds du lixiviat).
\end{itemize}

\courssub{Impacts urbains : caniveaux bouchés et inondations}

\begin{exemplebox}[Les inondations de Douala et Yaoundé : un lien direct avec les plastiques]
En saison des pluies (mars-juin, septembre-novembre), plusieurs quartiers de Douala (Makepe, Bepanda, New Bell) et de Yaoundé (Mvog-Ada, Essos, Mokolo) connaissent des inondations récurrentes. Une cause majeure, documentée par la CUC et la CUY : les \cle{caniveaux bouchés par les plastiques} (sachets, bouteilles) et les ordures, qui réduisent drastiquement la section hydraulique et empêchent l'écoulement des eaux de ruissellement. L'eau monte alors dans les maisons, détruit les biens, coupe les routes et stagne ensuite pendant des jours — créant au passage de nouveaux gîtes à moustiques. Un simple sachet plastique jeté dans la rue peut donc contribuer à inonder tout un quartier.
\end{exemplebox}

Au-delà des inondations, l'accumulation des déchets dégrade le cadre de vie : odeurs nauséabondes (putréfaction, \ce{H2S}, \ce{NH3}), paysages souillés, prolifération de rongeurs (rats, vecteurs de leptospirose), baisse de la valeur foncière, image négative des quartiers. C'est ce qu'on appelle la \cle{dépréciation du cadre de vie}, qui a un coût économique réel (santé, perte de productivité, coût de curage des caniveaux).

\courssub{Impacts sur la faune et la biodiversité}

Les animaux paient aussi un lourd tribut, souvent invisible :
\begin{itemize}
\item \textbf{Faune domestique / élevage urbain} : chèvres, moutons et vaches en divagation dans nos villes ingèrent des sachets plastiques qui s'accumulent dans le rumen (bourrelets plastiques), obstruent le tube digestif, provoquent une fausse satiété, des occlusions intestinales et la mort de l'animal. Perte économique directe pour l'éleveur.
\item \textbf{Faune sauvage et marine} : les déchets plastiques charriés par les fleuves (Wouri, Sanaga, Nyong) atteignent l'océan Atlantique. Au large de Kribi ou de Limbé, tortues marines (confondant sacs plastiques et méduses), oiseaux de mer et poissons ingèrent des débris plastiques : fausse satiété, blocage digestif, transfert de polluants adsorbés (PCB, DDT) dans les tissus. Le plastique jeté à Douala peut ainsi finir dans l'estomac d'une tortue au large de Kribi.
\end{itemize}

\courssub{Synthèse : l'arbre des impacts}

\begin{popfigure}
\begin{tikzpicture}[font=\small,
  racine/.style={draw=popdark, fill=popink, text=white, rounded corners, align=center, minimum width=3.5cm, minimum height=1.1cm},
  pollu/.style={draw=popblue, fill=popblueL, rounded corners, align=center, minimum width=2.8cm, minimum height=1cm},
  conseq/.style={draw=poporange, fill=poporangeL, rounded corners, align=center, minimum width=2.8cm, minimum height=1cm},
  every edge/.style={draw=popmuted, -{Stealth}, thick}]
\node[racine, align=center] (R){\textbf{Déchet mal géré}\\(abandon, brûlage, décharge sauvage)};
\node[pollu, below left=1.5cm and 4cm of R, align=center] (sol){\textbf{Pollution}\\ \textbf{des sols}};
\node[pollu, below=1.5cm of R, align=center] (eau){\textbf{Pollution}\\ \textbf{de l'eau}};
\node[pollu, below right=1.5cm and 4cm of R, align=center] (air){\textbf{Pollution}\\ \textbf{de l'air}};
\node[conseq, below=1.2cm of sol, align=center] (c1){Imperméabilisation,\\ stérilisation, lixiviat\\ $\rightarrow$ nappes contaminées};
\node[conseq, below=1.2cm of eau, align=center] (c2){Eutrophisation, asphyxie,\\ choléra, typhoïde,\\ eau impropre};
\node[conseq, below=1.2cm of air, align=center] (c3){Fumées toxiques (dioxines),\\ GES (méthane),\\ maladies respiratoires};
\draw (R) edge (sol);
\draw (R) edge (eau);
\draw (R) edge (air);
\draw (sol) edge (c1);
\draw (eau) edge (c2);
\draw (air) edge (c3);
\end{tikzpicture}
\legende{Arbre des impacts : un déchet mal géré dégrade les trois compartiments environnementaux (sol, eau, air), avec des conséquences sanitaires (maladies), écologiques (biodiversité) et économiques (inondations, dépréciation).}
\end{popfigure}

\begin{exoresolu}[Analyser une situation de pollution locale]
\textbf{Énoncé.} Dans un village de la région du Centre, les habitants jettent leurs ordures ménagères, y compris des piles usagées et des médicaments périmés, dans une dépression du terrain (ancienne carrière) située à 50 m en amont (pente) du puits principal du village. Après la saison des pluies, plusieurs enfants présentent des douleurs abdominales, des nausées et une anémie inexpliquée. 
\begin{enumerate}
\item En t'appuyant sur les mécanismes vus en cours, explique comment les piles et médicaments ont pu rendre les enfants malades. Nomme précisément les phénomènes physiques et chimiques en jeu.
\item Propose trois mesures concrètes, applicables au niveau villageois, pour protéger durablement la ressource en eau.
\end{enumerate}
\tcblower
\textbf{Solution détaillée.}
\begin{enumerate}
\item \textbf{Chaîne de contamination :}
\begin{enumerate}
\item \textbf{Lixiviation} : Les eaux de pluie infiltrées (percolation) traversent la masse de déchets. Elles solubilisent les métaux lourds (Pb, Hg, Cd des piles) et les principes actifs résiduels des médicaments (antibiotiques, hormones...). Ce liquide pollué est le \textbf{lixiviat}.
\item \textbf{Transfert vers la nappe} : Sous l'effet de la gravité et de la pente géologique, le lixiviat migre verticalement (zone non saturée) puis horizontalement dans la \textbf{nappe phréatique} (zone saturée) qui alimente le puits. La distance de 50 m est très courte pour une nappe libre ; le temps de transfert peut être de quelques jours à quelques semaines.
\item \textbf{Exposition humaine} : Les villageois pompent cette eau contaminée pour boire, cuisiner, se laver. L'ingestion chronique de métaux lourds (plomb, cadmium) et de résidus médicamenteux explique le tableau clinique : douleurs abdominales (colique saturnine pour le Pb), nausées, anémie (inhibition de la synthèse de l'hème par le Pb), néphrotoxicité (Cd).
\end{enumerate}
\item \textbf{Mesures de protection (niveau villageois) :}
\begin{enumerate}
\item \textbf{Suppression de la source} : Fermeture immédiate de la décharge sauvage. Mise en place d'un tri à la source : collecte séparée des piles/batteries (stockage en bidons fermés pour évacuation vers centre de traitement), médicaments périmés (retour au centre de santé / pharmacie pour incinération), plastiques (réutilisation/vente aux récupérateurs), organiques (compostage).
\item \textbf{Protection du point de captage} : Réalisation d'une \textbf{marge de protection} autour du puits (rayon 15-20 m minimum) : clôture, dallage étanche en pente vers l'extérieur, couverture hermétique, pompe à main (pas de seau/ corde souillés). Curage régulier du puits.
\item \textbf{Traitement de l'eau à domicile} : En attendant l'assainissement de la nappe, promotion du traitement de l'eau de boisson au point d'usage : chloration (pastilles \emph{Aquatabs} / \emph{WaterGuard} disponibles au Cameroun), ébullition 10 min, ou filtration céramique. Sensibilisation continue par le comité de gestion de l'eau / ASC (Agent de Santé Communautaire).
\end{enumerate}
\end{enumerate}
\end{exoresolu}

\begin{aretenir}[L'essentiel de la section — À connaître pour le Probatoire]
Un déchet mal géré pollue les trois milieux physiques :
\begin{itemize}
\item \textbf{Sols} : imperméabilisation par plastiques (persistance 100-1000 ans), fragmentation en microplastiques, \cle{lixiviation} des métaux lourds (Pb, Hg, Cd) $\rightarrow$ stérilisation, contamination nappes.
\item \textbf{Eau} : eutrophisation par matière organique (consommation $\ce{O2}$, asphyxie faune), contamination chimique (lixiviat) et microbiologique (choléra, typhoïde) $\rightarrow$ eau impropre.
\item \textbf{Air} : combustion sauvage $\rightarrow$ dioxines (cancérigènes), particules fines (maladies respiratoires) ; décharge anaérobie $\rightarrow$ \ce{CH4} (GES, PRG $\approx$ 25-30 $\times$ \ce{CO2}) $\rightarrow$ changement climatique.
\end{itemize}
\textbf{Conséquences sanitaires majeures au Cameroun} : paludisme (gîtes larvaires plastiques), choléra/typhoïde (eau souillée), intoxications métaux lourds, blessures/tétanos.
\textbf{Conséquences urbaines} : inondations par bouchage caniveaux (plastiques), dépréciation cadre de vie.
\textbf{Chaîne type à retenir} : \textbf{Décharge sauvage $\xrightarrow{\text{pluie}}$ Lixiviat $\xrightarrow{\text{infiltration}}$ Nappe phréatique $\xrightarrow{\text{captage}}$ Puits $\xrightarrow{\text{consommation}}$ Homme malade}.
\end{aretenir}

Maintenant que tu comprends \emph{pourquoi} les déchets sont dangereux, la prochaine section te montrera \emph{comment} agir concrètement : le tri des déchets et la fameuse règle des 3R (Réduire, Réutiliser, Recycler).

\coursec{Le tri des déchets et la règle des 3R}

Dans la section précédente, nous avons vu que les déchets mal gérés polluent les sols, les eaux et l'air, et qu'ils menacent directement la santé humaine : eaux stagnantes favorisant les moustiques vecteurs du paludisme, fumées toxiques des plastiques brûlés, contamination des nappes phréatiques par les piles et les produits chimiques. Face à ces dangers, une question simple se pose : \emph{que faire concrètement de nos déchets ?} La réponse commence par un geste élémentaire, à la portée de chacun, dès la maison, l'école ou le marché : le \cle{tri}. C'est l'objet de cette section, qui te donnera aussi un cadre de pensée universel : la \cle{règle des 3R}.

\courssub{Le tri à la source : définition et mise en pratique}

\textbf{Une observation pour commencer.} Dans une concession de Yaoundé, on jette tout dans le même sac : épluchures de manioc, sachets plastiques, boîtes de conserve vides, vieilles piles de radio, restes de médicaments. Une fois mélangés, ces déchets deviennent quasiment inutilisables : les épluchures souillent le plastique, les piles contaminent l'ensemble. À l'inverse, au marché de Mokolo, certains vendeurs séparent les feuilles et tiges abîmées (revendues comme fourrage ou compostées) des emballages. Le même déchet peut donc être une \emph{nuisance} ou une \emph{ressource} : tout dépend du moment où l'on décide de le séparer.

\begin{definition}[Tri à la source]
Le \cle{tri à la source} est la séparation des déchets selon leur nature, effectuée \textbf{à l'endroit même où ils sont produits} (cuisine, salle de classe, atelier, étal de marché), en vue de leur permettre un traitement adapté : compostage, recyclage, valorisation ou élimination sécurisée.
\end{definition}

L'adjectif « à la source » est essentiel : trier un sac déjà mélangé est long, coûteux, dangereux et souvent impossible. Trier au moment où l'on jette ne prend que quelques secondes.

\textbf{Les quatre catégories pratiques.} Dans le cadre scolaire et domestique, on retient quatre catégories, chacune associée à un bac ou un sac distinct :

\begin{center}
\begin{tabularx}{\linewidth}{|l|X|X|}\hline
\rowcolor{popblueL} \textbf{Catégorie} & \textbf{Exemples} & \textbf{Destination} \\ \hline
Déchets organiques & épluchures, restes de repas, feuilles, coquilles d'{\oe}ufs & compostage, biogaz, alimentation animale \\ \hline
Plastiques & sachets, bouteilles en plastique, bidons & recyclage, valorisation \\ \hline
Métaux et verre & boîtes de conserve, canettes, bouteilles en verre & réutilisation (consigne), recyclage \\ \hline
Déchets dangereux & piles, ampoules, médicaments périmés, produits chimiques & collecte spéciale, jamais en décharge ordinaire \\ \hline
\end{tabularx}
\end{center}

\textbf{Le code couleur des bacs.} Pour rendre le tri intuitif, on associe souvent une couleur à chaque catégorie. Le code le plus répandu (adaptable selon les communes) est le suivant :

\begin{itemize}
\item \textbf{bac vert} : déchets organiques (la couleur de la végétation) ;
\item \textbf{bac jaune} : déchets recyclables (plastiques, métaux, verre, papier) ;
\item \textbf{bac rouge} : déchets dangereux (la couleur de l'alerte).
\end{itemize}

\begin{attention}[Ne jamais mélanger les déchets dangereux aux ordures ménagères]
L'erreur la plus fréquente — et la plus grave — consiste à jeter les \textbf{piles, ampoules, médicaments périmés et restes de peinture} dans la poubelle ordinaire. Une seule pile bouton peut contaminer jusqu'à \SI{1}{m^3} de terre et des milliers de litres d'eau par le mercure, le plomb ou le cadmium qu'elle contient. Les déchets dangereux doivent être stockés à part et remis à une filière de collecte spéciale (mairie, pharmacie pour les médicaments, points de collecte des piles).
\end{attention}

\begin{methode}[Trier correctement les déchets d'une cuisine ou d'une salle de classe]
\begin{enumerate}
\item \textbf{Identifier} chaque déchet avant de le jeter : de quoi est-il fait ? Est-il souillé, toxique, réutilisable ?
\item \textbf{Préparer quatre contenants} (bacs, seaux ou sacs) étiquetés : organiques / plastiques / métaux-verre / dangereux. Utiliser le code couleur si possible.
\item \textbf{Déposer immédiatement} chaque déchet dans le bon contenant : épluchures au vert, sachets propres au jaune, boîtes vides rincées au jaune, piles au rouge.
\item \textbf{Vérifier les cas douteux} : un emballage souillé de nourriture n'est pas recyclable tel quel ; on le rince ou on le jette avec les ordures résiduelles.
\item \textbf{Orienter chaque contenant} vers sa filière : composteur ou fosse pour le vert, point de collecte ou récupérateur pour le jaune, collecte spéciale pour le rouge.
\end{enumerate}
\end{methode}

\courssub{La règle des 3R : Réduire, Réutiliser, Recycler}

\textbf{Situation problème.} Deux familles consomment la même quantité de nourriture par semaine. La première achète tout en sachets jetables, jette les bouteilles après usage et met tout à la poubelle : elle produit un grand sac de déchets par jour. La seconde utilise un panier au marché, consigne ses bouteilles en verre et composte ses épluchures : elle ne remplit qu'un petit sac tous les trois jours. La différence ne tient ni au revenu ni au hasard, mais à l'application d'une règle simple : les \cle{3R}.

\begin{definition}[La règle des 3R]
La règle des \cle{3R} est un principe de gestion durable des déchets fondé sur trois actions hiérarchisées :
\begin{itemize}
\item \textbf{Réduire} : limiter la production de déchets à la source, en évitant ce qui deviendra un déchet ;
\item \textbf{Réutiliser} : donner une seconde vie à un objet sans le transformer, pour le même usage ou un usage nouveau ;
\item \textbf{Recycler} : transformer un déchet en nouvelle matière première pour fabriquer un nouvel objet.
\end{itemize}
\end{definition}

\begin{definition}[Recyclage et valorisation]
Le \cle{recyclage} est l'ensemble des procédés qui transforment un déchet en matière première ou en produit nouveau (ex. : plastique fondu en pavés, ferraille refondue).
La \cle{valorisation} est plus large : c'est toute opération qui donne une \textbf{valeur} à un déchet — matière (recyclage, compost) ou énergie (biogaz, incinération avec récupération de chaleur).
\end{definition}

\textbf{Réduire : le premier et le meilleur des gestes.} Le meilleur déchet est celui qu'on ne produit pas. Exemples concrets au quotidien :
\begin{itemize}
\item utiliser un \textbf{panier ou un sac réutilisable} au marché au lieu d'accumuler les sachets plastiques ;
\item \textbf{refuser les emballages superflus} (double sachet, suremballage) ;
\item \textbf{acheter en vrac} (riz, haricots, huile) avec ses propres contenants ;
\item réparer plutôt que jeter (savate, radio, marmite).
\end{itemize}

\textbf{Réutiliser : la seconde vie.} Avant de recycler, on réutilise : cela ne coûte aucune énergie de transformation.
\begin{itemize}
\item les \textbf{bouteilles en verre consignées} (boissons gazeuses, bière) sont rapportées, lavées et remplies à nouveau ;
\item les \textbf{boîtes de conserve} nettoyées deviennent des pots à crayons, des pots de fleurs ou des mesures ;
\item les \textbf{pneus usés} sont transformés en fauteuils, jardinières ou balançoires — un artisanat très visible dans nos villes ;
\item les vêtements passent d'aîné en cadet, puis deviennent chiffons.
\end{itemize}

\textbf{Recycler : la transformation.} Quand l'objet ne peut plus servir, sa matière peut renaître : le plastique fondu devient pavés de rue ou ardoises, la ferraille est refondue, le papier repasse en pâte. Le recyclage exige que le déchet soit \emph{propre et trié} — d'où l'importance du tri à la source.

\begin{popfigure}
\begin{center}
\begin{tikzpicture}[scale=1]
\draw[thick, popblue, fill=popblue!15] (0,2.2) ellipse (2.6 and 0.55);
\draw[thick, popgreen, fill=popgreen!15] (0,0) ellipse (2.6 and 0.55);
\draw[thick, poporange, fill=poporange!15] (0,-2.2) ellipse (2.6 and 0.55);
\node[font=\bfseries] at (0,2.2) {RÉDUIRE};
\node[font=\small] at (0,1.75) {\'eviter le déchet à la source};
\node[font=\bfseries] at (0,0) {RÉUTILISER};
\node[font=\small] at (0,-0.45) {donner une seconde vie};
\node[font=\bfseries] at (0,-2.2) {RECYCLER};
\node[font=\small] at (0,-2.65) {transformer en nouvelle matière};
\draw[-{Stealth[length=3mm]}, very thick, poppurple] (3.1,1.9) arc (60:-20:1.4);
\draw[-{Stealth[length=3mm]}, very thick, poppurple] (3.1,-0.3) arc (60:-20:1.4);
\draw[-{Stealth[length=3mm]}, very thick, poppurple] (-3.1,-2.5) arc (240:200:1.4);
\draw[-{Stealth[length=3mm]}, very thick, poppurple] (-3.1,-0.3) arc (240:200:1.4);
\node[font=\small\itshape, popmuted] at (4.6,0.9) {si on ne peut pas réduire};
\node[font=\small\itshape, popmuted] at (4.6,-1.3) {si on ne peut pas réutiliser};
\node[font=\small\itshape, popmuted] at (-4.9,-1.3) {la matière renaît};
\node[font=\small\itshape, popmuted] at (-4.9,0.9) {et repart en cycle};
\end{tikzpicture}
\end{center}
\legende{La règle des 3R en cercle : on cherche d'abord à Réduire, puis à Réutiliser, puis à Recycler ; la matière recyclée réinjecte le cycle de production.}
\end{popfigure}

\courssub{La hiérarchie des modes de traitement}

Les 3R s'inscrivent dans une échelle plus générale, appelée \cle{hiérarchie des modes de gestion des déchets}. Toutes les solutions ne se valent pas : on classe les options de la plus souhaitable (en haut) à la moins souhaitable (en bas) :

\begin{enumerate}
\item \textbf{Prévention} : ne pas produire le déchet (réduire) — priorité absolue ;
\item \textbf{Réutilisation} : resservir l'objet tel quel ;
\item \textbf{Recyclage} : transformer la matière ;
\item \textbf{Valorisation énergétique} : récupérer l'énergie (biogaz, incinération avec récupération de chaleur) ;
\item \textbf{Mise en décharge} : enfouissement ou dépôt — \textbf{solution de dernier recours}, réservée aux déchets ultimes qu'on ne peut ni réutiliser, ni recycler, ni valoriser.
\end{enumerate}

\begin{popfigure}
\begin{center}
\begin{tikzpicture}[scale=1]
\fill[popgreen!70] (-0.7,4.4) -- (0.7,4.4) -- (1.4,3.4) -- (-1.4,3.4) -- cycle;
\fill[popgreen!40] (-1.4,3.4) -- (1.4,3.4) -- (2.1,2.4) -- (-2.1,2.4) -- cycle;
\fill[popgold!50] (-2.1,2.4) -- (2.1,2.4) -- (2.8,1.4) -- (-2.8,1.4) -- cycle;
\fill[poporange!45] (-2.8,1.4) -- (2.8,1.4) -- (3.5,0.4) -- (-3.5,0.4) -- cycle;
\fill[popink!25] (-3.5,0.4) -- (3.5,0.4) -- (4.2,-0.6) -- (-4.2,-0.6) -- cycle;
\node[font=\small\bfseries] at (0,3.9) {Prévention};
\node[font=\small\bfseries] at (0,2.9) {Réutilisation};
\node[font=\small\bfseries] at (0,1.9) {Recyclage};
\node[font=\small\bfseries] at (0,0.9) {Valorisation énergétique};
\node[font=\small\bfseries] at (0,-0.1) {Mise en décharge};
\draw[-{Stealth[length=3mm]}, very thick, popblue] (-4.8,-0.6) -- (-4.8,4.4);
\node[rotate=90, font=\small\itshape, popblue] at (-5.2,1.9) {priorité croissante};
\node[font=\small\itshape, popmuted, align=left] at (6.3,3.9) {à privilégier};
\node[font=\small\itshape, popmuted, align=left] at (6.3,-0.1) {dernier recours};
\end{tikzpicture}
\end{center}
\legende{Pyramide de la hiérarchie des modes de gestion des déchets : la prévention est au sommet (option la plus favorable pour l'environnement), la mise en décharge à la base (dernier recours).}
\end{popfigure}

\begin{aretenir}[L'essentiel]
Le tri à la source sépare les déchets selon leur nature dès leur production, grâce à des bacs différenciés (vert : organiques ; jaune : recyclables ; rouge : dangereux). La règle des 3R — \textbf{Réduire, Réutiliser, Recycler} — et la hiérarchie des traitements guident chaque décision : le meilleur déchet est celui qu'on évite, la décharge n'est que le dernier recours.
\end{aretenir}

\begin{exemplebox}[Un exemple chiffré : le compostage divise par deux les ordures]
Dans une famille camerounaise moyenne, les déchets organiques (épluchures, restes de repas, feuilles) représentent environ \textbf{50 à 60\,\% du poids} des ordures ménagères. Si une famille produit \SI{10}{kg} de déchets par semaine, trier et composter la fraction organique retire environ \SI{6}{kg} du sac poubelle : le volume à collecter est réduit de plus de moitié, la décharge est moins sollicitée, et la famille récupère un engrais naturel gratuit pour son jardin ou son champ.
\end{exemplebox}

\begin{exoresolu}[Appliquer les 3R à une salle de classe]
Énoncé. Une salle de Première A produit chaque semaine : des brouillons de papier, des sachets de beignets, des trognons et pelures de fruits de la récréation, deux piles usagées de la calculatrice et des stylos vides. Propose une gestion conforme aux 3R et à la hiérarchie des traitements.
\tcblower
Solution détaillée.
\begin{enumerate}
\item \textbf{Réduire} : écrire au recto-verso, utiliser une gourde plutôt que des sachets d'eau, apporter les fruits dans un bol réutilisable.
\item \textbf{Réutiliser} : les feuilles imprimées d'un seul côté servent de brouillon ; les stylos vides peuvent être rechargés.
\item \textbf{Trier} : installer trois contenants — vert (pelures et trognons, vers le compost de l'école), jaune (papiers propres et sachets plastiques propres, vers un récupérateur), rouge (les deux piles, remises à un point de collecte spécial).
\item \textbf{Recycler / valoriser} : le papier trié part en filière de recyclage ; les organiques deviennent du compost pour le jardin scolaire.
\item \textbf{Dernier recours} : seuls les résidus non valorisables (emballages souillés) rejoignent la collecte municipale vers la décharge.
\end{enumerate}
Conformément à la hiérarchie, la décharge ne reçoit que les déchets ultimes.
\end{exoresolu}

\courssub{Les acteurs du tri et de la collecte}

Le tri ne fonctionne que si chacun joue son rôle, de la cuisine à la décharge contrôlée :

\begin{itemize}
\item \textbf{les ménages} : premiers trieurs, ils séparent à la source et réduisent leurs achats d'emballages ;
\item \textbf{les écoles} : lieux d'éducation et de pratique (compost scolaire, bacs de tri, clubs environnement) ;
\item \textbf{les mairies} : organisent la collecte, installent les points d'apport et réglementent les dépôts ;
\item \textbf{les entreprises de collecte} : au Cameroun, \textbf{HYSACAM} assure l'enlèvement des ordures dans les grandes villes comme Douala et Yaoundé ;
\item \textbf{les récupérateurs informels} : trient et revendent plastiques, métaux et verre aux filières de recyclage ;
\item \textbf{les associations et ONG} : sensibilisent, forment au compostage et organisent des journées de nettoyage.
\end{itemize}

\begin{savaistu}[Les récupérateurs informels, premiers recycleurs de nos villes]
Dans les décharges de Yaoundé (Nkolfoulou) et de Douala, des centaines de récupérateurs informels trient chaque jour, à la main, les plastiques, les métaux et le verre. Leur travail détourne déjà une part importante de ces matériaux de l'enfouissement et alimente toute une économie du recyclage, souvent méconnue. Reconnaître, organiser et protéger ces acteurs (gants, vaccination, coopératives) est un enjeu majeur de la gestion durable des déchets au Cameroun.
\end{savaistu}

\begin{exercice}[À toi de jouer]
\begin{enumerate}
\item Définis le tri à la source et explique pourquoi il est plus efficace que le tri en décharge.
\item Classe les déchets suivants dans les quatre catégories (organiques / plastiques / métaux-verre / dangereux) : épluchure de banane plantain, pile usagée, bouteille en verre, sachet plastique propre, médicament périmé, boîte de sardine vide, tige de manioc.
\item Pour chacun des objets suivants, propose une action de Réduction, de Réutilisation ou de Recyclage : sachet plastique, pneu usé, bouteille en verre consignée.
\item Rappelle les cinq niveaux de la hiérarchie des modes de traitement, du plus souhaitable au moins souhaitable.
\end{enumerate}
\end{exercice}

\begin{corrige}[Exercice : À toi de jouer]
\begin{enumerate}
\item Le tri à la source est la séparation des déchets selon leur nature à l'endroit où ils sont produits. Il est plus efficace car les déchets restent propres et non mélangés, donc directement valorisables ; en décharge, le mélange et la souillure rendent le tri lent, dangereux et souvent impossible.
\item Organiques : épluchure de plantain, tige de manioc. Plastiques : sachet propre. Métaux-verre : bouteille en verre, boîte de sardine vide (rincée). Dangereux : pile usagée, médicament périmé.
\item Sachet plastique : réduire (panier réutilisable) ou recycler (filière plastique). Pneu usé : réutiliser (fauteuil, jardinière). Bouteille consignée : réutiliser (retour au consignataire).
\item Prévention $>$ réutilisation $>$ recyclage $>$ valorisation énergétique $>$ mise en décharge (dernier recours).
\end{enumerate}
\end{corrige}

En résumé, trier à la source et appliquer la règle des 3R transforment chacun de nous en acteur de la gestion durable des déchets. Mais que deviennent concrètement les déchets triés ? C'est l'objet de la section suivante, consacrée aux principes et aux filières de transformation et de recyclage.

\coursec{Transformation et recyclage : principes et filières}

Dans la section précédente, tu as appris à trier les déchets selon leur nature et à appliquer la règle des 3R : \cle{Réduire}, \cle{Réutiliser}, \cle{Recycler}. Mais une fois le tri effectué, que devient réellement un sac de bouteilles plastiques ramassé dans une cour d'école à Bafoussam ? Comment un vieux pneu abandonné sur le bord de la route de Douala peut-il se transformer en sandale ou en fauteuil ? C'est toute la question de cette section : comprendre les \cle{principes} de la transformation et du recyclage, et découvrir les \cle{filières} concrètes qui donnent une seconde vie aux déchets.

\courssub{Les notions fondamentales : transformation, recyclage, surcyclage}

\textbf{Une observation pour commencer.} Sur le marché de Mokolo à Yaoundé, on trouve des arrosoirs fabriqués à partir de bidons d'huile vides, des nattes tressées avec des sachets plastiques usagés et des chaises dont l'assise est faite de chutes de bois. Rien de tout cela n'a été acheté « neuf » : des objets destinés à la poubelle ont été \emph{transformés} en objets utiles. Cette pratique quotidienne illustre parfaitement les notions que nous allons définir rigoureusement.

\begin{definition}[Transformation d'un déchet]
La \cle{transformation} d'un déchet est l'ensemble des opérations qui modifient sa \textbf{forme} ou sa \textbf{nature} (broyage, fonte, découpe, tressage, compostage\ldots) afin d'obtenir un \textbf{nouveau produit} utilisable, différent du déchet de départ.
\end{definition}

\begin{definition}[Recyclage]
Le \cle{recyclage} est la \textbf{réintroduction d'un déchet comme matière première} dans un cycle de production, afin de fabriquer un produit neuf, identique ou différent du produit d'origine. Exemple : des bouteilles en verre fondues pour fabriquer de nouvelles bouteilles.
\end{definition}

\begin{definition}[Surcyclage (upcycling)]
Le \cle{surcyclage} (ou \emph{upcycling}) est une forme particulière de transformation dans laquelle le déchet est converti en un produit de \textbf{valeur supérieure ou égale} à celle du produit d'origine, souvent par des techniques artisanales, sans destruction chimique du matériau. Exemple : un pneu usagé transformé en fauteuil ou en jardinière.
\end{definition}

\begin{attention}[Ne pas confondre réutilisation et recyclage]
C'est l'erreur la plus fréquente au Probatoire. La \textbf{réutilisation} consiste à employer à nouveau un objet \emph{tel quel}, pour le même usage ou un usage direct, \textbf{sans transformation de sa matière} : remplir à nouveau une bouteille consignée, utiliser un bocal vide comme boîte de rangement. Le \textbf{recyclage}, lui, implique une \textbf{transformation de la matière} (fonte, broyage, défibrage) : la bouteille est fondue pour fabriquer un nouveau récipient. Retiens : réutiliser = même objet ; recycler = nouvelle matière.
\end{attention}

\courssub{Les grandes filières de recyclage}

Une \cle{filière de recyclage} est la chaîne organisée d'opérations qui conduit du déchet collecté au produit fini. Étudions les quatre filières classiques, puis la transformation artisanale.

\textbf{1. Le recyclage du plastique.} C'est la filière la plus visible dans nos villes, car les déchets plastiques (sachets, bouteilles) envahissent les rues et les caniveaux. La chaîne comprend six étapes :

\begin{enumerate}
\item \textbf{Collecte} : ramassage des plastiques auprès des ménages, des marchés ou par les associations de jeunes ;
\item \textbf{Tri par type} : on sépare les plastiques selon leur nature (PET des bouteilles d'eau, PEHD des bidons, sachets\ldots), car ils ne fondent pas à la même température ;
\item \textbf{Lavage} : élimination des saletés, restes alimentaires et étiquettes ;
\item \textbf{Broyage} : le plastique propre est réduit en petits fragments (paillettes) ;
\item \textbf{Fonte} : les paillettes sont chauffées jusqu'à fusion ;
\item \textbf{Fabrication} : la matière fondue est moulée en \textbf{granulats} (matière première revendue à l'industrie) ou directement en objets nouveaux : pavés, tuiles, arrosoirs, chaises, seaux.
\end{enumerate}

\begin{popfigure}
\begin{center}
\begin{tikzpicture}[
  etape/.style={rectangle, rounded corners=3pt, draw=popblue, fill=popblueL, text width=1.9cm, align=center, minimum height=1cm, font=\small\bfseries},
  fleche/.style={-{Stealth[length=2.5mm]}, thick, poporange}
]
\node[etape] (a) at (0,0) {Collecte};
\node[etape] (b) at (2.6,0) {Tri par type};
\node[etape] (c) at (5.2,0) {Lavage};
\node[etape] (d) at (7.8,0) {Broyage};
\node[etape] (e) at (10.4,0) {Fonte};
\node[etape, fill=popgreenL, draw=popgreen] (f) at (13,0) {Produit fini};
\draw[fleche] (a) -- (b);
\draw[fleche] (b) -- (c);
\draw[fleche] (c) -- (d);
\draw[fleche] (d) -- (e);
\draw[fleche] (e) -- (f);
\node[below=2pt of f, font=\scriptsize, text width=2.2cm, align=center, popdark] {pavés, tuiles, arrosoirs, chaises};
\end{tikzpicture}
\end{center}
\legende{Chaîne de la filière de recyclage du plastique : six étapes successives conduisent du déchet plastique collecté au produit fini (granulats ou objets manufacturés).}
\end{popfigure}

\textbf{2. Le recyclage du papier et du carton.} Les vieux cahiers, journaux et cartons sont d'abord triés puis soumis au \cle{défibrage} : ils sont trempés dans l'eau et malaxés pour séparer les fibres de cellulose. On obtient une \cle{pâte à papier}, qui est ensuite égouttée, pressée et séchée en feuilles : c'est le \textbf{papier recyclé} ou le \textbf{carton recyclé}. Le papier peut être recyclé environ 5 à 7 fois, car les fibres raccourcissent à chaque cycle.

\textbf{3. Le recyclage des métaux.} Le fer (ferraille) et l'aluminium (canettes, casseroles usagées) sont collectés, triés puis \textbf{fondus} dans des fours à haute température. Le métal liquide est coulé en lingots puis \textbf{refabriqué} en objets neufs. C'est la filière la plus \textbf{rentable} : le métal se recycle presque indéfiniment sans perte de qualité, et recycler l'aluminium consomme environ 20 fois moins d'énergie que de le produire à partir du minerai (bauxite). C'est pourquoi les fers et métaux sont activement recherchés par les récupérateurs dans nos villes.

\textbf{4. Le recyclage du verre.} Deux voies existent :
\begin{itemize}
\item la \cle{consigne} : la bouteille vide est rapportée, lavée et \textbf{réutilisée directement} pour le même usage (système longtemps pratiqué pour les bouteilles de boissons au Cameroun) ;
\item la \cle{fusion} : le verre cassé (calcin) est fondu vers \SI{1500}{\celsius} pour fabriquer de \textbf{nouveaux récipients}. Le verre est recyclable à l'infini.
\end{itemize}

\textbf{5. La transformation artisanale (upcycling).} Dans nos quartiers et villages, l'ingéniosité artisanale valorise les déchets sans industrie lourde :
\begin{itemize}
\item les \textbf{pneus usagés} deviennent des sandales (les célèbres « pata-pata »), des meubles, des jardinières ou des balançoires ;
\item les \textbf{sachets plastiques} sont lavés, coupés en lanières puis \textbf{tressés} en sacs à provisions, nattes et chapeaux ;
\item les \textbf{bouteilles plastiques} deviennent des arrosoirs percés, des mangeoires pour volailles, ou, remplies de sable, des « briques écologiques » utilisées comme matériaux de construction de murs légers.
\end{itemize}

\courssub{Conditions de viabilité et limites du recyclage}

\begin{propriete}[Conditions d'une filière de recyclage viable]
Une filière de recyclage ne fonctionne durablement que si quatre conditions sont réunies :
\begin{enumerate}
\item une \textbf{collecte régulière} qui garantit un approvisionnement constant en déchets ;
\item un \textbf{tri de qualité} à la source (un plastique souillé ou mélangé à d'autres matières devient inutilisable) ;
\item un \textbf{débouché économique} : il faut des acheteurs pour les produits recyclés (pavés, granulats, pâte à papier) ;
\item une \textbf{main-d'œuvre formée} aux techniques de tri, de transformation et d'hygiène.
\end{enumerate}
Si l'une de ces conditions fait défaut, la filière s'effondre et les déchets retournent à la décharge.
\end{propriete}

\begin{attention}[Les limites du recyclage]
Le recyclage n'est pas une solution magique. Certains plastiques (sachets très fins, emballages multicouches) sont \textbf{techniquement non recyclables} ; les coûts de collecte et de traitement peuvent dépasser la valeur du produit obtenu ; enfin, la \textbf{fonte artisanale} du plastique ou des métaux, mal maîtrisée, dégage des fumées toxiques dangereuses pour les artisans et le voisinage. C'est pourquoi, dans la hiérarchie des 3R, la \textbf{réduction} des déchets à la source reste toujours la priorité : le meilleur déchet est celui qu'on ne produit pas.
\end{attention}

\begin{exemplebox}[Un chiffre qui fait réfléchir]
Selon les données industrielles de la filière papetière, recycler \textbf{1 tonne de papier} permet d'économiser environ \textbf{17 arbres}, ainsi que de grandes quantités d'eau et d'énergie. À l'échelle d'un lycée de 2\,000 élèves qui jetterait chacun 25~kg de papiers et cahiers par an, cela représente 50 tonnes de papier, soit près de \textbf{850 arbres épargnés} si ce papier est recyclé au lieu d'être brûlé ou enfoui. Ce résultat est admis : il illustre l'impact environnemental considérable d'un simple geste de tri.
\end{exemplebox}

\begin{savaistu}[Des pavés en plastique dans nos villes]
À Douala et à Yaoundé, plusieurs entreprises et coopératives fabriquent désormais des \textbf{pavés de rue} en mélangeant du \textbf{plastique fondu et du sable}. Ces pavés, moulés puis refroidis, se révèlent \textbf{plus résistants} à la compression et à l'usure que les pavés classiques en ciment, tout en étant moins coûteux. Chaque mètre carré de pavage valorise des centaines de sachets plastiques qui auraient sinon bouché les caniveaux et aggravé les inondations en saison des pluies. Le déchet devient ainsi une ressource pour la ville.
\end{savaistu}

\courssub{Choisir le bon mode de transformation : méthode et application}

\begin{methode}[Choisir le mode de transformation ou de recyclage adapté à un déchet]
Face à un déchet donné, raisonne en quatre étapes :
\begin{enumerate}
\item \textbf{Identifier la nature du déchet} : organique, plastique, papier, métal, verre, dangereux ? La nature détermine les techniques possibles (fonte pour le métal, défibrage pour le papier, compostage pour l'organique).
\item \textbf{Estimer la quantité disponible} : une petite quantité se prête à la transformation artisanale (upcycling) ; une grande quantité régulière justifie une filière industrielle.
\item \textbf{Vérifier l'existence d'un débouché local} : y a-t-il un acheteur ou un usage pour le produit obtenu (pavés, compost, ferraille) ?
\item \textbf{Choisir la technique la plus simple, la plus sûre et la plus rentable} compatible avec les trois réponses précédentes, en privilégiant toujours la réutilisation si elle est possible.
\end{enumerate}
\end{methode}

\begin{center}
\begin{tabularx}{\linewidth}{|l|X|X|}
\hline
\rowcolor{popblueL} \textbf{Déchet} & \textbf{Mode de transformation adapté} & \textbf{Produit obtenu} \\
\hline
Bouteilles plastiques & Broyage, fonte, moulage (ou découpe artisanale) & Granulats, pavés, arrosoirs, tuiles \\
\hline
Papiers, cartons & Défibrage, fabrication de pâte à papier & Papier et carton recyclés \\
\hline
Ferraille, canettes & Fonte en four, refabrication & Lingots, objets métalliques neufs \\
\hline
Bouteilles en verre & Consigne (réutilisation) ou fusion & Récipients neufs en verre \\
\hline
Pneus usagés & Découpe, assemblage (upcycling) & Sandales, meubles, jardinières \\
\hline
Sachets plastiques & Lavage, découpe, tressage & Sacs, nattes, chapeaux \\
\hline
Épluchures, restes de cuisine & Compostage ou méthanisation & Compost, biogaz \\
\hline
Piles, huiles usagées & Collecte spécialisée, traitement contrôlé & Matières valorisées ou neutralisées \\
\hline
\end{tabularx}
\end{center}

\begin{exoresolu}[Proposer un mode de recyclage adapté]
Énoncé. Le foyer des jeunes d'un quartier de Bafoussam collecte chaque semaine environ 300~kg de bouteilles plastiques et 50~kg de vieux papiers. Le président du foyer demande conseil : quel mode de transformation proposer pour chaque type de déchet, et pourquoi ?
\tcblower
Solution détaillée. Appliquons la méthode en quatre étapes.
\begin{enumerate}
\item \textbf{Nature des déchets} : il s'agit de plastique (bouteilles) et de papier, deux matières recyclables par des techniques différentes.
\item \textbf{Quantité} : 300~kg de plastique par semaine est une quantité importante et régulière, qui justifie une filière semi-industrielle ; 50~kg de papier peuvent être vendus ou transformés artisanalement.
\item \textbf{Débouché local} : à Bafoussam, les chantiers de construction achètent des pavés et tuiles ; les papetiers et usines achètent la pâte à papier ou le papier trié.
\item \textbf{Choix} : pour le plastique, on propose la chaîne complète \emph{tri par type $\rightarrow$ lavage $\rightarrow$ broyage $\rightarrow$ fonte $\rightarrow$ moulage} pour produire des \textbf{pavés et arrosoirs} revendus localement. Pour le papier, on propose la \textbf{revente à une filière de défibrage} ou la fabrication artisanale de carton recyclé.
\end{enumerate}
Conclusion : le plastique sera recyclé en pavés par fonte et moulage, le papier sera orienté vers la filière pâte à papier ; dans les deux cas, la régularité de la collecte et le débouché économique garantissent la viabilité du projet.
\end{exoresolu}

\begin{exercice}[À toi de jouer]
Dans ton village, on observe trois types de déchets abandonnés : des pneus usagés près du garage, des bouteilles en verre derrière les buvettes, et des épluchures de manioc au marché. Pour chaque déchet, propose un mode de transformation ou de valorisation adapté, en justifiant ton choix par la nature du déchet et le débouché local.
\end{exercice}

\begin{corrige}[Exercice]
\textbf{Pneus usagés} : caoutchouc non biodégradable, non recyclable industriellement sur place $\rightarrow$ transformation artisanale (upcycling) en sandales, meubles ou jardinières ; débouché : marché local de l'artisanat.
\textbf{Bouteilles en verre} : verre réutilisable et recyclable $\rightarrow$ priorité à la \textbf{consigne} (réutilisation directe après lavage) ; les bouteilles cassées peuvent être fondues pour fabriquer de nouveaux récipients.
\textbf{Épluchures de manioc} : déchet organique biodégradable $\rightarrow$ \textbf{compostage} pour produire un engrais naturel destiné aux champs, ou alimentation du bétail et méthanisation (biogaz). Dans chaque cas, le choix découle de la nature du déchet, de la quantité disponible et de l'existence d'un usage local du produit obtenu.
\end{corrige}

\begin{aretenir}[L'essentiel de la section]
\begin{itemize}
\item \textbf{Transformer} un déchet, c'est modifier sa forme ou sa nature pour obtenir un nouveau produit ; \textbf{recycler}, c'est le réintroduire comme matière première dans un cycle de production ; le \textbf{surcyclage} lui donne une valeur supérieure par l'artisanat.
\item Chaque matière a sa filière : fonte et moulage pour le \textbf{plastique}, défibrage pour le \textbf{papier}, fonte et refabrication pour les \textbf{métaux} (filière la plus rentable), consigne ou fusion pour le \textbf{verre}.
\item Une filière viable exige : collecte régulière, tri de qualité, débouché économique, main-d'œuvre formée.
\item Le recyclage a des limites (plastiques non recyclables, coûts, fumées toxiques) : la \textbf{réduction à la source} reste la priorité des 3R.
\end{itemize}
\end{aretenir}

\coursec{Valorisation des déchets organiques : compostage et biogaz}

Dans la section précédente, nous avons vu que la transformation et le recyclage permettent de redonner une seconde vie aux déchets secs : le plastique devient pavé ou objet neuf, le verre est refondu, le métal est refondu. Mais que faire des déchets \cle{organiques} — épluchures de manioc et de banane plantain, feuilles de légumes fanées, restes de repas, fumier des animaux — qui représentent souvent plus de la moitié des ordures d'une famille camerounaise ? Ces déchets, trop humides et trop putrescibles pour être recyclés comme le plastique, possèdent pourtant une richesse cachée : la matière organique qu'ils contiennent peut être transformée en \cle{engrais naturel} (le compost) ou en \cle{énergie} (le biogaz). C'est l'objet de cette section : comprendre comment deux procédés naturels, le compostage et la méthanisation, transforment un problème en ressource.

\courssub{Le compostage : transformer les déchets organiques en engrais naturel}

\textbf{Une observation de la vie courante.}\par

Derrière la case de grand-mère, au village, il y a souvent un coin où l'on jette les épluchures, les fanes de légumes et les feuilles mortes. Quelques mois plus tard, ce tas a diminué de volume et s'est transformé en une terre noire, friable et sans mauvaise odeur, que l'on mélange au sol de la plantation de macabo. Les tomates et les folong y poussent vigoureusement. Que s'est-il passé ? Comment des déchets malodorants ont-ils pu devenir un « engrais » gratuit ? C'est exactement ce que le compostage permet de contrôler et d'accélérer.

\begin{definition}[Compost et compostage]
Le \cle{compost} est un amendement organique, riche en \cle{humus} (matière organique stable), obtenu à partir de déchets organiques (épluchures, feuilles, herbes, restes de cuisine). Le \cle{compostage} est la transformation contrôlée des déchets organiques en compost par \textbf{décomposition aérobie}, c'est-à-dire en présence d'oxygène, sous l'action de micro-organismes (bactéries, champignons) et de décomposeurs (vers de terre, insectes).
\end{definition}

\textbf{Le mécanisme : qui fait le travail ?}\par

Le compostage n'est pas de la magie : c'est le travail de milliards d'êtres vivants. Les \cle{bactéries} et les \cle{champignons} attaquent les premiers la matière organique : ils digèrent les sucres, les protéines et la cellulose, et cette activité dégage de la \textbf{chaleur} — un tas de compost actif peut atteindre \SI{50-65}{\celsius} en son centre, ce qui détruit la plupart des germes pathogènes et des graines d'adventices. Puis les \cle{vers de terre} et autres petits décomposeurs fragmentent les résidus et les mélangent. Trois conditions sont indispensables :
\begin{itemize}
\item l'\cle{oxygène} : les micro-organismes du compostage sont aérobies ; sans air, le tas fermente et sent mauvais ;
\item l'\cle{humidité} : le tas doit rester humide comme une éponge essorée, ni sec ni détrempé ;
\item un bon équilibre \cle{carbone/azote} (rapport C/N) : les micro-organismes ont besoin à la fois de carbone (source d'énergie) et d'azote (pour synthétiser leurs protéines).
\end{itemize}

\begin{propriete}[Décomposition aérobie ou anaérobie (résultat admis)]
La décomposition des déchets organiques par les micro-organismes produit :
\begin{itemize}
\item du \textbf{compost} (humus) en \textbf{présence d'oxygène} (décomposition aérobie) : \ce{C6H10O5 + O2 -> CO2 + H2O + Energie + Humus} (équation globale simplifiée de la minéralisation de la cellulose) ;
\item du \textbf{méthane} (\ce{CH4}, gaz combustible) en \textbf{absence d'oxygène} (fermentation anaérobie / méthanisation) : \ce{C6H10O5 -> CH4 + CO2 + Digestat} (équation globale simplifiée).
\end{itemize}
Ce résultat de microbiologie est \emph{admis} au niveau Première : il n'est pas demandé d'en démontrer les mécanismes biochimiques détaillés, mais il faut savoir l'utiliser pour choisir le bon mode de valorisation.
\end{propriete}

\textbf{Les étapes du compostage.}\par

Pour obtenir un bon compost, on ne jette pas les déchets au hasard : on construit le tas en couches alternées.
\begin{enumerate}
\item \textbf{Collecte} des déchets organiques triés (épluchures, fanes, feuilles, herbes fauchées, marc de café).
\item \textbf{Alternance des couches} : une couche \emph{verte} (riche en azote : épluchures fraîches, herbes, fumier) puis une couche \emph{brune} (riche en carbone : feuilles sèches, paille, sciure, carton non plastifié), et ainsi de suite.
\item \textbf{Humidification} : on arrose légèrement chaque couche si elle est sèche.
\item \textbf{Brassage régulier} (toutes les 1 à 2 semaines) : retourner le tas à la fourche apporte de l'oxygène et homogénéise la décomposition.
\item \textbf{Maturation} : après 2 à 6 mois, le compost est mûr : il est brun foncé, friable, et sent la terre de forêt.
\end{enumerate}

\begin{popfigure}
\begin{center}
\begin{tikzpicture}[scale=1.0, every node/.style={font=\small}]
% Composteur : caisson
\draw[thick, popdark] (0,0) -- (0,4.4) -- (6,4.4) -- (6,0);
\draw[thick, popdark] (-0.4,0) -- (6.4,0);
% Couches alternées
\fill[popgreenL] (0.05,0.05) rectangle (5.95,0.95);
\fill[popgoldL] (0.05,0.95) rectangle (5.95,1.85);
\fill[popgreenL] (0.05,1.85) rectangle (5.95,2.75);
\fill[popgoldL] (0.05,2.75) rectangle (5.95,3.65);
\fill[popgreenL] (0.05,3.65) rectangle (5.95,4.35);
% Labels couches
\node[popdark] at (3,0.5) {couche brune (carbonée) : feuilles sèches, paille};
\node[popdark] at (3,1.4) {couche verte (azotée) : épluchures, herbes};
\node[popdark] at (3,2.3) {couche brune};
\node[popdark] at (3,3.2) {couche verte};
\node[popdark] at (3,4.0) {couche verte récente};
% Flèches d'air
\draw[-{Stealth}, thick, popblue] (-1.2,0.6) -- (-0.1,0.6);
\draw[-{Stealth}, thick, popblue] (-1.2,2.2) -- (-0.1,2.2);
\draw[-{Stealth}, thick, popblue] (-1.2,3.4) -- (-0.1,3.4);
\node[popblue, align=center] at (-1.6,2.2) {air\\ (\ce{O2})};
% Chaleur au centre
\draw[-{Stealth}, thick, poporange] (3,2.3) -- (3,3.9);
\node[poporange, right] at (3.1,4.6) {chaleur dégagée : \SI{50-65}{\celsius}};
% Maturation
\draw[-{Stealth}, thick, poppurple] (6.3,4.2) -- (6.3,0.3);
\node[poppurple, rotate=90, anchor=center] at (6.8,2.2) {maturation : 2 à 6 mois};
% Gouttes d'eau
\node[popblue] at (1,5.0) {humidification régulière};
\draw[-{Stealth}, popblue] (1,4.8) -- (1,4.45);
\end{tikzpicture}
\end{center}
\legende{Coupe schématique d'un composteur : les couches vertes (azotées) et brunes (carbonées) alternent ; l'air entre latéralement, l'activité microbienne dégage de la chaleur, et le compost mûrit en 2 à 6 mois.}
\end{popfigure}

\begin{methode}[Réaliser un compost familial ou scolaire en 5 étapes]
\begin{enumerate}
\item \textbf{Choisir l'emplacement} : un coin d'ombre, sur le sol nu (pour que les vers et micro-organismes remontent), à l'écart des habitations.
\item \textbf{Trier et collecter} : ne garder que les déchets organiques (épluchures, restes de légumes, feuilles, herbes, fumier) ; exclure plastiques, métaux, verre, os et viande en grande quantité.
\item \textbf{Monter le tas en couches} : alterner une couche verte (azotée) et une couche brune (carbonée) d'environ \SI{15}{cm} chacune, en humidifiant légèrement.
\item \textbf{Entretenir} : brasser le tas à la fourche toutes les 1 à 2 semaines et vérifier l'humidité (le tas doit rester humide, jamais sec ni détrempé).
\item \textbf{Récolter} : au bout de 2 à 6 mois, le compost mûr (brun, friable, odeur de terre) est épandu dans le jardin, la plantation ou les jardinières de l'école.
\end{enumerate}
\end{methode}

\begin{attention}[Erreurs fréquentes à éviter]
\begin{itemize}
\item Ne jamais mettre dans le compost des \textbf{plastiques, métaux, verre, piles} : ils ne se décomposent pas et polluent le compost.
\item Éviter les \textbf{os} et les gros restes de viande : ils se décomposent très lentement et attirent les rongeurs.
\item Ne pas laisser le compost \textbf{sécher complètement} : sans humidité, les micro-organismes meurent et la décomposition s'arrête. À l'inverse, un tas détrempé fermente et sent mauvais.
\item Ne pas oublier de \textbf{brasser} : sans oxygène, le tas passe en fermentation anaérobie et dégage des odeurs nauséabondes (\ce{H2S}, \ce{NH3}).
\end{itemize}
\end{attention}

\begin{exemplebox}[Un exemple chiffré]
Une famille produit \SI{10}{kg} de déchets de cuisine (épluchures, restes de légumes, marc de café). Après compostage et maturation, la majeure partie de l'eau s'est évaporée et une partie de la matière organique a été minéralisée en \ce{CO2} et en eau par les micro-organismes : on récolte environ \textbf{3 à 4 kg de compost}. Le volume des ordures à jeter diminue donc fortement, et l'on obtient gratuitement un amendement pour le jardin.
\end{exemplebox}

\textbf{Les intérêts du compost.}\par

Le compost est un \cle{amendement} du sol : il améliore sa structure (le sol devient plus meuble, retient mieux l'eau) et apporte des éléments nutritifs (azote, phosphore, potassium) libérés lentement. Son utilisation permet de \textbf{réduire les engrais chimiques}, coûteux et parfois polluants pour les nappes phréatiques (lessivage des nitrates), et de \textbf{diminuer le volume des ordures} envoyées en décharge. Au marché, les coopératives qui compostent les fanes et épluchures invendues réduisent à la fois les nuisances (odeurs, mouches) et leurs dépenses d'engrais.

\courssub{Le biogaz : produire de l'énergie à partir des déchets organiques}

\textbf{Une situation concrète.}\par

Dans une ferme de la région de l'Ouest, la cuisinière n'achète ni bois de chauffe ni gaz en bouteille : elle cuisine avec un gaz produit sur place, à partir des déjections de ses porcs. Comment des excréments peuvent-ils produire une flamme bleue ? La réponse tient en un mot : la \cle{fermentation anaérobie}.

\begin{definition}[Biogaz, fermentation anaérobie et digesteur]
Le \cle{biogaz} est un gaz combustible, riche en \cle{méthane} (\ce{CH4}, environ 50 à 70 \%), produit par la \cle{fermentation anaérobie} (ou méthanisation) des déchets organiques, c'est-à-dire leur décomposition par des bactéries \textbf{en l'absence d'oxygène}. Le \cle{digesteur} est la cuve fermée (étanche à l'air) dans laquelle s'effectue cette fermentation. Le résidu restant après fermentation, appelé \cle{digestat}, est un excellent engrais liquide.
\end{definition}

\textbf{Le principe de fonctionnement d'un digesteur domestique.}\par

Un digesteur domestique est une cuve enterrée en maçonnerie, totalement fermée pour empêcher l'air d'entrer. Son fonctionnement se décrit en quatre temps :
\begin{enumerate}
\item \textbf{Introduction de la matière} : on verse chaque jour par un tuyau d'entrée un mélange de déchets organiques (déjections de porcs ou de volailles, fumier, épluchures) et d'\textbf{eau}, en quantités à peu près égales (taux de matière sèche optimal).
\item \textbf{Fermentation} : dans la cuve, à l'abri de l'air et à la chaleur du climat (température mésophile \SI{30-40}{\celsius}), les bactéries anaérobies dégradent la matière organique en quatre phases (hydrolyse, acidogenèse, acétogenèse, méthanogenèse) et dégagent du biogaz, qui s'accumule dans la partie haute de la cuve.
\item \textbf{Récupération du gaz} : un tuyau relié au sommet du digesteur conduit le biogaz jusqu'à la cuisine, où il alimente un brûleur (cuisson) ou une lampe à gaz (éclairage).
\item \textbf{Évacuation du digestat} : le résidu liquide, débarrassé de la plupart des germes et des odeurs, s'écoule par un tuyau de sortie vers une fosse ; il est ensuite épandu au champ comme \textbf{engrais}.
\end{enumerate}

\begin{popfigure}
\begin{center}
\begin{tikzpicture}[scale=1.0, every node/.style={font=\small}]
% Sol
\draw[thick, popdark] (-0.5,0) -- (11.5,0);
\fill[pattern=north east lines, pattern color=popmuted] (-0.5,-0.35) rectangle (11.5,0);
% Cuve enterrée
\draw[thick, popdark, fill=popblueL] (2,-3.2) -- (2,-0.4) arc[start angle=180, end angle=0, radius=2] -- (6,-3.2) -- cycle;
% Liquide dans la cuve
\fill[popgreenL] (2.1,-3.1) rectangle (5.9,-1.6);
\draw[popgreen, thick] (2.1,-1.6) -- (5.9,-1.6);
\node[popdark] at (4,-2.4) {boue en fermentation};
\node[popdark] at (4,-0.85) {biogaz (\ce{CH4} + \ce{CO2})};
% Tuyau d'entrée
\draw[thick, popdark] (0.3,1.2) -- (0.3,-1.9) -- (2.05,-1.9);
\draw[-{Stealth}, thick, popgreen] (0.3,1.6) -- (0.3,1.25);
\node[popgreen, align=center, anchor=south west] at (-0.4,1.7) {entrée : déchets\\ organiques + eau};
% Tuyau de gaz vers cuisinière
\draw[thick, poporange] (4,0.0) -- (4,1.6) -- (8.6,1.6) -- (8.6,0.75);
\draw[-{Stealth}, thick, poporange] (8.6,1.6) -- (8.6,0.8);
% Cuisinière
\draw[thick, popdark, fill=popcream] (8.1,0) rectangle (9.1,0.55);
\draw[poporange, thick] (8.6,0.55) -- (8.6,0.75);
\draw[poporange, thick] (8.45,0.95) .. controls (8.5,0.75) and (8.7,0.75) .. (8.75,0.95);
\node[poporange, align=center] at (8.6,1.35) {flamme};
\node[popdark] at (8.6,-0.6) {cuisinière};
% Sortie digestat
\draw[thick, popdark] (5.95,-2.6) -- (7.4,-2.6) -- (7.4,-0.9);
\draw[-{Stealth}, thick, poppurple] (7.4,-0.9) -- (7.4,-0.45);
\node[poppurple, align=center] at (7.4,-1.5) [right]{sortie du digestat};
% Champ
\draw[thick, popgreen] (10.2,0) -- (10.2,0.5) (10.6,0) -- (10.6,0.6) (11.0,0) -- (11.0,0.45);
\draw[-{Stealth}, thick, poppurple] (7.7,-0.3) -- (10.0,-0.3);
\node[popgreen] at (10.6,0.95) {champ fertilisé};
% Label cuve
\node[popdark, align=center] at (4,-3.7) {cuve enterrée, étanche à l'air (digesteur)};
\end{tikzpicture}
\end{center}
\legende{Digesteur à biogaz domestique : les déchets organiques mélangés à l'eau entrent dans la cuve enterrée ; la fermentation anaérobie y produit du biogaz, conduit vers la cuisinière ; le digestat résiduel est récupéré et épandu comme engrais au champ.}
\end{popfigure}

\textbf{Pourquoi le biogaz est-il une double victoire ?}\par

Dans une décharge ou un tas d'ordures à l'air libre, les déchets organiques finissent par fermenter et dégager du méthane, un gaz à effet de serre au pouvoir de réchauffement global \textbf{28 fois supérieur} à celui du \ce{CO2} sur 100 ans, directement dans l'atmosphère. Le digesteur \textbf{capture} ce méthane et le transforme en énergie utile : le problème devient une ressource. De plus, le digestat remplace une partie des engrais achetés, et la cuisson au biogaz réduit la coupe de bois, donc la déforestation et les fumées nocives dans les cuisines (pollution de l'air intérieur).

\begin{savaistu}[Le biogaz au Cameroun]
Certaines fermes des régions de l'Ouest et du Nord-Ouest du Cameroun cuisinent \textbf{entièrement au biogaz} produit par les déjections de leurs porcs et de leurs volailles : quelques dizaines de kilogrammes de déjections par jour suffisent à alimenter les brûleurs d'une famille. Dans les zones rurales non électrifiées, le biogaz sert aussi à l'éclairage. De même, des coopératives valorisent les résidus de marchés (fanes, fruits avariés) par le compostage, et des projets pilotes transforment les déchets de scieries et de bananeraies (sciure, rejets de bananes) en compost ou en combustible solide (briquettes).
\end{savaistu}

\begin{exoresolu}[Choisir entre compostage et biogaz]
Énoncé. Une coopérative agricole de Bafoussam dispose chaque semaine de : (a) \SI{200}{kg} de feuilles sèches et de paille ; (b) \SI{300}{kg} d'épluchures de manioc et de légumes fanés ; (c) les déjections quotidiennes de 40 porcs. La coopérative veut à la fois améliorer la fertilité de ses champs et réduire sa facture de gaz en bouteille. Proposer un mode de valorisation adapté à chaque type de déchet, en justifiant.
\tcblower
Solution détaillée.
\begin{itemize}
\item \textbf{Déjections des porcs (c)} : elles sont fraîches, très humides et disponibles chaque jour en quantité régulière — conditions idéales pour un \textbf{digesteur à biogaz}. Mélangées à l'eau et fermentées sans oxygène, elles produiront du biogaz pour la cuisson (réduction de la facture de gaz) et un digestat utilisable comme engrais.
\item \textbf{Épluchures et légumes fanés (b)} : ce sont des déchets \emph{verts}, riches en azote. Ils seront \textbf{compostés} en couches alternées avec les feuilles sèches et la paille.
\item \textbf{Feuilles sèches et paille (a)} : déchets \emph{bruns}, riches en carbone ; ils constituent la couche carbonée indispensable à l'équilibre du compost (rapport C/N).
\end{itemize}
Conclusion : le biogaz répond au besoin d'\textbf{énergie}, le compostage répond au besoin de \textbf{fertilité des sols}. Les deux procédés sont complémentaires : le digestat du digesteur peut même rejoindre le tas de compost. On vérifie ainsi la règle d'or : à chaque déchet, sa filière de valorisation adaptée.
\end{exoresolu}

\courssub{Bilan : transformer un problème en ressources}

\begin{aretenir}[L'essentiel de la section]
\begin{itemize}
\item Le \cle{compostage} est la décomposition \textbf{aérobie} (avec oxygène) et contrôlée des déchets organiques ; il produit du \cle{compost} (amendement humique), en 2 à 6 mois, grâce aux bactéries, champignons et vers de terre, dans un tas en couches vertes/brunes alternées, humide et brassé régulièrement. Équation globale : \ce{Matiere organique + O2 -> Humus + CO2 + H2O + Chaleur}.
\item Le \cle{biogaz} est produit par la fermentation \textbf{anaérobie} (sans oxygène) des déchets organiques dans un \cle{digesteur} ; riche en méthane (\ce{CH4}), il sert à la cuisson et à l'éclairage, et son résidu (le \cle{digestat}) est un engrais liquide. Équation globale : \ce{Matiere organique -> CH4 + CO2 + Digestat}.
\item Ces deux procédés transforment un problème (déchets organiques encombrants, méthane des décharges, odeurs) en \textbf{ressources} : engrais naturels, énergie domestique, réduction des engrais chimiques et du volume des ordures.
\item Résultat admis à retenir : \textbf{avec oxygène} les déchets organiques donnent du compost ; \textbf{sans oxygène} ils donnent du méthane.
\end{itemize}
\end{aretenir}

\begin{exercice}[À toi de jouer]
Dans ton quartier, le marché produit chaque jour de grandes quantités de fanes de légumes, d'épluchures et de fruits avariés, actuellement brûlés ou abandonnés au bord de la route. Rédige un court texte de sensibilisation (8 à 10 lignes) à destination des commerçants, expliquant deux modes de valorisation possibles de ces déchets, leurs conditions de réussite et leurs avantages pour le marché et pour l'environnement.
\end{exercice}

\begin{corrige}[Exercice : À toi de jouer]
Éléments attendus dans le texte de sensibilisation :
\begin{itemize}
\item \textbf{Compostage} : trier les déchets organiques du marché, les disposer en couches alternées (épluchures et fanes = couches vertes ; feuilles sèches et cartons = couches brunes), humidifier et brasser régulièrement ; en 2 à 6 mois on obtient du compost, vendable ou utilisable par les maraîchers, ce qui réduit les ordures, les odeurs et les mouches.
\item \textbf{Biogaz} : installer un digesteur alimenté par les déchets organiques mélangés à l'eau ; la fermentation sans oxygène produit du biogaz (cuisson, éclairage du marché) et un digestat-engrais.
\item \textbf{Avantages} : marché plus propre et plus sain (moins de mouches et de maladies), revenus supplémentaires (vente du compost), économies d'énergie, et protection de l'environnement (moins de méthane et de fumées de brûlage).
\end{itemize}
Le texte doit employer un vocabulaire précis (compost, fermentation anaérobie, digesteur, digestat) et un ton adapté à une campagne de sensibilisation.
\end{corrige}

\coursec{Agir pour une gestion responsable des déchets}

Dans la section précédente, tu as découvert comment les déchets organiques peuvent être transformés en compost ou en biogaz au lieu de pourrir dans la nature. Mais ces techniques ne fonctionnent que si quelqu'un les met en œuvre : un ménage qui trie, une école qui composte, une commune qui organise la collecte. La gestion des déchets n'est donc pas seulement une affaire de technique : c'est d'abord une affaire de \cle{comportements} et d'\cle{organisation collective}. Cette dernière section répond à une question très concrète : que peux-tu faire, toi, ta famille, ton lycée, ton quartier, pour gérer durablement les déchets ?

\courssub{Responsabilité individuelle et collective}

\textbf{Situation concrète.} Un samedi matin à Yaoundé, après une forte pluie, un caniveau bouché par des sachets plastiques déborde : l'eau stagne dans la rue, les moustiques pullulent quelques jours plus tard, et deux cas de paludisme sont déclarés dans le quartier. Qui est responsable ? La mairie qui n'a pas curé le caniveau ? Les habitants qui ont jeté les sachets ? Les commerçants qui les distribuent ? En réalité, tout le monde : chaque acteur a une part de responsabilité.

\begin{definition}[Éco-citoyenneté]
L'\cle{éco-citoyenneté} est l'ensemble des comportements et des engagements d'un citoyen qui assume sa responsabilité vis-à-vis de l'environnement : il réduit sa production de déchets, les trie, participe aux actions collectives de propreté et respecte les règles de gestion des déchets de sa communauté.
\end{definition}

Le principe fondamental est simple : \textbf{chaque citoyen est à la fois producteur de déchets et acteur de leur gestion}. Un Camerounais urbain produit en moyenne environ \num{0,5} à \SI{1}{kg} de déchets par jour. Pour une ville comme Douala, cela représente des milliers de tonnes quotidiennes. Aucune entreprise de collecte ne peut tout gérer seule : la gestion durable des déchets repose sur le partage des responsabilités entre tous les acteurs.

\begin{popfigure}
\begin{center}
\begin{tikzpicture}[scale=1, every node/.style={font=\small}]
\node[draw, fill=popgreen!25, rounded corners, thick, minimum width=3.4cm, minimum height=1.1cm, align=center, font=\bfseries] (centre) at (0,0) {Gestion durable\\des déchets};
\node[draw, fill=popblueL, rounded corners, align=center] (menage) at (-4.2,2.2) {Ménage\\(tri, compostage)};
\node[draw, fill=popblueL, rounded corners, align=center] (ecole) at (0,3.1) {École\\(club environnement)};
\node[draw, fill=popblueL, rounded corners, align=center] (commune) at (4.2,2.2) {Commune\\(collecte, points de dépôt)};
\node[draw, fill=poporangeL, rounded corners, align=center] (entreprise) at (-4.2,-2.2) {Entreprise de collecte\\(HYSACAM)};
\node[draw, fill=poporangeL, rounded corners, align=center] (recycleurs) at (0,-3.1) {Recycleurs\\(coopératives, artisans)};
\node[draw, fill=poporangeL, rounded corners, align=center] (etat) at (4.2,-2.2) {État\\(MINEPDED, MINESEC)};
\draw[thick, popline] (centre) -- (menage);
\draw[thick, popline] (centre) -- (ecole);
\draw[thick, popline] (centre) -- (commune);
\draw[thick, popline] (centre) -- (entreprise);
\draw[thick, popline] (centre) -- (recycleurs);
\draw[thick, popline] (centre) -- (etat);
\end{tikzpicture}
\end{center}
\legende{Carte mentale des acteurs de la gestion des déchets. Aucun acteur ne peut réussir seul : la gestion durable exige la coordination du ménage, de l'école, de la commune, des entreprises de collecte, des recycleurs et de l'État.}
\end{popfigure}

\courssub{Agir à l'échelle du ménage}

Le ménage est la première unité de production des déchets : c'est donc là que tout commence. Quatre gestes essentiels :

\begin{enumerate}
\item \textbf{Trier à la source} : séparer les déchets organiques (épluchures, restes de repas), les plastiques, le verre, les métaux et les déchets dangereux dans des récipients distincts. Le tri à la source est beaucoup plus efficace que le tri après mélange, car les déchets mélangés se contaminent mutuellement.
\item \textbf{Composter} les déchets organiques dans un coin du jardin ou dans un composteur familial (voir section 5) : cela réduit d'environ un tiers le volume des ordures à évacuer.
\item \textbf{Réduire les emballages} : préférer les sacs réutilisables (sacs en tissu, paniers), acheter en vrac quand c'est possible, refuser les sachets plastiques superflus au marché.
\item \textbf{Déposer les déchets dangereux en points de collecte} : piles usagées, ampoules, médicaments périmés, huiles de vidange et peintures ne doivent jamais finir dans la poubelle ordinaire ni dans la nature, car ils contiennent des substances toxiques (métaux lourds, solvants) qui contaminent les sols et les nappes phréatiques.
\end{enumerate}

\begin{attention}[Le danger invisible des déchets dangereux]
Une pile jetée dans la nature peut polluer environ \SI{1}{m^3} de terre pendant des dizaines d'années avec du mercure, du cadmium ou du plomb. Ne jamais brûler les piles ni les plastiques : leur combustion libère des gaz toxiques (dioxines) dangereux pour les voies respiratoires.
\end{attention}

\courssub{Agir à l'échelle de l'école}

Le lycée est un lieu privilégié d'apprentissage de l'éco-citoyenneté. Voici les actions concrètes qu'un établissement peut mener :

\begin{itemize}
\item \textbf{Créer un club environnement} : une association d'élèves encadrée par un enseignant (souvent le professeur de SVT), qui organise les activités de sensibilisation et de propreté.
\item \textbf{Installer des poubelles de tri} : au minimum deux bacs clairement étiquetés — « organique » et « plastique/autres » — dans la cour, les couloirs et près du réfectoire.
\item \textbf{Organiser des journées de nettoyage} : ramassage des déchets dans l'enceinte du lycée et dans le quartier environnant, avec gants et sacs fournis.
\item \textbf{Installer un composteur scolaire} : les restes du réfectoire et les feuilles mortes y sont déposés et transformés en compost en quelques mois.
\item \textbf{Créer un jardin pédagogique alimenté au compost} : le compost produit fertilise un potager scolaire (tomates, légumes-feuilles, maïs), ce qui boucle le cycle et sert de support concret aux cours de SVT.
\end{itemize}

\begin{exemplebox}[Une journée de nettoyage, ça pèse lourd]
Lors d'une journée de nettoyage d'un grand marché urbain au Cameroun, quelques centaines de volontaires (élèves, commerçants, associations) peuvent collecter \textbf{plusieurs tonnes de déchets} en une matinée. Sur un chantier typique de ce type, environ la moitié du volume ramassé est constituée de déchets organiques valorisables par compostage, et une fraction notable de plastiques revendables aux recycleurs. Une partie des déchets collectés n'est donc pas une perte, mais une \cle{ressource}.
\end{exemplebox}

\begin{savaistu}[Le compost paie les activités du lycée]
Plusieurs lycées camerounais, notamment dans les régions de l'Ouest et du Centre, financent une partie de leurs activités (matériel sportif, récompenses, excursions) grâce à la \textbf{vente du compost} produit par leur composteur scolaire et à la \textbf{vente des plastiques triés} aux coopératives de recyclage. Les déchets deviennent ainsi une source de revenus : c'est l'économie circulaire en action, à l'échelle d'un établissement scolaire.
\end{savaistu}

\courssub{Agir à l'échelle de la communauté}

Au-delà du ménage et de l'école, la communauté (quartier, village, commune) dispose de leviers puissants :

\begin{itemize}
\item \textbf{Les associations et comités de quartier} organisent les nettoyages collectifs, négocient l'emplacement des points de collecte et relaient les campagnes de sensibilisation.
\item \textbf{Les coopératives de recyclage} regroupent des ramasseurs et trieurs qui collectent plastiques, métaux et cartons, les conditionnent et les revendent aux filières de recyclage. Elles créent des emplois tout en assainissant l'environnement.
\item \textbf{Les partenariats avec la mairie et les entreprises de collecte} permettent d'organiser des points de dépôt réguliers, la desserte des quartiers éloignés et la gestion des décharges.
\end{itemize}

\textbf{Cadre institutionnel au Cameroun.} La gestion des déchets s'inscrit dans un cadre précis :
\begin{itemize}
\item les \textbf{communes} sont responsables de l'organisation de la collecte et de la propreté urbaine ;
\item \textbf{HYSACAM} (Hygiène et Salubrité du Cameroun) est la principale entreprise chargée de la collecte et du traitement des ordures dans les grandes villes ;
\item le \textbf{MINEPDED} (Ministère de l'Environnement, de la Protection de la Nature et du Développement Durable) définit les politiques environnementales et encadre la gestion des déchets ;
\item le \textbf{MINESEC} (Ministère des Enseignements Secondaires) intègre l'éducation à l'environnement dans les programmes scolaires — cette leçon en est la preuve ;
\item les \textbf{journées nationales de propreté} mobilisent périodiquement citoyens, écoles et administrations pour des nettoyages collectifs.
\end{itemize}

\courssub{Sensibiliser : techniques et construction d'un message efficace}

\begin{definition}[Sensibilisation]
La \cle{sensibilisation} est l'ensemble des actions de communication visant à informer un public, à modifier ses attitudes et à l'amener à adopter des comportements favorables à un objectif donné — ici, la gestion responsable des déchets.
\end{definition}

\textbf{Techniques de sensibilisation.} Selon le public visé, on peut utiliser :
\begin{itemize}
\item des \textbf{affiches} placardées aux endroits stratégiques (marchés, écoles, chefferies) ;
\item des \textbf{sketches et saynètes} joués lors de cérémonies ou de journées portes ouvertes ;
\item des \textbf{chansons et slogans} en français et en langues locales, faciles à mémoriser ;
\item les \textbf{réseaux sociaux} (WhatsApp, Facebook) pour toucher les jeunes et diffuser photos et vidéos ;
\item le \textbf{porte-à-porte}, très efficace dans les quartiers pour expliquer directement les gestes du tri ;
\item les \textbf{démonstrations pratiques} : montrer un composteur en fonctionnement vaut mieux qu'un long discours.
\end{itemize}

\begin{methode}[Construire un message de sensibilisation en 4 points]
Un message de sensibilisation efficace se construit toujours en quatre étapes :
\begin{enumerate}
\item \textbf{Un problème précis} : identifier un seul problème, local et observable (ex. : « les sachets plastiques bouchent le caniveau du quartier et l'eau stagne »).
\item \textbf{Un public cible} : définir à qui l'on s'adresse (ménagères, élèves, commerçants du marché) pour adapter le langage et le support.
\item \textbf{Une solution concrète} : proposer un geste simple, réalisable localement et immédiatement (ex. : « utilise un panier ou un sac en tissu pour tes courses »).
\item \textbf{Un appel à l'action} : formuler une invitation claire, datée si possible (ex. : « rejoins la journée de nettoyage samedi à 7 h devant le marché »).
\end{enumerate}
\end{methode}

\begin{attention}[Erreur fréquente : sensibiliser sans solution]
Le piège classique est de dénoncer un problème (« il y a trop de déchets partout ! ») sans proposer de \textbf{solution concrète et réalisable localement}. Un tel message décourage au lieu de mobiliser. Proposer de « construire une usine de recyclage » à des élèves est irréaliste ; proposer de trier les sachets plastiques dans un bac dédié et de les vendre à une coopérative est faisable dès demain. Toujours vérifier : le geste proposé est-il simple, gratuit ou peu coûteux, et possible ici et maintenant ?
\end{attention}

\begin{popfigure}
\begin{center}
\begin{tikzpicture}
\draw[fill=popcream, draw=popdark, thick, rounded corners] (0,0) rectangle (10,7);
\node[font=\bfseries\large, text=popdark] at (5,6.4) {STOP AUX SACHETS DANS LE CANIVEAU !};
\draw[fill=popink!15, draw=popink, rounded corners] (0.4,4.2) rectangle (9.6,5.8);
\node[align=left, text width=8.6cm, font=\small] at (5,5.0) {\textbf{1. Le problème :} les sachets plastiques jetés bouchent le caniveau du quartier ; l'eau stagne et les moustiques se multiplient.};
\draw[fill=poporangeL, draw=poporange, rounded corners] (0.4,2.3) rectangle (9.6,3.9);
\node[align=left, text width=8.6cm, font=\small] at (5,3.1) {\textbf{2. La solution :} jette tes sachets dans la poubelle prévue devant le marché, ou utilise un panier réutilisable.};
\draw[fill=popgreenL, draw=popgreen, rounded corners] (0.4,0.4) rectangle (9.6,2.0);
\node[align=left, text width=8.6cm, font=\small] at (5,1.2) {\textbf{3. Agis :} journée de nettoyage samedi à 7 h, place du marché. Gants et sacs fournis. Un quartier propre, c'est la santé pour tous !};
\end{tikzpicture}
\end{center}
\legende{Maquette d'affiche de sensibilisation structurée en trois blocs : problème précis et local (en haut), solution concrète et réalisable (au milieu), appel à l'action daté (en bas).}
\end{popfigure}

\courssub{Complément utile au Probatoire : relier les déchets aux Objectifs de Développement Durable}

\begin{aretenir}[Gestion des déchets et ODD — complément au programme]
Les \textbf{Objectifs de Développement Durable} (ODD), adoptés par les Nations unies en 2015, fixent 17 objectifs mondiaux pour 2030. La gestion responsable des déchets contribue directement à trois d'entre eux :
\begin{itemize}
\item \textbf{ODD 11 — Villes et communautés durables} : réduire l'impact environnemental des villes, notamment grâce à une meilleure gestion des déchets ;
\item \textbf{ODD 12 — Consommation et production responsables} : réduire, réutiliser et recycler les déchets (la règle des 3R !) ;
\item \textbf{ODD 13 — Lutte contre le changement climatique} : le compostage et le biogaz évitent les émissions de méthane des décharges, un puissant gaz à effet de serre.
\end{itemize}
Au Probatoire, savoir citer ces liens montre ta capacité à situer la gestion des déchets dans le développement durable : c'est un atout pour la rédaction.
\end{aretenir}

\begin{exoresolu}[Construire un message de sensibilisation]
\textbf{Énoncé.} Ton lycée constate que de nombreux élèves jettent des sachets et des papiers dans la cour, malgré la présence de poubelles. Le club environnement te charge de concevoir un message de sensibilisation. Construis ce message en suivant la méthode en 4 points, puis rédige le slogan final.
\tcblower
\textbf{Solution.}
\begin{enumerate}
\item \textbf{Problème précis} : les sachets et papiers jetés dans la cour du lycée salissent l'environnement scolaire, bouchent les évacuations d'eau de pluie et attirent mouches et moustiques.
\item \textbf{Public cible} : les élèves du lycée, notamment ceux qui achètent des friandises et boissons à la pause. Le message doit être court, en langage simple, avec un support visuel (affiche près des points de vente).
\item \textbf{Solution concrète} : trois poubelles de tri étiquetées sont installées près du réfectoire et de la cour ; chaque élève y jette ses sachets et papiers. Les déchets organiques iront au composteur du jardin pédagogique.
\item \textbf{Appel à l'action} : « À partir de lundi, jette ton sachet dans le bac vert. Une cour propre commence par ton geste ! »
\end{enumerate}
\textbf{Slogan final} : « Un sachet dans le bac, un lycée propre, une santé protégée ! »
\textbf{Vérification (piège évité)} : la solution proposée est concrète (des bacs existants), gratuite pour l'élève et réalisable immédiatement — le message ne se contente pas de dénoncer.
\end{exoresolu}

\begin{exercice}[À toi de jouer]
Ton quartier souffre de dépôts sauvages d'ordures sur un terrain vague, à proximité d'un point d'eau. 1) Identifie deux risques pour la santé des habitants. 2) Propose une action à l'échelle du ménage, une action à l'échelle de la communauté et une action à mener avec la mairie. 3) Rédige un message de sensibilisation complet (problème, cible, solution, appel à l'action) destiné aux habitants du quartier.
\end{exercice}

\begin{corrige}[Exercice]
\textbf{1) Risques sanitaires} : la stagnation des eaux sur les déchets favorise la prolifération des moustiques (paludisme) ; les déchets en décomposition attirent mouches et rongeurs, vecteurs de maladies diarrhéiques (choléra, typhoïde) ; la proximité du point d'eau expose à la contamination de l'eau de boisson.
\textbf{2) Actions} : au niveau du ménage, trier les déchets et les déposer uniquement aux points de collecte officiels, composter les déchets organiques. Au niveau de la communauté, organiser avec l'association du quartier une journée de nettoyage du terrain vague et une surveillance collective contre les dépôts sauvages. Avec la mairie, demander l'installation d'un point de collecte desservi par l'entreprise de ramassage et l'aménagement du terrain (jardin communautaire, espace propre).
\textbf{3) Message} : Problème — les ordures du terrain vague menacent notre point d'eau et la santé de nos enfants. Cible — les habitants du quartier. Solution — déposer ses ordures au bac installé à l'entrée de la rue et composter les déchets de cuisine. Appel à l'action — « Grand nettoyage du terrain vague samedi à 6 h 30 : viens avec ta machette et tes gants. Un quartier propre, une eau saine ! »
\end{corrige}

\begin{aretenir}[L'essentiel de la section]
\begin{itemize}
\item Chaque citoyen est \textbf{producteur de déchets et acteur de leur gestion} : c'est le principe de l'\cle{éco-citoyenneté}.
\item On agit à trois échelles : le \textbf{ménage} (tri, compostage, réduction des emballages, dépôt des déchets dangereux), l'\textbf{école} (club environnement, poubelles de tri, composteur, jardin pédagogique) et la \textbf{communauté} (associations, coopératives, partenariats avec la mairie et HYSACAM).
\item Un message de \cle{sensibilisation} efficace comporte toujours : un \textbf{problème précis}, un \textbf{public cible}, une \textbf{solution concrète et réalisable localement}, un \textbf{appel à l'action}.
\item La gestion des déchets s'inscrit dans les \textbf{ODD 11, 12 et 13} : villes durables, consommation responsable, lutte contre le changement climatique.
\end{itemize}
\end{aretenir}

\newpage

\coursec{Activité d'intégration}

\coursec{Transformation et recyclage des déchets}

\courssub{Activité d'intégration : Audit et plan de valorisation des déchets du lycée}

\courssub{Objectifs de l'activité}

À l'issue de cette activité d'intégration, tu seras capable de :
\begin{itemize}
\item \cle{trier} et \cle{classifier} des déchets réels ou simulés selon leur nature (organiques, plastiques, papiers, dangereux) ;
\item \cle{calculer} des proportions (pourcentages) et les représenter dans un tableau ou un diagramme ;
\item \cle{proposer} pour chaque catégorie de déchet un mode de transformation ou de recyclage adapté (compostage, biogaz, recyclage, point de collecte spécialisé) ;
\item \cle{rédiger} un message de sensibilisation structuré (problème, cible, solution, appel à l'action) ;
\item \cle{argumenter} et défendre un plan de gestion réaliste devant la classe.
\end{itemize}

Cette activité mobilise l'ensemble des savoir-faire de la leçon dans une situation authentique : la gestion des déchets de ton propre lycée.

\courssub{Situation}

\textbf{Le lycée de Mfoundi face à ses poubelles.}

Le lycée bilingue de Mfoundi, à Yaoundé, accueille \num{1800} élèves et dispose d'une cantine, d'un terrain de sport et de salles d'informatique. Chaque jour, les poubelles débordent : restes de repas, sachets plastiques, feuilles de papier, et même de vieilles piles et ampoules jetées pêle-mêle. Les ordures sont brûlées dans un coin de la cour, ce qui dégage une fumée âcre, ou ramassées irrégulièrement par HYSACAM, l'entreprise d'hygiène et salubrité du Cameroun. Le proviseur, inquiet des risques sanitaires (moustiques, mouches, maladies diarrhéiques) et environnementaux, lance un concours entre les classes de Première : \emph{« Élaborez le meilleur plan de gestion et de valorisation des déchets du lycée. »}

Un groupe d'élèves a réalisé un \cle{audit} : pendant une journée type, il a collecté, trié et pesé tous les déchets produits. Voici les résultats :

\begin{center}
\begin{tabularx}{\linewidth}{|l|X|c|}
\hline
\rowcolor{popblueL}
\textbf{Nature du déchet} & \textbf{Exemples observés} & \textbf{Masse par jour} \\
\hline
Déchets organiques & Restes de repas de la cantine (ndolé, riz, arachides), épluchures de manioc et de banane plantain, herbes fauchées & \SI{40}{kg} \\
\hline
Déchets plastiques & Sachets d'eau et de jus, bouteilles en plastique, emballages de beignets & \SI{12}{kg} \\
\hline
Papiers et cartons & Vieilles feuilles, cahiers usagés, cartons d'emballage de la cantine & \SI{8}{kg} \\
\hline
Déchets dangereux & Piles usagées, ampoules grillées, cartouches d'encre vides & \SI{2}{kg} \\
\hline
\end{tabularx}
\end{center}

\textbf{Données complémentaires utiles :}
\begin{itemize}
\item Un récupérateur de quartier achète le plastique propre et trié à \SI{50}{FCFA} le kilogramme et le papier-carton à \SI{25}{FCFA} le kilogramme.
\item Le compost mûr peut être vendu \SI{100}{FCFA} le kilogramme aux jardiniers du quartier ; le compostage transforme environ $60\,\%$ de la masse organique en compost utilisable.
\item Le lycée fonctionne \num{200} jours par an.
\end{itemize}

\textbf{Ta mission :} avec ton groupe, élaborer un plan complet de tri, de valorisation et de sensibilisation, puis le présenter en 3 minutes devant la classe. Le meilleur plan sera proposé à la direction.

\courssub{Consignes}

\begin{enumerate}
\item \textbf{Classer.} À partir du tableau de l'audit, classe chaque déchet cité dans l'une des catégories suivantes : déchet organique, déchet recyclable (plastique, papier), déchet dangereux. Justifie ton classement pour les piles usagées.
\item \textbf{Calculer.} Calcule la masse totale de déchets produite par jour, puis le pourcentage de chaque catégorie. Présente les résultats dans un tableau à trois colonnes (catégorie, masse, pourcentage).
\item \textbf{Représenter.} Construis un diagramme circulaire (camembert) ou un diagramme en bâtons représentant la répartition des déchets. Indique comment tu passes du pourcentage à l'angle du secteur.
\item \textbf{Valoriser.} Pour chaque catégorie, propose un mode de transformation ou de recyclage adapté et réaliste dans le contexte du lycée. Calcule le revenu annuel potentiel de la vente des plastiques, des papiers et du compost.
\item \textbf{Sensibiliser.} Rédige le texte d'une affiche de sensibilisation destinée aux élèves, structurée en quatre parties : le problème, la cible, la solution, l'appel à l'action.
\item \textbf{Présenter et comparer.} Présente ton plan en 3 minutes. La classe compare les plans selon trois critères : réalisme, coût, impact environnemental, puis élit le meilleur.
\end{enumerate}

\courssub{Ressources à mobiliser}

\begin{aretenir}[Boîte à outils pour l'activité]
\begin{itemize}
\item \textbf{Typologie des déchets :} organiques (biodégradables), recyclables (plastiques, papiers, métaux, verre), dangereux (piles, ampoules, produits chimiques — toxiques pour le sol, l'eau et la santé).
\item \textbf{Règle des 3R :} \cle{Réduire} à la source, \cle{Réutiliser}, \cle{Recycler}.
\item \textbf{Valorisation de la matière organique :} le \cle{compostage} (décomposition aérobie contrôlée produisant un engrais naturel) et la méthanisation (production de \cle{biogaz} en anaérobie).
\item \textbf{Calcul d'un pourcentage :} $\text{pourcentage} = \dfrac{\text{masse de la catégorie}}{\text{masse totale}} \times 100$.
\item \textbf{Angle d'un secteur circulaire :} $\text{angle} = \dfrac{\text{pourcentage}}{100} \times 360^{\circ}$.
\item \textbf{Revenu annuel :} revenu journalier $\times$ nombre de jours de fonctionnement.
\end{itemize}
\end{aretenir}

\courssub{Démarche et corrigé commenté}

\begin{exoresolu}[Audit et plan de valorisation du lycée de Mfoundi]
\textbf{Énoncé.} À partir des données de l'audit (40 kg d'organiques, 12 kg de plastiques, 8 kg de papiers, 2 kg de dangereux ; 200 jours/an ; prix de rachat : plastique 50 FCFA/kg, papier 25 FCFA/kg, compost 100 FCFA/kg avec un rendement de compostage de $60\,\%$), réalise le classement, calcule les pourcentages, propose les filières de valorisation et estime le revenu annuel.

\tcblower

\textbf{Étape 1 — Classement des déchets.}

On applique la typologie vue en cours :
\begin{itemize}
\item \emph{Organiques :} restes de repas, épluchures, herbes fauchées (biodégradables).
\item \emph{Recyclables :} sachets et bouteilles plastiques, papiers et cartons.
\item \emph{Dangereux :} piles, ampoules, cartouches d'encre. \textbf{Justification :} les piles contiennent des métaux lourds (mercure, plomb, cadmium) qui, jetés dans la nature ou brûlés, contaminent les sols et les nappes phréatiques et provoquent des intoxications. Elles ne doivent jamais rejoindre la poubelle ordinaire ni le feu.
\end{itemize}

\textbf{Étape 2 — Calcul des pourcentages.}

Masse totale journalière :
\[ m_{\text{totale}} = 40 + 12 + 8 + 2 = \SI{62}{kg}. \]

\begin{align*}
\text{Organiques} &= \frac{40}{62}\times 100 \approx 64{,}5\,\% &
\text{Plastiques} &= \frac{12}{62}\times 100 \approx 19{,}4\,\% \\
\text{Papiers} &= \frac{8}{62}\times 100 \approx 12{,}9\,\% &
\text{Dangereux} &= \frac{2}{62}\times 100 \approx 3{,}2\,\%
\end{align*}

\begin{center}
\begin{tabularx}{\linewidth}{|l|X|c|c|}
\hline
\rowcolor{popblueL}
\textbf{Catégorie} & \textbf{Masse (kg/jour)} & \textbf{Pourcentage} & \textbf{Angle du secteur} \\
\hline
Organiques & 40 & $64{,}5\,\%$ & $232^{\circ}$ \\
\hline
Plastiques & 12 & $19{,}4\,\%$ & $70^{\circ}$ \\
\hline
Papiers & 8 & $12{,}9\,\%$ & $46^{\circ}$ \\
\hline
Dangereux & 2 & $3{,}2\,\%$ & $12^{\circ}$ \\
\hline
\textbf{Total} & \textbf{62} & \textbf{$100\,\%$} & \textbf{$360^{\circ}$} \\
\hline
\end{tabularx}
\end{center}

\textbf{Étape 3 — Représentation.} Pour le diagramme circulaire, chaque angle se calcule par $\text{angle} = \frac{\%}{100}\times 360^{\circ}$ : par exemple $\frac{64{,}5}{100}\times 360 \approx 232^{\circ}$ pour les organiques. On reporte les secteurs au rapporteur, du plus grand au plus petit, avec une couleur par catégorie.

\begin{popfigure}
\begin{tikzpicture}
\fill[popgreenL] (0,0) -- (90:2) arc (90:322:2) -- cycle;
\fill[poporangeL] (0,0) -- (322:2) arc (322:392:2) -- cycle;
\fill[popblueL] (0,0) -- (32:2) arc (32:78:2) -- cycle;
\fill[poppinkL] (0,0) -- (78:2) arc (78:90:2) -- cycle;
\draw (0,0) circle (2);
\draw[popdark] (206:2.4) node {\footnotesize Organiques $64{,}5\,\%$};
\draw[popdark] (-3:2.6) node {\footnotesize Plastiques $19{,}4\,\%$};
\draw[popdark] (55:2.7) node {\footnotesize Papiers $12{,}9\,\%$};
\draw[popdark] (84:2.7) node {\footnotesize Dangereux $3{,}2\,\%$};
\end{tikzpicture}
\legende{Diagramme circulaire de la répartition des déchets du lycée par jour. Les déchets organiques dominent largement : c'est la filière prioritaire de valorisation.}
\end{popfigure}

\textbf{Étape 4 — Proposition de valorisation et calcul du revenu annuel.}

\begin{itemize}
\item \emph{Organiques (40 kg/jour) :} mise en place d'un \textbf{composteur} derrière la cantine (alterner couches vertes et brunes, humidifier, retourner chaque semaine). Masse de compost par jour : $40 \times 0{,}60 = \SI{24}{kg}$. Revenu annuel : $24 \times 200 \times 100 = \SI{480000}{FCFA}$.
\item \emph{Plastiques (12 kg/jour) :} tri dans une poubelle jaune, vente au récupérateur. Revenu annuel : $12 \times 200 \times 50 = \SI{120000}{FCFA}$.
\item \emph{Papiers (8 kg/jour) :} poubelle bleue, vente au récupérateur. Revenu annuel : $8 \times 200 \times 25 = \SI{40000}{FCFA}$.
\item \emph{Dangereux (2 kg/jour) :} \textbf{point de collecte sécurisé} (caisse fermée étiquetée « Piles et ampoules ») remise périodiquement à une filière spécialisée via HYSACAM. Aucun revenu, mais un gain sanitaire majeur.
\end{itemize}

Revenu annuel total potentiel :
\[ 480000 + 120000 + 40000 = \SI{640000}{FCFA}, \]
de quoi financer des fournitures scolaires ou l'entretien du jardin du lycée.

\textbf{Étape 5 — Affiche de sensibilisation (modèle).}
\begin{itemize}
\item \emph{Problème :} « Chaque jour, notre lycée produit 62 kg de déchets brûlés ou abandonnés : fumées toxiques, moustiques, maladies. »
\item \emph{Cible :} « Élèves et personnel du lycée de Mfoundi. »
\item \emph{Solution :} « Trois poubelles, trois gestes : vert pour le compost, jaune pour le plastique, bleu pour le papier. Les piles vont à la caisse rouge. »
\item \emph{Appel à l'action :} « Dès lundi, trie ton déchet : notre cour propre, notre santé, et 640 000 FCFA pour notre lycée ! »
\end{itemize}

\textbf{Étape 6 — Présentation.} Le plan est présenté en 3 minutes : résultats de l'audit (1 min), filières et revenus (1 min), affiche et appel à l'action (1 min). La classe vote selon les critères de réalisme, de coût et d'impact.
\end{exoresolu}

\begin{attention}[Pièges fréquents dans cette activité]
\begin{itemize}
\item Oublier d'arrondir proprement : la somme des pourcentages doit rester égale à $100\,\%$ (ajuste le dernier arrondi si nécessaire).
\item Mettre les piles au compost ou au feu : un déchet dangereux n'est \textbf{jamais} composté, recyclé avec les plastiques ni brûlé.
\item Confondre la masse de déchets organiques (40 kg) avec la masse de compost produit (24 kg) : applique toujours le rendement de $60\,\%$.
\item Proposer des solutions irréalistes (usine de recyclage, incinérateur) : un bon plan s'appuie sur des acteurs locaux existants (récupérateurs, HYSACAM, jardiniers du quartier).
\end{itemize}
\end{attention}

\courssub{Bilan de l'activité}

Cet audit a d'abord rendu visible une réalité chiffrée : près des \textbf{deux tiers} des déchets du lycée ($64{,}5\,\%$) sont organiques et donc \cle{valorisable} sur place par compostage, sans aucun transport. Il a ensuite montré que le \cle{tri à la source} est la condition de toute valorisation : un plastique souillé de restes de nourriture ne peut être ni vendu ni recyclé. L'activité a aussi révélé la dimension économique de la gestion des déchets : ce qui encombrait la cour peut rapporter environ \SI{640000}{FCFA} par an, illustrant concrètement l'économie circulaire. Enfin, la rédaction de l'affiche et la présentation orale ont développé ton savoir-être : convaincre ses camarades et la direction, c'est déjà agir en \emph{citoyen éco-responsable}. Le plan élu par la classe pourra être réellement proposé au proviseur : la leçon sort alors de la salle de classe pour transformer durablement ton environnement quotidien.

\newpage

\coursec{Méthodes et exercices résolus}

\coursec{Les déchets : définition, sources et typologie}

\begin{definition}[Déchet]
Tout résidu d'un processus de production, de transformation ou d'utilisation, toute substance, matière, produit ou plus généralement tout bien meuble abandonné ou que son destinataire a l'intention ou l'obligation d'abandonner (Loi n°96/12 du 5 août 1996 relative à la gestion de l'environnement au Cameroun).
\end{definition}

\begin{definition}[Typologie des déchets selon leur nature]
\begin{itemize}
\item \textbf{Déchets organiques (biodégradables)} : d'origine végétale ou animale, décomposables par les micro-organismes (épluchures, restes alimentaires, tontes, fumier, papier/carton non souillé).
\item \textbf{Déchets recyclables non dangereux} : matières valorisables par transformation industrielle : plastiques (PET, PEHD, PP), métaux (fer, aluminium), verre, papiers/cartons propres.
\item \textbf{Déchets dangereux} : présentent un risque pour la santé ou l'environnement du fait de leur toxicité, inflammabilité, corrosivité, réactivité ou infectiosité (piles, batteries, médicaments périmés, solvants, huiles de vidange, déchets d'activités de soins à risques infectieux -- DASRI).
\item \textbf{Déchets ultimes} : résidus non valorisables techniquement ni économiquement dans les conditions actuelles, destinés à l'enfouissement en centre de stockage technique (CST) ou à l'incinération avec récupération d'énergie.
\end{itemize}
\end{definition}

\begin{savaistu}[Contexte camerounais : HYSACAM et la gestion urbaine]
Au Cameroun, la société HYSACAM (Hygiène et Salubrité du Cameroun) assure la collecte et le traitement des ordures ménagères dans les grandes villes (Yaoundé, Douala, Bafoussam, Bamenda, Garoua, Maroua, Ngaoundéré). Les décharges contrôlées (ex. : Nkolfoulou à Yaoundé, PK10 à Douala) reçoivent les déchets mélangés. Le tri à la source reste marginal, bien que des initiatives locales (ONG, groupements de femmes, startups comme \emph{Kemit Ecology} pour le charbon écologique à Douala) développent le recyclage et la valorisation organique.
\end{savaistu}

\courssub{Savoir-faire 1 : Trier des déchets selon leur nature}

\begin{methode}[Réaliser un tri correct des déchets]
Pour classer un déchet dans la bonne catégorie, suis toujours la même démarche :
\begin{enumerate}
\item \textbf{Identifier la matière dominante} du déchet : matière organique (d'origine végétale ou animale), plastique, verre, métal, papier/carton, ou substance chimique.
\item \textbf{Se poser la question de la biodégradabilité} : le déchet peut-il être décomposé par les micro-organismes ? Si oui, c'est un déchet \cle{organique} (biodégradable).
\item \textbf{Rechercher un danger potentiel} : toxicité, inflammabilité, caractère corrosif ou infectieux. Si oui, c'est un déchet \cle{dangereux} (piles, médicaments périmés, produits d'entretien, déchets hospitaliers, huiles de vidange).
\item \textbf{Vérifier la recyclabilité} : plastiques (bouteilles, sachets), verre, métaux et papiers sont des déchets \cle{recyclables} non dangereux.
\item \textbf{Attribuer le déchet à une filière} : compostage/biogaz (organique), recyclage (plastique, verre, métal, papier), collecte spécialisée (dangereux), enfouissement ou incinération contrôlée (ultime).
\end{enumerate}
\textbf{Réflexe Probatoire} : dans un tableau de tri, justifie toujours chaque classement par un critère (biodégradabilité, toxicité, recyclabilité) et non par une simple intuition.
\end{methode}

\begin{experience}[Tri des déchets de la cantine -- TP réalisable en classe]
\textbf{Matériel} : échantillons de déchets (épluchures, bouteilles plastiques, canettes, papier, piles usagées, emballages souillés), 4 bacs étiquetés (Organique, Recyclable, Dangereux, Ultime), gants, balances.
\textbf{Protocole} :
\begin{enumerate}
\item Peser la masse totale du mélange de déchets.
\item Répartir chaque déchet dans le bac correspondant en appliquant la méthode ci-dessus.
\item Peser la masse de chaque catégorie.
\item Calculer le pourcentage de chaque catégorie.
\end{enumerate}
\textbf{Observations attendues} : prédominance de la fraction organique (souvent $>50\,\%$ en milieu scolaire/urbain camerounais), fraction recyclable significative, fraction dangereuse faible en masse mais à risque élevé.
\textbf{Interprétation} : le tri à la source permet d'orienter chaque flux vers sa filière de valorisation adaptée.
\end{experience}

\begin{exoresolu}[Tri des déchets d'un marché de Yaoundé]
Après une journée au marché Mokolo, on ramasse les déchets suivants : épluchures de manioc, sachets plastiques, piles usagées, boîtes de conserve vides, restes de poisson, médicaments périmés, cartons d'emballage, bouteilles en verre.
\begin{enumerate}
\item Classe ces déchets en trois catégories : organiques, recyclables non dangereux, dangereux.
\item Pour chaque catégorie, propose une filière de traitement adaptée.
\end{enumerate}
\tcblower
\textbf{1. Classement des déchets.}

On applique la démarche : nature de la matière $\rightarrow$ biodégradabilité $\rightarrow$ danger $\rightarrow$ recyclabilité.

\begin{center}
\begin{tabularx}{\linewidth}{|l|X|X|}
\hline
\rowcolor{popblueL} \textbf{Déchet} & \textbf{Catégorie} & \textbf{Justification} \\
\hline
Épluchures de manioc & Organique & Matière végétale, biodégradable par les micro-organismes \\
\hline
Restes de poisson & Organique & Matière animale, putrescible et biodégradable \\
\hline
Sachets plastiques & Recyclable non dangereux & Polymère non biodégradable mais recyclable (PE/PP) \\
\hline
Boîtes de conserve vides & Recyclable non dangereux & Métal (fer/aluminium) recyclable à l'infini \\
\hline
Cartons d'emballage & Recyclable non dangereux & Cellulose recyclable (ou compostable si souillé) \\
\hline
Bouteilles en verre & Recyclable non dangereux & Verre recyclable à l'infini sans perte de qualité \\
\hline
Piles usagées & Dangereux & Contiennent des métaux lourds toxiques (mercure, cadmium, plomb) \\
\hline
Médicaments périmés & Dangereux & Substances chimiques actives toxiques pour les sols et les eaux \\
\hline
\end{tabularx}
\end{center}

\textbf{2. Filières de traitement adaptées.}
\begin{itemize}
\item \textbf{Déchets organiques} (épluchures, restes de poisson) : \cle{compostage} pour produire un engrais naturel destiné aux champs et jardins, ou méthanisation pour produire du \cle{biogaz}.
\item \textbf{Déchets recyclables} (plastiques, métaux, cartons, verre) : orientation vers les \cle{filières de recyclage} (les plastiques peuvent être transformés en pavés ou objets, le verre refondu, les métaux refondus).
\item \textbf{Déchets dangereux} (piles, médicaments) : \cle{collecte spécialisée} séparée, puis traitement ou stockage sécurisé ; ils ne doivent jamais être jetés avec les ordures ménagères ni brûlés à l'air libre.
\end{itemize}
\textbf{Conclusion} : un tri à la source permet de valoriser la majorité des déchets du marché et d'isoler les déchets dangereux qui menacent la santé publique et l'environnement.
\end{exoresolu}

\begin{aretenir}[Tri des déchets]
\begin{itemize}
\item Le tri se fonde sur trois critères scientifiques : \textbf{biodégradabilité}, \textbf{dangerosité}, \textbf{recyclabilité}.
\item Trois filières principales : \textbf{organique} $\rightarrow$ compostage/méthanisation ; \textbf{recyclable} $\rightarrow$ recyclage matière ; \textbf{dangereux} $\rightarrow$ filière spécialisée.
\item Tout déchet non trié devient un déchet ultime (enfouissement/incinération), perte de ressources et source de pollution.
\end{itemize}
\end{aretenir}

\coursec{Impacts des déchets sur l'environnement et la santé}

\begin{definition}[Lixiviat (jus de décharge)]
Liquide pollué résultant de la percolation de l'eau de pluie à travers une masse de déchets, chargé en matières organiques, sels minéraux, métaux lourds et micro-organismes pathogènes. Il contamine sols et nappes phréatiques s'il n'est pas collecté et traité.
\end{definition}

\begin{propriete}[Chaîne d'impacts « Déchet $\rightarrow$ Mécanisme $\rightarrow$ Conséquence »]
\begin{center}
\begin{tabularx}{\linewidth}{|l|X|X|}
\hline
\rowcolor{popblueL} \textbf{Déchet / Propriété} & \textbf{Mécanisme de nuisance} & \textbf{Conséquence sanitaire / environnementale} \\
\hline
Plastiques (non biodégradables) & Obstruction des caniveaux et cours d'eau & Inondations urbaines, asphyxie des milieux aquatiques \\
\hline
Déchets divers (récipients) & Rétention d'eau stagnante & Gîtes larvaires \emph{Anopheles} $\rightarrow$ \textbf{Paludisme} ; \emph{Aedes} $\rightarrow$ Dengue, Chikungunya \\
\hline
Déchets organiques (putrescibles) & Attraction mouches/rongeurs + production de lixiviat & Contamination eau/aliments $\rightarrow$ \textbf{Diarrées, Choléra, Typhoïde} \\
\hline
Piles, batteries, produits chimiques & Lessivage métaux lourds (Hg, Cd, Pb) + substances toxiques & Pollution durable sols/nappes $\rightarrow$ \textbf{Intoxications chroniques} (rénales, neurologiques), stérilisation sols \\
\hline
Brûlage à l'air libre (plastiques, pneus) & Émission dioxines, furanes, particules fines (PM2,5), CO, CO$_2$ & \textbf{Maladies respiratoires, cancers}, réchauffement climatique, retombées acides \\
\hline
\end{tabularx}
\end{center}
\end{propriete}

\courssub{Savoir-faire 5 : Analyser les impacts des déchets sur l'environnement et la santé}

\begin{methode}[Construire une chaîne « déchet $\rightarrow$ impact »]
Pour analyser rigoureusement l'impact d'un déchet, raisonne toujours en trois maillons :
\begin{enumerate}
\item \textbf{Le déchet et sa propriété} : non biodégradabilité (plastique : plusieurs siècles), toxicité (métaux lourds des piles), putrescibilité (déchets organiques), pouvoir calorifique.
\item \textbf{Le mécanisme de nuisance} : contamination des sols et des nappes phréatiques (lixiviat), obstruction des caniveaux $\rightarrow$ inondations, eaux stagnantes $\rightarrow$ prolifération des moustiques, fumées toxiques lors de brûlages sauvages (dioxines, furanes), ingestion par les animaux (faune sauvage, bétail).
\item \textbf{La conséquence sanitaire ou environnementale} : paludisme, choléra, diarrées, intoxications aiguës ou chroniques, mortalité de la faune, stérilisation des sols, émissions de gaz à effet de serre (méthane \ce{CH4} des décharges anaérobies, \ce{CO2} de l'incinération).
\end{enumerate}
\textbf{Réflexe} : une bonne réponse au Probatoire relie toujours explicitement le déchet à la maladie ou à la dégradation par un mécanisme précis, et non par une affirmation vague (« ça pollue »).
\end{methode}

\begin{exoresolu}[Décharge sauvage et santé dans un quartier de Yaoundé]
Dans un quartier de Yaoundé, les habitants déposent leurs ordures dans un ravin qui draine les eaux de pluie vers un puits. On y trouve des déchets organiques, des sachets plastiques, des piles et des bidons de produits d'entretien. Pendant la saison des pluies, le quartier connaît des inondations et une recrudescence de cas de paludisme et de diarrées.
\begin{enumerate}
\item Explique, par un mécanisme précis, le lien entre la décharge et chacun des problèmes observés : inondations, paludisme, diarrées.
\item Explique pourquoi la présence de piles et de produits d'entretien aggrave la situation à long terme.
\item Propose deux mesures pour résoudre durablement le problème.
\end{enumerate}
\tcblower
\textbf{1. Mécanismes de nuisance.}
\begin{itemize}
\item \textbf{Inondations} : les sachets plastiques, non biodégradables, \emph{obstruent le ravin et les caniveaux} ; l'eau de pluie ne s'écoule plus et déborde dans les habitations.
\item \textbf{Paludisme} : les déchets (sachets, boîtes, bidons) \emph{retiennent l'eau de pluie en petites collections stagnantes}, qui deviennent des gîtes larvaires où les femelles d'\emph{Anopheles} pondent ; la multiplication des moustiques vecteurs augmente la transmission du paludisme.
\item \textbf{Diarrées} : les déchets organiques en décomposition attirent mouches et rongeurs et libèrent un jus de décharge (\cle{lixiviat}) qui s'infiltre et \emph{contamine le puits} par des microbes pathogènes (bactéries comme \emph{Vibrio cholerae}, \emph{Salmonella}, virus, parasites) ; la consommation de cette eau provoque des diarrées, voire le choléra.
\end{itemize}
\textbf{2. Danger à long terme des déchets dangereux.}

Les piles contiennent des \cle{métaux lourds} (mercure, cadmium, plomb) et les produits d'entretien des substances chimiques toxiques (solvants, tensioactifs, chlore). Lessivés par les pluies, ces polluants \emph{s'accumulent durablement dans les sols et la nappe phréatique} car ils ne se dégradent pas (non biodégradables, non métabolisables) ; ils provoquent des intoxications chroniques (atteintes rénales, neurologiques, cancérigènes) et stérilisent les sols pendant des décennies (phytotoxicité).

\textbf{3. Mesures durables.}
\begin{itemize}
\item Organiser le \textbf{tri à la source} et une \textbf{collecte régulière} des ordures, avec compostage des déchets organiques et recyclage des plastiques ;
\item Assainir le ravin, \textbf{interdire et surveiller les dépôts sauvages}, protéger le puits (margelle, couvercle, éloignement des sources de pollution $>30$ m) et \textbf{sensibiliser} les habitants à la gestion responsable des déchets.
\end{itemize}
\textbf{Conclusion} : la décharge sauvage dégrade à la fois l'environnement (inondations, pollution durable des sols et des eaux) et la santé (paludisme, diarrées, intoxications chroniques) ; seuls le tri, la collecte organisée et la sensibilisation rompent durablement cette chaîne de nuisances.
\end{exoresolu}

\begin{aretenir}[Impacts majeurs des déchets non gérés]
\begin{itemize}
\item \textbf{Environnement} : pollution sols/eaux (lixiviat, métaux lourds), gaz à effet de serre (\ce{CH4}, \ce{CO2}, \ce{N2O}), perte de biodiversité, dégradation paysagère.
\item \textbf{Santé} : maladies vectorielles (paludisme, arboviroses), maladies hydriques (choléra, typhoïde), toxiques (métaux lourds, dioxines), blessures/infections (objets tranchants, DASRI).
\item \textbf{Économie} : coûts de santé, perte de productivité agricole, coûts de dépollution, impact touristique.
\end{itemize}
\end{aretenir}

\coursec{Le tri des déchets et la règle des 3R}

\begin{propriete}[Hiérarchie des 3R -- Ordre de priorité environnementale]
\[
\textbf{Réduire} \;>\; \textbf{Réutiliser} \;>\; \textbf{Recycler}
\]
\begin{itemize}
\item \textbf{Réduire} (priorité absolue) : éviter la production du déchet $\rightarrow$ économie de matières premières, d'énergie, d'eau, pas de transport ni traitement.
\item \textbf{Réutiliser} : prolonger la durée de vie de l'objet tel quel $\rightarrow$ faible coût énergétique (nettoyage, réparation).
\item \textbf{Recycler} : transformer la matière $\rightarrow$ coût énergétique et industriel significatif, pertes de matière, mais évite l'extraction de ressources vierges.
\end{itemize}
\end{propriete}

\courssub{Savoir-faire 2 : Appliquer la règle des 3R}

\begin{methode}[Analyser une situation à l'aide de la règle des 3R]
La gestion durable des déchets repose sur trois principes hiérarchisés :
\begin{enumerate}
\item \textbf{\cle{Réduire}} : diminuer la quantité de déchets produits à la source (achat en vrac, refus des emballages superflus, limitation du gaspillage alimentaire, éco-conception). C'est le principe \emph{prioritaire}.
\item \textbf{\cle{Réutiliser}} : employer à nouveau un objet pour le même usage ou un autre usage sans transformation industrielle (bouteilles consignées, sacs réutilisables, réparation, détournement d'usage : pneu $\rightarrow$ jardinière).
\item \textbf{\cle{Recycler}} : transformer la matière du déchet pour fabriquer un nouveau produit (fonte du verre, régénération du plastique, compostage des biodéchets -- recyclage organique).
\end{enumerate}
\textbf{Réflexe} : pour classer une action, demande-toi : « la matière est-elle transformée industriellement ? ». Si oui $\rightarrow$ recyclage ; si l'objet est réemployé tel quel $\rightarrow$ réutilisation ; si on évite de produire le déchet $\rightarrow$ réduction. L'ordre de priorité est toujours : Réduire $>$ Réutiliser $>$ Recycler.
\end{methode}

\begin{exoresolu}[Les 3R dans une famille de Douala]
Une famille de Douala adopte les comportements suivants :
\begin{itemize}
\item[(a)] Elle achète le riz en vrac avec son propre contenant au lieu de sachets plastiques individuels.
\item[(b)] Elle transforme les vieux pneus en fauteuils et en jardinières.
\item[(c)] Elle apporte ses bouteilles en verre vides chez le revendeur qui les consigne.
\item[(d)] Elle envoie ses déchets plastiques à une entreprise qui les fond pour fabriquer des pavés de cour.
\end{itemize}
\begin{enumerate}
\item Attribue à chaque action le principe des 3R qu'elle illustre, en justifiant.
\item Classe ces actions par ordre de priorité environnementale.
\end{enumerate}
\tcblower
\textbf{1. Attribution des principes.}
\begin{itemize}
\item[(a)] \textbf{Réduire} : l'achat en vrac évite la production des sachets plastiques ; le déchet n'est \emph{jamais créé}. C'est une réduction à la source.
\item[(b)] \textbf{Réutiliser} : le pneu est réemployé tel quel (comme fauteuil ou jardinière) sans transformation industrielle de sa matière ; on prolonge sa durée de vie.
\item[(c)] \textbf{Réutiliser} : la consigne permet de remplir à nouveau la même bouteille en verre, sans refonte ; c'est le réemploi classique (circuit court).
\item[(d)] \textbf{Recycler} : le plastique est fondu, c'est-à-dire que sa \emph{matière} est transformée pour fabriquer un produit nouveau (les pavés).
\end{itemize}
\textbf{2. Ordre de priorité environnementale.}

La hiérarchie des 3R est : Réduire $>$ Réutiliser $>$ Recycler, car réduire évite toute consommation de matière et d'énergie, réutiliser ne demande qu'un nettoyage/réparation, tandis que recycler exige une transformation industrielle coûteuse en énergie et en eau. On obtient donc :
\[ (a) \;>\; (b) \text{ et } (c) \;>\; (d) \]
\textbf{Conclusion} : la réduction à la source (a) est l'action la plus vertueuse ; le recyclage (d), bien qu'utile, n'intervient qu'en dernier recours avant l'élimination.
\end{exoresolu}

\begin{exercice}[Application des 3R à la maison]
Pour chacune des actions suivantes, indique le principe des 3R illustré et justifie en une phrase.
\begin{enumerate}
\item J'utilise une gourde en inox au lieu de bouteilles d'eau jetables.
\item Je donne mes vieux vêtements à une association caritative.
\item Je répare mon téléphone portable au lieu d'en racheter un neuf.
\item Je mets mes bouteilles en plastique dans le bac jaune de tri sélectif.
\item J'achète des légumes sans emballage au marché.
\item Je transforme des bocaux en verre en pots de rangement pour la cuisine.
\end{enumerate}
\end{exercice}

\begin{corrige}[Exercice Application des 3R]
\begin{enumerate}
\item \textbf{Réutiliser} : la gourde remplace des centaines de bouteilles jetables ; l'objet est réemployé tel quel.
\item \textbf{Réutiliser} : les vêtements sont réemployés par d'autres personnes sans transformation.
\item \textbf{Réutiliser} (voire Réduire) : la réparation prolonge la durée de vie, évitant la production d'un nouvel appareil (réduction indirecte).
\item \textbf{Recycler} : le plastique est collecté, broyé, lavé, fondu pour faire de nouveaux objets (transformation de la matière).
\item \textbf{Réduire} : l'absence d'emballage évite la création du déchet à la source.
\item \textbf{Réutiliser} : le bocal est réemployé tel quel pour un nouvel usage (détournement).
\end{enumerate}
\end{corrige}

\coursec{Transformation et recyclage : principes et filières}

\begin{definition}[Recyclage]
Opération de valorisation des déchets par laquelle des matières sont retraitées en substances, matières ou produits pour leur fonction d'origine ou à d'autres fins. Il exclut la valorisation énergétique (incinération avec récupération de chaleur) et le remblayage.
\end{definition}

\begin{propriete}[Principales filières de recyclage matière au Cameroun]
\begin{center}
\begin{tabularx}{\linewidth}{|l|X|X|}
\hline
\rowcolor{popblueL} \textbf{Matière} & \textbf{Procédé principal} & \textbf{Produits recyclés (exemples locaux)} \\
\hline
Verre & Tri par couleur $\rightarrow$ broyage (calcin) $\rightarrow$ fusion ($1500\,^\circ\mathrm{C}$) & Nouvelles bouteilles, laine de verre, isolants \\
\hline
Métaux (Fer, Al) & Tri magnétique/courants de Foucault $\rightarrow$ broyage $\rightarrow$ fusion & Fers à béton (aciéries locales), ustensiles, pièces auto \\
\hline
Papier/Carton & Désencrage $\rightarrow$ pâte à papier $\rightarrow$ séchage & Carton ondulé, papier journal, mouchoirs, boîtes d'œufs \\
\hline
Plastiques (PET, PEHD, PP) & Tri $\rightarrow$ broyage $\rightarrow$ lavage $\rightarrow$ extrusion/granulés & Pavés écologiques, tuyaux, arrosoirs, sacs, fils polyester (PET) \\
\hline
\end{tabularx}
\end{center}
\end{propriete}

\begin{savaistu}[Initiatives locales de recyclage]
\begin{itemize}
\item \textbf{Kemit Ecology (Douala)} : transforme les déchets plastiques en pavés écologiques et en charbon écologique (briquettes de déchets organiques carbonisés).
\item \textbf{Plastic Waste Recycling (Yaoundé)} : collecte et recycle le PET en paillettes pour l'export ou la filière textile.
\item \textbf{Groupements de femmes (région Ouest, Nord)} : tressage de sacs en sachets plastiques récupérés (artisanat d'art).
\item \textbf{Projet "Zéro Déchet" dans certains lycées} : compostage des déchets de cantine, recyclage papier/plastique avec partenaires privés.
\end{itemize}
\end{savaistu}

\courssub{Savoir-faire 3 : Proposer un mode de transformation ou de recyclage adapté}

\begin{methode}[Choisir la filière de valorisation adaptée à un déchet]
\begin{enumerate}
\item \textbf{Caractériser le déchet} : nature (organique, plastique, métal, verre, dangereux), quantité disponible, régularité de la production, humidité, propreté.
\item \textbf{Identifier les propriétés exploitables} : fermentescibilité (déchets organiques), fusibilité (verre, métal, plastique), pouvoir calorifique, présence de nutriments (azote, phosphore, potassium pour le compost).
\item \textbf{Associer le déchet à une technologie adaptée au contexte local} :
\begin{itemize}
\item déchets organiques humides $\rightarrow$ \cle{compostage} (aérobie) ou \cle{méthanisation} (biogaz, anaérobie) ;
\item plastiques $\rightarrow$ tri par type (code résine), broyage, lavage, fonte (pavés, tuyaux, objets) ;
\item métaux et verre $\rightarrow$ refonte en fonderie/verrerie ;
\item déchets dangereux $\rightarrow$ collecte et traitement spécialisés (incinération haute température $>1200\,^\circ\mathrm{C}$ pour DASRI, stabilisation/solidification pour métaux lourds).
\end{itemize}
\item \textbf{Évaluer les retombées} : produit obtenu (compost, biogaz, matière première secondaire), emplois créés, réduction de la pollution, économie circulaire.
\item \textbf{Justifier le choix} en comparant avec les alternatives (enfouissement, incinération sauvage) et leurs nuisances.
\end{enumerate}
\textbf{Réflexe} : ne propose jamais une filière sans préciser le \emph{produit valorisé} obtenu et la \emph{condition biologique ou technique} du procédé (aérobie/anaérobie, température, durée, tri amont).
\end{methode}

\begin{exoresolu}[Valoriser les déchets d'un lycée de Bafoussam]
Le lycée de Bafoussam produit chaque semaine environ $200$ kg de déchets : $60\,\%$ de déchets de cuisine et d'herbes de tonte (organiques), $25\,\%$ de plastiques, $10\,\%$ de papiers et $5\,\%$ de divers non recyclables. Le lycée dispose d'un grand jardin pédagogique et d'une cantine.
\begin{enumerate}
\item Calcule la masse hebdomadaire de déchets organiques produite par le lycée.
\item Propose deux modes de valorisation de ces déchets organiques adaptés au contexte, en précisant le produit obtenu et son utilisation au lycée.
\item Indique le devenir à réserver aux plastiques et aux papiers.
\end{enumerate}
\tcblower
\textbf{1. Masse de déchets organiques.}
\[ m_{\text{organique}} = 200 \times \dfrac{60}{100} = 120 \text{ kg par semaine.} \]
Le lycée produit donc $120$ kg de déchets organiques valorisables chaque semaine.

\textbf{2. Deux modes de valorisation des déchets organiques.}
\begin{itemize}
\item \textbf{Le compostage} : les déchets de cuisine et l'herbe de tonte sont empilés en andains ou placés en composteur, en présence d'\cle{air} (décomposition aérobie), d'humidité ($50\text{--}60\,\%$) et de micro-organismes (bactéries, champignons, actinomycètes). Après quelques semaines à quelques mois (maturation), on obtient du \cle{compost}, un amendement organique riche en humus stable et en éléments nutritifs (azote, phosphore, potassium). \emph{Utilisation au lycée} : fertiliser le jardin pédagogique et les espaces verts, ce qui supprime l'achat d'engrais chimiques et améliore la structure du sol.
\item \textbf{La méthanisation (production de biogaz)} : les déchets de cuisine (plus riches en énergie que l'herbe) sont introduits dans un \cle{digesteur} étanche (béton, géomembrane), où ils fermentent \emph{en l'absence d'air} (anaérobie) grâce à des bactéries hydrolytiques, acétogènes et méthanogènes (méso $35\,^\circ\mathrm{C}$ ou thermophile $55\,^\circ\mathrm{C}$). On obtient du \cle{biogaz} (mélange $\approx 60\text{--}70\,\%$ méthane \ce{CH4}, $30\text{--}40\,\%$ \ce{CO2}, traces \ce{H2S}) et un \cle{digestat} (liquide/solide) utilisable comme engrais. \emph{Utilisation au lycée} : le biogaz alimente les foyers de la cantine (cuisson), réduisant l'achat de bois ou de gaz butane ; le digestat fertilise le jardin.
\end{itemize}
\textbf{3. Devenir des plastiques et des papiers.}
\begin{itemize}
\item Les \textbf{plastiques} ($200 \times 0{,}25 = 50$ kg/semaine) sont triés par type (bouteilles PET, sachets PE), lavés, puis vendus ou apportés à une filière de \cle{recyclage} (broyage et fonte pour fabriquer des pavés, arrosoirs, tuyaux).
\item Les \textbf{papiers} ($200 \times 0{,}10 = 20$ kg/semaine) sont triés et orientés vers le recyclage papetier ; les papiers souillés (serviettes, emballages gras) peuvent rejoindre le compost (apport de carbone).
\end{itemize}
\textbf{Conclusion} : en combinant compostage, méthanisation et recyclage, le lycée valorise $95\,\%$ de ses déchets, produit de l'engrais et de l'énergie, et ne rejette plus que $10$ kg de déchets ultimes par semaine.
\end{exoresolu}

\begin{experience}[Fabrication d'un composteur scolaire -- Projet pédagogique]
\textbf{Objectif} : produire du compost à partir des déchets de la cantine et des tontes.
\textbf{Matériel} : palettes en bois ou grillage (1 m$^3$), bâche, fourche, thermomètre à compost, seau, eau.
\textbf{Protocole} :
\begin{enumerate}
\item Construire un bac à compost (3 compartiments : frais, en cours, mûr).
\item Alterner couches \textbf{riches en azote} (déchets cuisine, herbe fraîche) et couches \textbf{riches en carbone} (feuilles mortes, paille, carton déchiqueté, sciure) -- ratio C/N $\approx 25\text{--}30$.
\item Maintenir l'humidité (essorage au poing : quelques gouttes) et l'aération (retourner toutes les 2--3 semaines).
\item Surveiller la température : montée à $55\text{--}65\,^\circ\mathrm{C}$ (phase thermophile, hygiénisation), puis descente.
\item Après $3\text{--}6$ mois : compost brun, odorant humus, température ambiante $\rightarrow$ tamisage et utilisation.
\end{enumerate}
\textbf{Observations} : réduction du volume ($\approx 50\,\%$), disparition des déchets reconnaissables, apparition de vers de terre (signe de maturité).
\end{experience}

\coursec{Valorisation des déchets organiques : compostage et biogaz}

\begin{definition}[Compostage]
Procédé biologique aérobie de décomposition et de transformation de la matière organique en un produit stable, hygiénique, riche en humus : le compost. Il met en jeu des successions de populations microbiennes (méso $\rightarrow$ thermophile $\rightarrow$ méso) et de la macrofaune (vers, collemboles).
\end{definition}

\begin{definition}[Méthanisation (Digestion anaérobie)]
Procédé biologique anaérobie de dégradation de la matière organique par un consortium bactérien (hydrolyse $\rightarrow$ acido-genèse $\rightarrow$ acéto-genèse $\rightarrow$ méthano-genèse) produisant du biogaz (\ce{CH4} + \ce{CO2}) et un digestat. Nécessite un réacteur étanche (digesteur), température contrôlée, pH neutre, absence d'inhibiteurs (NH$_3$, H$_2$S, métaux lourds).
\end{definition}

\begin{propriete}[Comparaison Compostage vs Méthanisation]
\begin{center}
\begin{tabularx}{\linewidth}{|l|X|X|}
\hline
\rowcolor{popblueL} \textbf{Critère} & \textbf{Compostage (Aérobie)} & \textbf{Méthanisation (Anaérobie)} \\
\hline
Condition en O$_2$ & Présence d'air (aération naturelle ou forcée) & Absence d'air (digesteur étanche) \\
\hline
Micro-organismes & Bactéries aérobies, champignons, actinomycètes, macrofaune & Bactéries anaérobies strictes (méthanogènes archées) \\
\hline
Produit principal & Compost (amendement organique solide) & Biogaz (gaz combustible) + Digestat (engrais liquide/solide) \\
\hline
Énergie & Perte de chaleur (exothermique, non récupérée) & Récupération énergie chimique (méthane) \\
\hline
Durée & $2$ à $6$ mois (selon suivi) & $20$ à $60$ jours (temps de séjour hydraulique) \\
\hline
Investissement & Faible (composteur, surface) & Élevé (digesteur, gazomètre, désulfurisation, cogénération) \\
\hline
Adaptation locale & Très adaptée (rural, périurbain, établissement) & Adaptée si gisement régulier $\&$ besoin énergie (cantine, ferme, abattoir) \\
\hline
\end{tabularx}
\end{center}
\end{propriete}

\begin{popfigure}
\begin{tikzpicture}[scale=0.9, every node/.style={font=\small}]
% Compostage
\begin{scope}[xshift=0cm]
\node[draw, rounded corners, fill=popgreenL, align=center] (c1) at (0,3) {Déchets organiques\\ (Cuisine, jardin)};
\node[draw, rounded corners, fill=poporangeL, align=center] (c2) at (0,1.5) {Composteur / Andain\\ \textbf{Aérobie} : Air, Eau, Micro-organismes};
\node[draw, rounded corners, fill=popblueL, align=center] (c3) at (0,0) {CO$_2$ + Chaleur + H$_2$O};
\node[draw, rounded corners, fill=popgoldL, align=center] (c4) at (0,-1.5) {COMPOST\\ (Humus, NPK)};
\draw[->, thick] (c1) -- (c2);
\draw[->, thick] (c2) -- (c3);
\draw[->, thick] (c2) -- (c4);
\node[align=center] at (0,4) {\textbf{COMPOSTAGE}};
\end{scope}

% Méthanisation
\begin{scope}[xshift=7cm]
\node[draw, rounded corners, fill=popgreenL, align=center] (m1) at (0,3) {Déchets organiques\\ (Lisier, cuisine, agro-industrie)};
\node[draw, rounded corners, fill=poppurpleL, align=center] (m2) at (0,1.5) {DIGESTEUR\\ \textbf{Anaérobie} : Étanche, 35--55°C, pH 7--8};
\node[draw, rounded corners, fill=popblueL, align=center] (m3) at (0,0) {BIOGAZ\\ (CH$_4$ 60--70\%, CO$_2$, H$_2$S)};
\node[draw, rounded corners, fill=popgoldL, align=center] (m4) at (0,-1.5) {DIGESTAT\\ (Engrais liquide)};
\draw[->, thick] (m1) -- (m2);
\draw[->, thick] (m2) -- (m3);
\draw[->, thick] (m2) -- (m4);
\node[align=center] at (0,4) {\textbf{MÉTHANISATION}};
\end{scope}

% Usages
\node[draw, rounded corners, fill=poppinkL, align=center] at (3.5,-3) {USAGES : Cuisson, Éclairage, Électricité, Chauffage, Carburant (bioGNV)};
\draw[->, thick, poppink] (3.5,-1.5) -- (3.5,-2.5);
\end{tikzpicture}
\legende{Schéma comparatif des deux voies de valorisation des déchets organiques : compostage (voie aérobie) et méthanisation (voie anaérobie).}
\end{popfigure}

\courssub{Savoir-faire 4 : Exploiter quantitativement le compostage et le biogaz}

\begin{methode}[Traiter les calculs liés à la valorisation des déchets]
\begin{enumerate}
\item \textbf{Repérer les données} : masse totale, pourcentages de chaque catégorie, rendement de transformation (fraction de la masse convertie en produit utile).
\item \textbf{Appliquer les relations de base} :
\[ m_{\text{catégorie}} = m_{\text{totale}} \times \dfrac{\text{pourcentage}}{100} \qquad ; \qquad m_{\text{produit}} = m_{\text{matière}} \times \text{rendement}. \]
\item \textbf{Pour le compostage} : la décomposition aérobie s'accompagne d'une perte de masse (eau évaporée, \ce{CO2} dégagé) ; le compost représente typiquement $30$ à $50\,\%$ de la masse initiale de biodéchets mis en composteur.
\item \textbf{Pour le biogaz} : retiens l'équation simplifiée de la méthanisation de la matière organique (glucose comme modèle) :
\[ \ce{C6H12O6 -> 3CH4 + 3CO2} \]
Le biogaz contient environ $50$ à $70\,\%$ de méthane \ce{CH4}, combustible. Le rendement volumique typique est de $0{,}05$ à $0{,}15$ m$^3$ biogaz par kg de matière volatile (selon substrat).
\item \textbf{Vérifier les unités et la cohérence} : une masse de produit ne peut jamais dépasser la masse de matière initiale ; convertis les pourcentages en fractions avant de calculer.
\end{enumerate}
\textbf{Réflexe} : au Probatoire, écris toujours la formule littérale, puis l'application numérique, puis le résultat avec son unité.
\end{methode}

\begin{exoresolu}[Compost et biogaz d'une coopérative de maraîchers]
Une coopérative de maraîchers de la région de l'Ouest collecte $800$ kg de déchets organiques par mois (fanures, fumier, résidus de récolte).
\begin{enumerate}
\item La coopérative composte $75\,\%$ de ces déchets. Sachant que le rendement du compostage est de $40\,\%$ (masse de compost obtenue par rapport à la masse de déchets mis à composter), calcule la masse mensuelle de compost produite.
\item Les $25\,\%$ restants sont méthanisés dans un digesteur qui produit $0{,}08$ m$^3$ de biogaz par kilogramme de déchets. Calcule le volume mensuel de biogaz produit.
\item Le biogaz contient $60\,\%$ de méthane. Calcule le volume de méthane disponible chaque mois, puis cite un usage de ce biogaz pour la coopérative.
\end{enumerate}
\tcblower
\textbf{1. Masse mensuelle de compost.}

Masse de déchets compostés :
\[ m_c = 800 \times \dfrac{75}{100} = 600 \text{ kg.} \]
Masse de compost obtenue (rendement $40\,\%$) :
\[ m_{\text{compost}} = 600 \times \dfrac{40}{100} = 240 \text{ kg.} \]
La coopérative produit \textbf{240 kg de compost par mois}, utilisé comme engrais naturel sur les parcelles maraîchères.

\textbf{2. Volume mensuel de biogaz.}

Masse de déchets méthanisés :
\[ m_m = 800 \times \dfrac{25}{100} = 200 \text{ kg.} \]
Volume de biogaz :
\[ V_{\text{biogaz}} = 200 \times 0{,}08 = 16 \text{ m}^3. \]
La coopérative produit \textbf{16 m$^3$ de biogaz par mois}.

\textbf{3. Volume de méthane et usage.}

Le méthane représente $60\,\%$ du biogaz :
\[ V_{\ce{CH4}} = 16 \times \dfrac{60}{100} = 9{,}6 \text{ m}^3. \]
La coopérative dispose de \textbf{9,6 m$^3$ de méthane par mois}. Ce biogaz peut servir à la \emph{cuisson} des aliments dans les foyers améliorés, à l'\emph{éclairage} (lampes à biogaz) ou au chauffage, réduisant ainsi la consommation de bois de chauffe et la déforestation.

\textbf{Conclusion} : la valorisation combinée (compostage + méthanisation) transforme la totalité des $800$ kg de déchets en $240$ kg d'engrais naturel et $16$ m$^3$ de gaz combustible : aucun déchet organique n'est perdu.
\end{exoresolu}

\begin{exoresolu}[Calcul de rendement compost -- Cas d'un lycée]
Un lycée met en place un composteur. Sur un mois, $500$ kg de déchets de cantine et $300$ kg de tontes sont compostés. On récupère $280$ kg de compost tamisé.
\begin{enumerate}
\item Calcule le rendement massique global du compostage.
\item Si le lycée a besoin de $1$ tonne de compost par an pour son jardin, quelle masse de déchets organiques doit-il composter par an (même rendement) ?
\end{enumerate}
\tcblower
\textbf{1. Rendement massique.}
Masse initiale $m_i = 500 + 300 = 800$ kg.
Masse finale $m_f = 280$ kg.
Rendement $R = \dfrac{m_f}{m_i} \times 100 = \dfrac{280}{800} \times 100 = 35\,\%$.
\textbf{2. Masse de déchets nécessaire pour 1 tonne de compost.}
$m_{\text{déchets}} = \dfrac{m_{\text{compost visé}}}{R} = \dfrac{1000}{0{,}35} \approx 2857$ kg $\approx 2{,}86$ tonnes par an.
\end{exoresolu}

\begin{aretenir}[Valorisation organique : points clés chiffrés]
\begin{itemize}
\item Compostage : rendement massique $30\text{--}50\,\%$ ; produit = amendement organique (humus, NPK) ; condition = \textbf{aérobie} (air, eau, C/N $\approx 30$).
\item Méthanisation : rendement volumique $0{,}05\text{--}0{,}15$ m$^3$ biogaz/kg déchets ; biogaz = $50\text{--}70\,\%$ \ce{CH4} ; condition = \textbf{anaérobie} (étanche, $35\text{--}55\,^\circ\mathrm{C}$, pH $7\text{--}8$).
\item Digestat = engrais liquide riche en N minéral (NH$_4^+$), utilisable directement ou après séparation de phase.
\item Équation type : $\ce{C6H12O6 -> 3CH4 + 3CO2}$ (bilan global simplifié).
\end{itemize}
\end{aretenir}

\coursec{Agir pour une gestion responsable des déchets}

\courssub{Savoir-faire 6 : Sensibiliser à la gestion responsable des déchets}

\begin{methode}[Concevoir une action de sensibilisation efficace]
\begin{enumerate}
\item \textbf{Définir la cible} (élèves, commerçants du marché, habitants d'un quartier, personnel de cantine) et \textbf{le problème précis} à résoudre (brûlage des ordures, absence de tri, dépôts sauvages, gaspillage alimentaire).
\item \textbf{Formuler un message simple, positif et actionnable} : une action concrète à adopter (« trie tes déchets organiques pour le compost ») plutôt qu'une interdiction vague (« ne jette pas n'importe où »).
\item \textbf{Choisir les supports adaptés} au contexte : affiches illustrées (pictogrammes), sketches/théâtre-forum, radio communautaire (émissions en langues locales), démonstration pratique (fabrication d'un composteur, tri en direct), journée de nettoyage citoyen (\emph{Clean-up day}), réseaux sociaux (WhatsApp, Facebook pour les jeunes).
\item \textbf{S'appuyer sur des arguments locaux et chiffrés} : coût des maladies évitées (paludisme, diarrées), compost gratuit pour les champs (économie engrais), revenus du recyclage (vente plastiques/métaux), embellissement du quartier, fierté collective.
\item \textbf{Prévoir l'évaluation} : observation du changement de comportement (propreté du site, nombre de composteurs installés, taux de tri), enquêtes de satisfaction, relance de l'action à intervalles réguliers.
\end{enumerate}
\textbf{Réflexe} : une bonne sensibilisation ne se contente pas d'informer : elle \emph{montre un geste}, donne un \emph{bénéfice concret et immédiat} et \emph{vérifie} que le comportement a changé durablement.
\end{methode}

\begin{exoresolu}[Campagne de sensibilisation au lycée]
Le club environnement d'un lycée de Buea veut lancer une campagne intitulée « Zéro plastique brûlé, tout organique composté ». Les élèves jettent leurs déchets en vrac dans une poubelle unique et le concierge brûle chaque soir le contenu à l'air libre.
\begin{enumerate}
\item Explique pourquoi le brûlage à l'air libre des déchets, notamment des plastiques, est dangereux.
\item Propose un plan de campagne en quatre étapes (cible, message, supports, évaluation).
\item Rédige un slogan complémentaire incitant au compostage des déchets de la cantine.
\end{enumerate}
\tcblower
\textbf{1. Danger du brûlage à l'air libre.}

La combustion des plastiques à basse température ($< 800\,^\circ\mathrm{C}$) et sans contrôle de l'air dégage des \textbf{fumées toxiques} contenant des dioxines (PCDD), des furanes (PCDF), des hydrocarbures aromatiques polycycliques (HAP), du monoxyde de carbone (CO) et des particules fines (PM2,5). Inhalées, ces substances provoquent des maladies respiratoires aiguës et chroniques (asthme, bronchites) et sont \emph{cancérigènes avérés} (CIRC Groupe 1 pour les dioxines) ; elles retombent aussi sur les sols et les cultures voisines, contaminant la chaîne alimentaire (bioaccumulation). Le brûlage détruit en outre des matières qui auraient pu être recyclées ou compostées, gaspillant des ressources.

\textbf{2. Plan de campagne en quatre étapes.}
\begin{enumerate}
\item \textbf{Cible et problème} : les élèves (tous niveaux), le personnel de cantine et le concierge ; le problème est l'absence de tri et le brûlage quotidien des déchets mélangés.
\item \textbf{Message} : « Trie en trois bacs : \textbf{Vert} = Organique (compost), \textbf{Jaune} = Plastique/Métal (recyclage), \textbf{Gris} = Reste (ultime). L'organique nourrit le jardin, le plastique se recycle, \textbf{rien ne se brûle}. »
\item \textbf{Supports} :
\begin{itemize}
\item Installation de trois bacs de couleurs différentes (Vert, Jaune, Gris) à points stratégiques avec affiches explicatives illustrées (pictogrammes).
\item Démonstration pratique de construction et de suivi d'un composteur avec les déchets de la cantine (club environnement + professeurs SVT).
\item Sketch et annonces lors des cérémonies du lundi ; spots audio en anglais/français/pidgin diffusés à la sonnerie.
\item Implication du concierge : formation au nouveau circuit (vidange bacs verts $\rightarrow$ composteur, bacs jaunes $\rightarrow$ collecteur recyclage, bac gris $\rightarrow$ poubelle HYSACAM), valorisation de son rôle.
\end{itemize}
\item \textbf{Évaluation} : pesée hebdomadaire des déchets triés par catégorie (tableau de bord affiché), arrêt effectif du brûlage vérifié chaque semaine par le club, quantité de compost récoltée pour le jardin du lycée après deux mois, enquête de satisfaction élèves/professeurs.
\end{enumerate}
\textbf{3. Slogan pour le compostage.}

« \emph{Les restes de la cantine font les récoltes de demain : compostons !} »

\textbf{Conclusion} : en remplaçant le brûlage par le tri, le compostage et le recyclage, le lycée protège la santé de tous (élèves, personnel, riverains), produit de l'engrais gratuit pour son jardin, génère des revenus potentiels (vente recyclables) et devient un modèle de gestion responsable des déchets pour tout le quartier.
\end{exoresolu}

\begin{exercice}[Concevoir une affiche de sensibilisation]
Tu es membre du club environnement de ton établissement. Concevoir une affiche A4 (format portrait) pour inciter au tri des déchets à la cantine. L'affiche doit contenir :
\begin{enumerate}
\item Un titre accrocheur.
\item Les 3 bacs de tri avec leur code couleur et 3 exemples de déchets par bac (images ou pictogrammes).
\item Un bénéfice concret pour l'élève / l'établissement.
\item Un slogan court et mémorable.
\end{enumerate}
(Dessine le schéma de l'affiche sur une feuille, ou décris précisément sa mise en page et son contenu.)
\end{exercice}

\begin{corrige}[Exercice Affiche de sensibilisation]
\textbf{Proposition de structure :}
\begin{itemize}
\item \textbf{Titre} (haut, grande police, couleur verte) : « \textbf{TRIER, C'EST GAGNER !} »
\item \textbf{Zone centrale} : 3 colonnes verticales, chaque colonne = un bac.
\begin{itemize}
\item Colonne 1 (fond vert) : Icône bac vert + « ORGANIQUE » + images : épluchures, pain, fruit, herbe. Flèche $\rightarrow$ « COMPOST = Engrais gratuit pour NOTRE jardin ».
\item Colonne 2 (fond jaune) : Icône bac jaune + « RECYCLABLE » + images : bouteille eau, canette, carton. Flèche $\rightarrow$ « RECYCLAGE = Propreté + Revenus club ».
\item Colonne 3 (fond gris) : Icône bac gris + « AUTRES » + images : styromousse, couche, stylo. Flèche $\rightarrow$ « DÉCHETTERIE = Dernier recours ».
\end{itemize}
\item \textbf{Bandeau bas} (fond bleu, police blanche) : « \textbf{1 geste = 1 arbre sauvé, 1 repas au jardin, 1 franc gagné.} »
\item \textbf{Slogan} (très visible, police manuscrite) : « \emph{Je trie, tu tries, nous gagnons !} »
\end{itemize}
\end{corrige}

\begin{attention}[Les 5 erreurs qui coûtent des points au Probatoire]
\begin{enumerate}
\item \textbf{Confondre réutiliser et recycler.} Réutiliser, c'est réemployer l'objet \emph{tel quel} (bouteille consignée, pneu transformé en jardinière, sac en tissu) ; recycler, c'est \emph{transformer la matière} (fonte du verre, plastique fondu en pavés, papier $\rightarrow$ pâte à papier). Citer la consigne comme du recyclage est une erreur classique.
\item \textbf{Classer les piles et médicaments avec les ordures ménagères.} Ce sont des déchets \cle{dangereux} (métaux lourds, substances toxiques, écotoxicité) exigeant une collecte spécialisée ; les oublier dans un tableau de tri fait perdre des points.
\item \textbf{Oublier les conditions biologiques des procédés.} Le compostage se fait \emph{en présence d'air} (aérobie), la méthanisation \emph{en l'absence d'air} (anaérobie, dans un digesteur étanche). Inverser les deux est sanctionné.
\item \textbf{Faire des calculs de pourcentages sans les convertir.} Écris toujours $m = 800 \times \dfrac{75}{100} = 600$ kg, avec la formule littérale, l'application numérique et l'\emph{unité} ; un résultat sans unité ou une masse de produit supérieure à la masse de départ (rendement $> 100\,\%$) trahit une erreur de raisonnement.
\item \textbf{Rester vague sur les impacts.} Écrire « les déchets polluent et rendent malade » sans mécanisme ne rapporte rien. Il faut expliciter la chaîne : sachets $\rightarrow$ caniveaux bouchés $\rightarrow$ inondations ; eaux stagnantes $\rightarrow$ larves d'\emph{Anopheles} $\rightarrow$ paludisme ; lixiviat $\rightarrow$ puits contaminé $\rightarrow$ diarrées/choléra ; métaux lourds $\rightarrow$ bioaccumulation $\rightarrow$ maladies chroniques.
\end{enumerate}
\end{attention}

\begin{aretenir}[Synthèse de la leçon -- Gestion durable des déchets]
\begin{itemize}
\item \textbf{Tri à la source} (organique / recyclable / dangereux) = condition \emph{sine qua non} de toute valorisation.
\item \textbf{Hiérarchie 3R} : Réduire (source) $>$ Réutiliser (objet) $>$ Recycler (matière) $>$ Valorisation énergétique $>$ Élimination.
\item \textbf{Organique} : Compostage (aérobie, compost) OU Méthanisation (anaérobie, biogaz + digestat) -- choix selon gisement, besoin énergie, moyens.
\item \textbf{Recyclables} : Filières matière (verre, métal, papier, plastique) -- nécessitent tri amont rigoureux.
\item \textbf{Dangereux} : Filière spécialisée exclusive -- jamais dans la nature ni au feu.
\item \textbf{Sensibilisation} : Cible + Message positif + Support adapté + Évaluation comportementale = Changement durable.
\item \textbf{Contexte Cameroun} : HYSACAM (collecte), initiatives locales (Kemit Ecology, groupements femmes), enjeux santé (paludisme, choléra) et environnement (déforestation, inondations).
\end{itemize}
\end{aretenir}

\newpage

\coursec{Exercices}

\courssub{Je vérifie mes connaissances}

\begin{exercice}[QCM : Typologie des déchets]
Parmi les propositions suivantes, cochez la ou les bonne(s) réponse(s) pour chaque question. Justifiez brièvement votre choix.

\begin{enumerate}
    \item Un déchet est considéré comme \emph{dangereux} s'il présente au moins une des propriétés suivantes (selon la nomenclature internationale) :
    \begin{itemize}
        \item[A.] Il est biodégradable en moins de 6 mois.
        \item[B.] Il est inflammable, corrosif, toxique, cancérigène ou écotoxique.
        \item[C.] Il provient d'un établissement de santé.
        \item[D.] Il est constitué de plastique non recyclable.
    \end{itemize}

    \item Le tri à la source consiste à :
    \begin{itemize}
        \item[A.] Séparer les déchets après leur collecte par les services municipaux.
        \item[B.] Séparer les déchets dès leur production, au niveau du producteur (ménage, entreprise).
        \item[C.] Brûler les déchets organiques pour réduire leur volume.
        \item[D.] Enfouir tous les déchets dans une décharge contrôlée.
    \end{itemize}

    \item Laquelle de ces affirmations sur le compostage est \textbf{fausse} ?
    \begin{itemize}
        \item[A.] Il nécessite la présence d'oxygène (processus aérobie).
        \item[B.] Il transforme la matière organique en un amendement riche en humus.
        \item[C.] Il produit du méthane ($\ce{CH4}$) comme gaz principal.
        \item[D.] Il réduit le volume initial des déchets d'environ 30 à 50 \%. 
    \end{itemize}
\end{enumerate}
\end{exercice}

\begin{exercice}[Vrai / Faux justifié]
Indiquez si les affirmations suivantes sont Vraies (V) ou Fausses (F). Justifiez chaque réponse par une phrase précise.

\begin{enumerate}
    \item Les sachets plastiques fins (type « sacs cabas » du marché) sont recyclables dans la filière classique de recyclage du plastique au Cameroun.
    \item Le biogaz produit par méthanisation est composé majoritairement de dioxyde de carbone ($\ce{CO2}$).
    \item Les piles usagées (alcalines, bouton) peuvent être jetées dans la poubelle des ordures ménagères résiduelles sans risque pour l'environnement.
    \item La règle des 3R (Réduire, Réutiliser, Recycler) place le recyclage comme la priorité absolue avant la réduction à la source.
\end{enumerate}
\end{exercice}

\begin{exercice}[Définitions et correspondance]
\begin{enumerate}
    \item Donnez une définition précise et scientifique des termes suivants :
    \begin{itemize}
        \item \textbf{Déchet ultime} :
        \item \textbf{Valorisation matière} :
        \item \textbf{Lixiviat (ou jus de décharge)} :
    \end{itemize}
    \item Reliez chaque déchet de la colonne de gauche à sa filière de traitement \textbf{la plus adaptée} dans la colonne de droite (une seule correspondance par déchet).
    \begin{center}
        \begin{tabular}{|l|l|}
            \hline
            \textbf{Déchet} & \textbf{Filière de traitement} \\ \hline
            1. Épluchures de manioc, feuilles de légumes & A. Collecte spécialisée déchets dangereux (ex: centre de traitement) \\
            2. Bouteilles en PET (eau, jus) & B. Compostage domestique ou industriel \\
            3. Pile bouton usagée (montre, calculatrice) & C. Recyclage filière plastique (granulats) \\
            4. Canettes en aluminium (boissons) & D. Recyclage filière verre (calcin) \\
            5. Bouteilles en verre (bière, vin) & E. Recyclage filière métaux (fonderie) \\ \hline
        \end{tabular}
    \end{center}
\end{enumerate}
\end{exercice}

\begin{exercice}[Calcul direct : Rendement de compostage]
Un ménage de Yaoundé produit en moyenne 1,2 kg de déchets de cuisine (organiques) par semaine. Il décide de composter ces déchets dans un composteur domestique. Le rendement massique du compostage (masse de compost mature obtenue / masse de déchets frais mis en entrée) est estimé à 30 \%.

\begin{enumerate}
    \item Calculez la masse de compost mature (en kg) que ce ménage peut espérer obtenir sur une année (52 semaines).
    \item Si ce compost est conditionné en sacs de 5 kg pour être vendu au marché local à 1 500 FCFA le sac, quel revenu annuel brut (en FCFA) le ménage peut-il espérer (arrondi au millier) ?
\end{enumerate}
\end{exercice}

\courssub{Je m'entraîne}

\begin{exercice}[Tri sélectif au marché : Mise en situation]
Vous êtes chargé d'organiser le tri des déchets au marché central de votre quartier. Vous disposez de quatre bacs de couleurs différentes selon le code couleur national (ou local) : \textbf{Vert} (organiques), \textbf{Jaune} (emballages recyclables : plastique, métal, carton), \textbf{Bleu} (verre), \textbf{Gris/Noir} (résidus ultimes).

\begin{enumerate}
    \item Pour chacun des déchets suivants, indiquez dans quel bac il doit être déposé :
    \begin{itemize}
        \item Feilles de bananier et tiges de légumes (déballage) \\
        \item Bouteilles d'eau en plastique vides (PET) \\
        \item Boîtes de conserve vides (tomates, sardines) \\
        \item Bouteilles de bière en verre consignées \\
        \item Papiers souillés (emballages graisseux de poisson, serviettes) \\
        \item Piles usagées des balances électroniques \\
        \item Cartons d'emballage (ondulés) non souillés \\
        \item Sacs plastiques fins (sacs cabas) déchirés
    \end{itemize}
    \item Justifiez pourquoi les papiers souillés (gras) ne vont pas dans le bac jaune (recyclables).
    \item Proposez une solution spécifique pour les piles usagées trouvées au marché.
\end{enumerate}
\end{exercice}

\begin{exercice}[Analyse d'une fiche technique : Méthanisation]
On souhaite installer un petit digesteur de biogaz (méthanisation) dans un lycée agricole disposant d'une porcherie de 50 porcs. La production de fumier frais est estimée à 4 kg par porc et par jour. La Matière Sèche (MS) du fumier représente 20 \% de sa masse fraîche. Le potentiel de méthane (Biochimical Methane Potential - BMP) de ce fumier est de 0,25 $\text{m}^3 \text{CH}_4 / \text{kg MS}$. Le biogaz produit contient 60 \% de méthane ($\ce{CH4}$).

\begin{enumerate}
    \item Calculez la production journalière de Matière Sèche (en kg) disponible pour le digesteur.
    \item Calculez le volume journalier de méthane théorique produit (en $\text{m}^3$).
    \item En déduire le volume journalier de biogaz brut produit (en $\text{m}^3$).
    \item Sachant qu'un réchaud à biogaz consomme environ 0,3 $\text{m}^3$ de biogaz par heure de cuisson, combien d'heures de cuisson par jour ce digesteur peut-il alimenter ?
\end{enumerate}
\end{exercice}

\begin{exercice}[Cycle de vie d'une bouteille en plastique : Analyse critique]
On compare deux scénarios pour une bouteille d'eau en PET de 1,5 L vide :
\begin{itemize}
    \item \textbf{Scénario A :} La bouteille est jetée dans la nature (caniveau). Elle met 400 ans à se fragmenter en microplastiques.
    \item \textbf{Scénario B :} La bouteille est triée (bac jaune), collectée, lavée, broyée en paillettes, puis extrudée en granulats pour fabriquer une polaire (vêtement) ou une nouvelle bouteille (recyclage bouteille-à-bouteille).
\end{itemize}

\begin{enumerate}
    \item Citez trois impacts environnementaux négatifs majeurs du Scénario A (pollution visuelle, physique, chimique).
    \item Pour le Scénario B, citez deux avantages environnementaux majeurs par rapport à la production de PET vierge (pétrole).
    \item Identifiez une limite majeure du recyclage mécanique du PET (bouteille $\rightarrow$ granulats $\rightarrow$ objet) par rapport à une économie circulaire idéale.
\end{enumerate}
\end{exercice}

\begin{exercice}[Argumentaire : Sensibilisation aux déchets dangereux]
Un habitant de votre quartier verse régulièrement l'huile de vidange de sa moto et les restes de peinture dans le caniveau devant chez lui, arguant que « l'eau de pluie emporte tout » et que « la terre filtre tout ».

\begin{enumerate}
    \item Expliquez scientifiquement pourquoi l'affirmation « la terre filtre tout » est fausse dans le cas des hydrocarbures (huile) et des métaux lourds (pigments peinture). Utilisez les termes \emph{infiltration}, \emph{nappe phréatique}, \emph{bioaccumulation}.
    \item Rédigez un court message de sensibilisation (5-6 lignes) à destination de cet habitant, structuré en : \textbf{Constat} $\rightarrow$ \textbf{Risque sanitaire/environnemental} $\rightarrow$ \textbf{Solution concrète locale} (ex: garage agréé, centre de collecte).
\end{enumerate}
\end{exercice}

\courssub{Je me prépare au Probatoire}

\begin{exercice}[Sujet type Probatoire : Gestion intégrée des déchets dans une commune urbaine (12 pts)]
\textbf{Contexte :} La commune de Dschang (Ouest Cameroun) produit environ 45 tonnes de déchets ménagers et assimilés par jour. Une caractérisation réalisée en saison sèche donne la composition gravimétrique moyenne suivante :
\begin{center}
\begin{tabular}{|l|c|}
\hline
\textbf{Fraction} & \textbf{Pourcentage massique (\%)} \\ \hline
Matières organiques (putrescibles : cuisine, marché, jardin) & 62 \\
Plastiques (films, bouteilles, rigides) & 12 \\
Papiers / Cartons & 8 \\
Verre & 4 \\
Métaux (ferreux, aluminium) & 3 \\
Textiles & 3 \\
Autres (poussières, fines, inertes, dangereux) & 8 \\ \hline
\textbf{Total} & \textbf{100} \\ \hline
\end{tabular}
\end{center}
La commune dispose d'une décharge à ciel ouvert non contrôlée en périphérie, source de nuisances (odeurs, fumées, lixiviats vers la rivière Mifi). Le conseil municipal envisage la construction d'un Centre de Valorisation des Déchets (CVD) intégrant une plateforme de compostage et un tri mécanique.

\textbf{Données complémentaires :}
\begin{itemize}
    \item Rendement massique du compostage industriel (matière organique entrante $\rightarrow$ compost mature NF U 44-051) : 35 \%.
    \item Taux de récupération des recyclables (plastiques, métaux, papier/carton, verre) par le tri mécanique + manuel : 80 \% des quantités présentes dans les déchets mélangés.
    \item Prix de vente moyen du compost : 150 FCFA/kg.
    \item Prix de vente moyen des balles de plastiques triés (PET/PE) : 120 FCFA/kg.
    \item Prix de vente moyen des métaux ferreux : 80 FCFA/kg ; Aluminium : 400 FCFA/kg.
    \item Coût d'enfouissement évité (transport + enfouissement en centre technique) : 5 000 FCFA/tonne.
\end{itemize}

\textbf{Questions :}

\begin{enumerate}
    \item \textbf{Caractérisation et enjeux (3 pts)}
    \begin{enumerate}
        \item Calculez la masse journalière (en tonnes) de matières organiques et de plastiques produits par la commune.
        \item Citez deux risques sanitaires et deux risques environnementaux majeurs liés à la décharge à ciel ouvert actuelle, en lien avec la composition des déchets (lixiviats, gaz, vecteurs).
    \end{enumerate}

    \item \textbf{Dimensionnement du Centre de Valorisation (4 pts)}
    \begin{enumerate}
        \item Calculez la production journalière théorique de compost mature (en tonnes) si toutes les matières organiques sont compostées.
        \item Calculez la masse journalière (en tonnes) de recyclables (plastiques + papiers/cartons + métaux + verre) potentiellement récupérables au CVD.
        \item Déduisez la masse journalière de déchets ultimes (résidus) qui devront encore être enfouis après tri et compostage (on suppose que les textiles et « autres » ne sont pas valorisés et partent en résidus, ainsi que les 20 \% de recyclables non récupérés et les 65 \% de pertes du compostage).
    \end{enumerate}

    \item \textbf{Bilan économique et développement durable (3 pts)}
    \begin{enumerate}
        \item Estimez le chiffre d'affaires journalier brut (en FCFA) généré par la vente du compost et des plastiques uniquement (on néglige papier, verre, métaux pour simplifier).
        \item Calculez l'économie journalière réalisée sur le coût d'enfouissement évité grâce à la diversion des organiques et recyclables (masse totale diversionnée $\times$ coût évité).
        \item En vous basant sur les 3 piliers du développement durable (Économique, Social, Environnemental), arguez en 3 points pourquoi ce projet de CVD est pertinent pour la commune de Dschang.
    \end{enumerate}

    \item \textbf{Sensibilisation et gouvernance (2 pts)}
    Proposez deux actions concrètes de communication/éducation à destination des populations de Dschang pour garantir le tri à la source (séparation organiques / secs) indispensable au bon fonctionnement du CVD. Précisez le support et le message clé pour chaque action.
\end{enumerate}
\end{exercice}

\begin{exercice}[Sujet type Probatoire : Pollution plastique et solutions locales (10 pts)]
\textbf{Contexte :} Le fleuve Wouri (Douala) charrie d'importantes quantités de déchets plastiques vers l'océan Atlantique. Une ONG locale, « \emph{Plastic Free Wouri} », mène une campagne de nettoyage des berges et de sensibilisation. Lors d'une opération de ramassage sur 500 m de berge, les bénévoles ont collecté et trié 1 200 kg de déchets. La répartition est la suivante : 55 \% bouteilles PET, 20 \% sachets PE (sacs cabas), 10 \% bouchons PP/PE, 5 \% polystyrène expansé (boîtes repas), 10 \% autres (verre, textile, bois).

\textbf{Données :}
\begin{itemize}
    \item Densité du PET : 1,38 g/cm$^3$ ; Densité du PE : 0,94 g/cm$^3$.
    \item Une bouteille PET 1,5 L pèse environ 35 g vide. Un sachet PE pèse environ 5 g.
    \item Le recyclage mécanique du PET permet d'économiser environ 2,5 kg de pétrole brut et 60 \% d'énergie par kg de PET recyclé vs PET vierge.
    \item Le Cameroun a interdit les sachets plastiques non biodégradables d'épaisseur $< 60 \mu m$ (Décret 2012), mais l'application reste difficile.
\end{itemize}

\textbf{Questions :}

\begin{enumerate}
    \item \textbf{Quantification et devenir (3 pts)}
    \begin{enumerate}
        \item Calculez le nombre approximatif de bouteilles PET et de sachets PE collectés lors de cette opération.
        \item Expliquez pourquoi les sachets PE (films fins) posent un problème technique majeur pour le tri automatisé et le recyclage mécanique industriel (encrassement, enroulement), contrairement aux bouteilles PET rigides.
        \item Si 80 \% des bouteilles PET collectées sont effectivement recyclées en granulats, calculez la masse de pétrole brut économisée (en kg) grâce à cette opération de nettoyage.
    \end{enumerate}

    \item \textbf{Impacts écotoxicologiques (3 pts)}
    \begin{enumerate}
        \item Décrivez le processus de fragmentation d'une bouteille PET en microplastiques ($< 5$ mm) dans le milieu aquatique (facteurs physiques, chimiques, biologiques).
        \item Citez deux mécanismes par lesquels les microplastiques ingérés par les poissons (ex: \emph{Barbus}, \emph{Silure}) peuvent affecter leur physiologie (ex: fausse satiété, transfert de polluants adsorbés).
        \item En quoi la présence de perturbateurs endocriniens (phtalates, bisphénol A) adsorbés sur les microplastiques représente-t-elle un risque pour la santé humaine via la chaîne alimentaire (bioamplification) ?
    \end{enumerate}

    \item \textbf{Solutions intégrées et réglementation (4 pts)}
    \begin{enumerate}
        \item Proposez une solution technique locale de valorisation matière pour les sachets PE collectés (non recyclables classiquement), en citant un exemple réel ou plausible au Cameroun (ex: pavés plastiques, bois plastique, pyrolyse). Décrivez brièvement le principe.
        \item Au-delà du nettoyage curatif, proposez trois mesures préventives (réglementaires, économiques, éducatives) que les autorités locales (Commune, MINEPDED) pourraient mettre en œuvre pour réduire drastiquement l'arrivée de plastiques dans le Wouri. Justifiez l'efficacité attendue de chaque mesure.
        \item Rédigez le slogan principal (percutant, 8 mots max) et les trois messages clés (une phrase chacun) d'une affiche de sensibilisation destinée aux piroguiers et riverains du Wouri, intégrant la notion de « \emph{bien commun} » et de « \emph{santé publique} ».
    \end{enumerate}
\end{enumerate}
\end{exercice}

\begin{exercice}[Sujet type Probatoire : Compostage et agriculture urbaine durable (8 pts)]
\textbf{Contexte :} À Yaoundé, le projet « \emph{Terre Verte} » collecte les déchets de cuisine des restaurants du quartier Mfoundi et les fientes de poulaillers périurbains pour produire du compost vendu aux jardiniers maraîchers des bas-fonds (zones marécageuses cultivées).

\textbf{Données sur le mélange initial (recette) :}
\begin{itemize}
    \item 60 \% (masse fraîche) de déchets de cuisine (légumes, fruits, épluchures) : Teneur en eau 85 \%, Rapport C/N = 15.
    \item 40 \% (masse fraîche) de fientes de poules (fumier) : Teneur en eau 60 \%, Rapport C/N = 10.
\end{itemize}
Le processus de compostage en andins aérés dure 90 jours. Les analyses du compost mature (norme NF U 44-051) donnent : Humidité 40 \%, Matière Sèche (MS) 60 \%, Matière Organique (MO) sur MS = 45 \%, Azote total (N) sur MS = 2 \%, Rapport C/N final = 12.

\textbf{Questions :}

\begin{enumerate}
    \item \textbf{Équilibrage de la recette (2 pts)}
    \begin{enumerate}
        \item Calculez le rapport C/N global du mélange initial (pondéré par les masses de Matière Sèche). Le mélange est-il équilibré pour un démarrage rapide du compostage (C/N optimal 25-35) ? Justifiez.
        \item Proposez une modification de la recette (ajout de carbone) pour atteindre un C/N initial de 30, en utilisant de la sciure de bois (C/N = 400, MS = 90 \%). Calculez le pourcentage de sciure à ajouter (en masse fraîche) par rapport au mélange initial total.
    \end{enumerate}

    \item \textbf{Suivi du processus et qualité (3 pts)}
    \begin{enumerate}
        \item Citez trois paramètres physico-chimiques à surveiller \textbf{régulièrement} pendant les 90 jours pour garantir une hygiénisation (montée en température $> 55^\circ\text{C}$) et une bonne dégradation. Pour chacun, donnez la valeur cible optimale.
        \item Le compost mature présente un rapport C/N = 12. Interprétez cette valeur au regard de la stabilité du compost et du risque de « faim d'azote » lors de l'épandage.
        \item L'analyse révèle une teneur en \emph{E. coli} de 500 UFC/g de compost frais. La norme NF U 44-051 impose $< 1000$ UFC/g. Le compost est-il conforme sur le plan hygiénique ? Expliquez le rôle de la phase thermophile ($> 55^\circ\text{C}$ pendant plusieurs jours) dans cette hygiénisation.
    \end{enumerate}

    \item \textbf{Modèle économique et impact agricole (3 pts)}
    \begin{enumerate}
        \item Le projet traite 2 tonnes de mélange initial par jour. Sachant une perte de masse totale (eau + $\ce{CO2}$) de 55 \% à la fin du processus, calculez la production mensuelle de compost mature (30 jours) en tonnes.
        \item Si le compost est vendu 100 FCFA/kg et que les charges d'exploitation (main d'œuvre, carburant broyeur, location, analyses) s'élèvent à 1 500 000 FCFA/mois, calculez le résultat net mensuel (bénéfice ou perte).
        \item Un maraîcher épand 10 t/ha de ce compost (base fraîche). Calculez les apports en unités d'azote (kg N/ha) et de matière organique (t MO/ha) apportés par cet épandage. Discutez brièvement de l'intérêt agronomique de cet apport sur des sols de bas-fonds souvent sableux et pauvres en humus.
    \end{enumerate}
\end{enumerate}
\end{exercice}

\newpage

\coursec{Corrigés des exercices}

\courssub{Je vérifie mes connaissances}

\begin{corrige}[Exercice 1 — QCM : Typologie des déchets]
\begin{enumerate}
    \item \textbf{Bonne réponse : B.}
    \textbf{Justification :} Selon la nomenclature internationale (ex: Directive européenne 2008/98/CE, reprise dans la réglementation camerounaise), un déchet est classé \emph{dangereux} s'il possède une ou plusieurs propriétés de danger (H1 à H15) : explosif, comburant, inflammable, irritant, toxique, cancérigène, corrosif, infectieux, mutagène, écotoxique, etc. 
    \begin{itemize}
        \item A est faux : la biodégradabilité caractérise les déchets organiques (non dangereux le plus souvent).
        \item C est faux : l'origine (établissement de santé) définit les DASRI (Déchets d'Activités de Soins à Risques Infectieux), une sous-catégorie, mais tous les déchets de santé ne sont pas dangereux (emballages, papier bureautique) et des déchets dangereux existent hors santé (industriels, ménagers : piles, solvants).
        \item D est faux : le plastique non recyclable est un déchet ultime ou à valoriser énergétiquement, pas nécessairement dangereux (sauf additifs spécifiques).
    \end{itemize}

    \item \textbf{Bonne réponse : B.}
    \textbf{Justification :} Le \emph{tri à la source} est l'action de séparer les déchets \textbf{dès leur lieu de production} (cuisine, atelier, bureau, marché) par le producteur lui-même, avant toute collecte. C'est la condition \emph{sine qua non} d'un recyclage de qualité.
    \begin{itemize}
        \item A décrit le tri en centre de tri (post-collecte), moins efficace car les déchets sont mélangés, souillés, cassés.
        \item C décrit l'incinération (avec ou sans valorisation énergétique), pas du tri.
        \item D décrit l'enfouissement (décharge), mode d'élimination final pour les ultimes.
    \end{itemize}

    \item \textbf{Bonne réponse : C (affirmation fausse).}
    \textbf{Justification :} Le compostage est un processus \emph{aérobie} (nécessitant O$_2$, affirmation A vraie) qui produit principalement du \ce{CO2}, de l'eau, de la chaleur et de l'humus (affirmation B vraie). Le \textbf{méthane (\ce{CH4})} est le produit majeur de la \emph{méthanisation} (digestion anaérobie, sans O$_2$). En compostage, la production de \ce{CH4} est résiduelle (zones anoxiques accidentelles) et signe un mauvais fonctionnement (manque de retournement). L'affirmation D est vraie : la perte de masse (eau + \ce{CO2}) est de 30 à 50\%.
\end{enumerate}
\end{corrige}

\begin{corrige}[Exercice 2 — Vrai / Faux justifié]
\begin{enumerate}
    \item \textbf{Faux.}
    \textbf{Justification :} Au Cameroun, les sachets plastiques fins (épaisseur < 60 µm, dits « sacs cabas ») sont \textbf{interdits} à la production, l'importation et la commercialisation (Décret n°2012/2631/PM du 24 août 2012). Ils ne sont \textbf{pas recyclables} dans la filière classique (ils s'enroulent dans les broyeurs, encrassent les chaînes de tri, sont trop sales). Ils finissent en résidus ultimes (poubelle grise/noire) ou, hélas, dans la nature.

    \item \textbf{Faux.}
    \textbf{Justification :} Le biogaz issu de la méthanisation est composé majoritairement de \textbf{méthane (\ce{CH4})} (généralement 50 à 70\%), le \ce{CO2} étant le second composant (30 à 50\%). C'est la teneur en \ce{CH4} qui lui confère son pouvoir calorifique (environ 6 kWh/m$^3$ pour 60\% \ce{CH4}).

    \item \textbf{Faux.}
    \textbf{Justification :} Les piles usagées (alcalines, bouton, lithium, Ni-Cd) contiennent des \textbf{métaux lourds} (mercure, cadmium, plomb, nickel, zinc) et des électrolytes corrosifs. Jetées en décharge ordinaire, leurs enveloppes se corrodent, relarguant ces polluants dans les \emph{lixiviats} (jus de décharge) qui contaminent sols et nappes phréatiques. Elles relèvent d'une \textbf{filiaire de collecte sélective spécifique} (DEEE / déchets dangereux ménagers) vers des centres de traitement (ex: centres de regroupement HYSACAM ou partenaires spécialisés).

    \item \textbf{Faux.}
    \textbf{Justification :} La règle des \textbf{3R} établit une \textbf{hiérarchie stricte} : 
    \textbf{1. Réduire} (priorité absolue : éviter de produire le déchet), 
    \textbf{2. Réutiliser} (prolonger la durée de vie), 
    \textbf{3. Recycler} (traiter la matière). 
    Le recyclage coûte de l'énergie et de l'eau ; il ne doit intervenir qu'après avoir épuisé les possibilités de réduction et réutilisation.
\end{enumerate}
\end{corrige}

\begin{corrige}[Exercice 3 — Définitions et correspondance]
\begin{enumerate}
    \item \textbf{Définitions précises :}
    \begin{itemize}
        \item \textbf{Déchet ultime :} Déchet qui ne peut plus être valorisé ni dans les conditions techniques et économiques actuelles, ni dans un avenir prévisible, notamment par recyclage, compostage ou valorisation énergétique. Il ne peut subir qu'un stockage final en installation de stockage des déchets non dangereux (ISDND) ou dangereux (ISDD).
        \item \textbf{Valorisation matière :} Opération de traitement des déchets permettant de récupérer la matière qui les constitue pour la réutiliser comme matière première secondaire (ex: broyage de bouteilles PET en paillettes/granulats pour refaire du plastique ; compostage transformant l'organique en amendement). S'oppose à la valorisation énergétique (incinération avec récupération chaleur).
        \item \textbf{Lixiviat (ou jus de décharge) :} Liquide pollué, chargé en matière organique dissoute, ammoniac, métaux lourds, sels minéraux et agents pathogènes, résultant de la percolation de l'eau de pluie (et de l'humidité des déchets) à travers la masse des déchets en décharge. Sa collecte et son traitement sont obligatoires en centre technique (ISDND).
    \end{itemize}

    \item \textbf{Correspondances (Déchet $\rightarrow$ Filière adaptée) :}
    \begin{center}
        \begin{tabular}{|l|l|l|}
            \hline
            \textbf{Déchet} & \textbf{Numéro} & \textbf{Filière (Lettre)} \\ \hline
            Épluchures de manioc, feuilles de légumes & 1 & \textbf{B} (Compostage domestique/industriel) \\ \hline
            Bouteilles en PET (eau, jus) & 2 & \textbf{C} (Recyclage filière plastique → granulats) \\ \hline
            Pile bouton usagée & 3 & \textbf{A} (Collecte spécialisée déchets dangereux) \\ \hline
            Canettes en aluminium & 4 & \textbf{E} (Recyclage filière métaux → fonderie) \\ \hline
            Bouteilles en verre & 5 & \textbf{D} (Recyclage filière verre → calcin) \\ \hline
        \end{tabular}
    \end{center}
\end{enumerate}
\end{corrige}

\begin{corrige}[Exercice 4 — Calcul direct : Rendement de compostage]
\textbf{Données :} Production hebdomadaire $m_{\text{sem}} = 1,2\,\text{kg}$ ; Rendement massique $R = 30\% = 0,30$ ; Semaines $n = 52$ ; Prix sac $5\,\text{kg} = 1\,500\,\text{FCFA}$.

\begin{enumerate}
    \item \textbf{Masse annuelle de déchets frais :}
    \[ m_{\text{an}} = m_{\text{sem}} \times n = 1,2 \times 52 = 62,4\,\text{kg} \]
    \textbf{Masse de compost mature :}
    \[ m_{\text{compost}} = m_{\text{an}} \times R = 62,4 \times 0,30 = \boxed{18,72\,\text{kg}} \]

    \item \textbf{Prix au kg :} $p_{\text{kg}} = \dfrac{1\,500}{5} = 300\,\text{FCFA/kg}$.
    \textbf{Revenu brut annuel :}
    \[ \text{Revenu} = m_{\text{compost}} \times p_{\text{kg}} = 18,72 \times 300 = 5\,616\,\text{FCFA} \]
    \textbf{Arrondi au millier :} $\boxed{6\,000\,\text{FCFA}}$.
    \textit{(Note : si vente uniquement par sacs entiers de 5 kg : 3 sacs = 15 kg vendus $\rightarrow$ 4\,500 FCFA $\approx$ 5\,000 FCFA. L'énoncé « conditionné en sacs... revenu espéré » suggère la valorisation de la production totale, d'où le calcul au kg.)}
\end{enumerate}
\end{corrige}

\courssub{Je m'entraîne}

\begin{corrige}[Exercice 5 — Tri sélectif au marché : Mise en situation]
\begin{enumerate}
    \item \textbf{Affectation des bacs (Code couleur national Cameroun) :}
    \begin{center}
        \begin{tabular}{|l|l|l|}
            \hline
            \textbf{Déchet} & \textbf{Bac (Couleur)} & \textbf{Justification brève} \\ \hline
            Feuilles de bananier, tiges de légumes & \textbf{Vert} (Organiques) & Matière putrescible, compostable \\ \hline
            Bouteilles d'eau PET vides & \textbf{Jaune} (Emballages recyclables) & Plastique rigide recyclable (PET) \\ \hline
            Boîtes de conserve vides & \textbf{Jaune} (Emballages recyclables) & Métal (fer blanc/aluminium) recyclable \\ \hline
            Bouteilles de bière en verre consignées & \textbf{Bleu} (Verre) \textit{ou retour consigne} & Verre d'emballage recyclable à l'infini \\ \hline
            Papiers souillés (gras, poisson) & \textbf{Gris/Noir} (Résidus ultimes) & Souillés = non recyclables (fibres collées) \\ \hline
            Piles usagées balances & \textbf{Collecte spéciale} (Point d'apport) & Déchets dangereux (métaux lourds) \\ \hline
            Cartons ondulés non souillés & \textbf{Jaune} (Emballages recyclables) & Carton propre, recyclable \\ \hline
            Sacs plastiques fins déchirés & \textbf{Gris/Noir} (Résidus ultimes) & Films fins non recyclables (interdits <60µm) \\ \hline
        \end{tabular}
    \end{center}

    \item \textbf{Justification papiers souillés :} Le recyclage papier/carton repose sur le \emph{défibrage} en eau (pulpeur). Les graisses (huiles, sauces) sont \textbf{hydrophobes} : elles empêchent la séparation des fibres, créent des « bâtons » (taches grasses) sur la feuille de papier recyclé, et perturbent le désencrage. Un papier gras contamine toute la balle de tri. $\rightarrow$ Bac résidus (ou compostage si très peu d'encre, mais généralement résidus au marché).

    \item \textbf{Solution pour piles usagées :} Mettre en place un \textbf{bac de collecte spécifique « Piles/Accus »} (souvent boîte verte/transparente) à l'entrée du marché ou chez les revendeurs de piles/électronique, avec affichage clair. Convention avec un collecteur agréé (ex: HYSACAM, ONG, entreprise spécialisée) pour enlèvement régulier vers centre de traitement (ex: Douala, Yaoundé) ou export filière spécialisée. Sensibiliser les commerçants (balances) à ne pas les jeter dans les bacs courants.
\end{enumerate}
\end{corrige}

\begin{corrige}[Exercice 6 — Analyse d'une fiche technique : Méthanisation]
\textbf{Données :} $N_{\text{porcs}} = 50$ ; $m_{\text{fumier frais/porc/j}} = 4\,\text{kg}$ ; $\%MS = 20\%$ ; $\text{BMP} = 0,25\,\text{m}^3\,\ce{CH4}/\text{kg MS}$ ; $\%\ce{CH4}$ (biogaz) $= 60\%$ ; Conso réchaud $= 0,3\,\text{m}^3/\text{h}$.

\begin{enumerate}
    \item \textbf{Production journalière de Matière Sèche (MS) :}
    \[ m_{\text{fumier frais}} = 50 \times 4 = 200\,\text{kg/j} \]
    \[ m_{\text{MS}} = 200 \times 0,20 = \boxed{40\,\text{kg MS/j}} \]

    \item \textbf{Volume journalier de méthane théorique :}
    \[ V_{\ce{CH4}} = m_{\text{MS}} \times \text{BMP} = 40 \times 0,25 = \boxed{10\,\text{m}^3\,\ce{CH4}/\text{j}} \]

    \item \textbf{Volume journalier de biogaz brut :}
    \[ V_{\text{biogaz}} = \frac{V_{\ce{CH4}}}{0,60} = \frac{10}{0,6} = \boxed{16,67\,\text{m}^3/\text{j}} \quad (\approx 16,7\,\text{m}^3/\text{j}) \]

    \item \textbf{Heures de cuisson par jour :}
    \[ t_{\text{cuisson}} = \frac{V_{\text{biogaz}}}{0,3} = \frac{16,67}{0,3} = \boxed{55,6\,\text{h/j}} \]
    \textit{Interprétation :} Ce volume permet d'alimenter \textbf{plusieurs foyers/réchauds} (ex: 3 réchauds $\times$ 6h = 18h ; ou cantine scolaire). Un seul réchaud ne peut pas brûler 55h/jour ! Nécessité de stockage (gasholder) ou multi-usages (éclairage, groupe électrogène).
\end{enumerate}
\end{corrige}

\begin{corrige}[Exercice 7 — Cycle de vie d'une bouteille en plastique : Analyse critique]
\begin{enumerate}
    \item \textbf{Trois impacts majeurs du Scénario A (abandon nature/caniveau) :}
    \begin{enumerate}
        \item \textbf{Pollution physique / Visuelle :} Encombrement des caniveaux $\rightarrow$ obstruction écoulement eaux pluviales $\rightarrow$ inondations urbaines (Yaoundé, Douala) ; dégradation paysagère, nuisance olfactive (eaux stagnantes).
        \item \textbf{Pollution physique / Faune :} Ingestion par animaux (bovins, chèvres, oiseaux, poissons) $\rightarrow$ occlusion intestinale, fausse satiété, mort. Piégeage (anneaux de bouteilles).
        \item \textbf{Pollution chimique / Microplastiques :} Fragmentation lente (UV, mécanique, bio) en \textbf{microplastiques} ($<5\,\text{mm}$) $\rightarrow$ relargage additifs (phtalates, bisphénol A, retardateurs de flamme) + adsorption polluants hydrophobes (HAP, PCB, pesticides) $\rightarrow$ transfert chaîne alimentaire jusqu'à l'homme.
    \end{enumerate}

    \item \textbf{Deux avantages majeurs du Scénario B (recyclage mécanique) vs PET vierge (pétrole) :}
    \begin{enumerate}
        \item \textbf{Économie de ressources fossiles non renouvelables :} $\approx 2,5\,\text{kg}$ pétrole brut économisés par kg PET recyclé (évite extraction, raffinage, polymérisation).
        \item \textbf{Réduction forte de l'empreinte carbone/énergie :} $\approx 60\%$ d'énergie primaire en moins (évite étapes craquage, purification, polymérisation haute T/P). Émissions \ce{CO2} évitées $\approx 2\text{--}3\,\text{kg}\ce{CO2eq}/\text{kg PET}$.
    \end{enumerate}

    \item \textbf{Limite majeure du recyclage mécanique (bouteille $\rightarrow$ granulats $\rightarrow$ objet) :}
    \textbf{Dégradation des propriétés polymères (downcycling).} Chaque cycle de chauffe/extrusion casse les chaînes macromoléculaires (baisse poids moléculaire, viscosité) $\rightarrow$ perte de résistance mécanique, transparence, pureté. 
    \begin{itemize}
        \item Bouteille $\rightarrow$ Granulats $\rightarrow$ Polaire / Tapis / Sangle (pas bouteille alimentaire sans décontamination lourde).
        \item Bouteille $\rightarrow$ Bouteille (bouteille-à-bouteille) possible mais limité (besoin \emph{superclean} process, ajout PET vierge).
    \end{itemize}
    $\rightarrow$ Pas de circularité infinie « fermée » comme le verre ou l'aluminium. Nécessite \textbf{recyclage chimique (dépolymérisation)} pour boucler vraiment la boucle.
\end{enumerate}
\end{corrige}

\begin{corrige}[Exercice 8 — Argumentaire : Sensibilisation aux déchets dangereux]
\begin{enumerate}
    \item \textbf{Explication scientifique :}
    L'affirmation « la terre filtre tout » est \textbf{fausse} pour les hydrocarbures et métaux lourds.
    \begin{itemize}
        \item \textbf{Infiltration :} L'huile de vidange (hydrocarbures) et les solvants/métaux de la peinture sont \textbf{lipophiles} et peu solubles dans l'eau. Ils ne sont pas retenus par la matrice argileuse (échange cationique inefficace pour molécules neutres/hydrophobes). Ils migrent verticalement par \emph{infiltration} préférentielle (fissures, macropores) vers la \textbf{nappe phréatique}.
        \item \textbf{Non-biodégradabilité / Persistance :} Les métaux lourds (Pb, Cd, Cr, Hg des pigments) sont des \textbf{éléments chimiques} : ils ne se dégradent pas. Ils s'adsorbent sur la matière organique/argiles mais peuvent être remobilisés (changement pH, redox). Les hydrocarbures (HAP) sont persistants, cancérigènes.
        \item \textbf{Bioaccumulation / Bioamplification :} Une fois dans la nappe $\rightarrow$ eau de forage/puits $\rightarrow$ consommation humaine/animale. Ou nappe $\rightarrow$ rivière $\rightarrow$ plancton $\rightarrow$ poisson $\rightarrow$ homme. Les métaux (Hg, Cd) et HAP s'accumulent dans les graisses/foie/reins (\emph{bioaccumulation}) et se concentrent aux niveaux trophiques supérieurs (\emph{bioamplification}). Doses chroniques faibles $\rightarrow$ cancers, troubles neurologiques, rénaux, endocriniens.
    \end{itemize}

    \item \textbf{Message de sensibilisation (5-6 lignes) :}
    \begin{quote}
    \textbf{Constat :} Monsieur, verser l'huile de vidange et la peinture dans le caniveau pollue directement notre eau et nos sols. \\
    \textbf{Risque :} Ces produits contiennent des métaux lourds et hydrocarbures cancérigènes qui infiltrent la nappe phréatique, contaminent l'eau de boisson et s'accumulent dans les poissons que nous mangeons (bioamplification), mettant en danger la santé de vos enfants et du quartier. \\
    \textbf{Solution :} Apportez vos huiles usagées au garage agréé \textit{[Nom garage local]} ou au point de collecte HYSACAM/Commune le plus proche (ex: Centre de regroupement \textit{[Quartier]}) ; les pots de peinture vides/pleins vont en déchetterie déchets dangereux. Protégeons notre bien commun : l'eau !
    \end{quote}
\end{enumerate}
\end{corrige}

\courssub{Je me prépare au Probatoire}

\begin{corrige}[Exercice 9 — Sujet type Probatoire : Gestion intégrée des déchets à Dschang (12 pts)]
\textbf{Barème indicatif :} Q1 (3 pts) + Q2 (4 pts) + Q3 (3 pts) + Q4 (2 pts) = 12 pts.

\begin{enumerate}
    \item \textbf{Caractérisation et enjeux (3 pts)}
    \begin{enumerate}
        \item \textbf{Calculs masses journalières (Total = 45 t/j) :}
        \begin{align*}
            \text{Organiques (62\%)} &: 45 \times 0,62 = \boxed{27,9\,\text{t/j}} \\
            \text{Plastiques (12\%)} &: 45 \times 0,12 = \boxed{5,4\,\text{t/j}}
        \end{align*}
        \item \textbf{Risques décharge à ciel ouvert (liens composition) :}
        \begin{itemize}
            \item \textbf{Santé 1 :} Prolifération \emph{vecteurs} (mouches, moustiques, rats) attirés par putrescibles (62\%) $\rightarrow$ maladies diarrhéiques, paludisme, leptospirose, fièvre typhoïde pour riverains.
            \item \textbf{Santé 2 :} Fumées toxiques (brûlages spontanés ou volontaires) : dioxines/furannes (plastiques PVC, 12\%), particules fines (PM2,5) $\rightarrow$ affections respiratoires aiguës/chroniques (asthme, BPCO).
            \item \textbf{Environnement 1 :} \textbf{Lixiviats} concentrés (charge organique forte DBO/DCO, \ce{NH4+}, métaux lourds piles/déchets dangereux « autres » 8\%) $\rightarrow$ pollution rivière Mifi (eutrophisation, mortalité poisson, eau impropre aval).
            \item \textbf{Environnement 2 :} \textbf{Gaz à effet de serre} : Méthanisation anarchique anaérobie de l'organique (62\%) $\rightarrow$ émissions \ce{CH4} (GES 28$\times$ \ce{CO2}) + \ce{CO2} $\rightarrow$ changement climatique ; odeurs \ce{H2S}, mercaptans.
        \end{itemize}
    \end{enumerate}

    \item \textbf{Dimensionnement du CVD (4 pts)}
    \begin{enumerate}
        \item \textbf{Compost mature théorique (rendement 35\% sur organiques entrantes) :}
        \[ m_{\text{compost}} = 27,9 \times 0,35 = \boxed{9,765\,\text{t/j}} \quad (\approx 9,8\,\text{t/j}) \]
        \item \textbf{Recyclables récupérables (taux 80\% sur gisement triable) :}
        \begin{align*}
            \text{Gisement total recyclables} &= \text{Plastiques} + \text{Papier/Carton} + \text{Métaux} + \text{Verre} \\
            &= 5,4 + (45\times0,08) + (45\times0,03) + (45\times0,04) \\
            &= 5,4 + 3,6 + 1,35 + 1,8 = \boxed{12,15\,\text{t/j}} \\
            \text{Récupérés (80\%)} &= 12,15 \times 0,80 = \boxed{9,72\,\text{t/j}}
        \end{align*}
        \item \textbf{Déchets ultimes (résidus) à enfouir :}
        \begin{itemize}
            \item Pertes compostage (65\% organiques) : $27,9 \times 0,65 = 18,135\,\text{t/j}$ (eau + \ce{CO2} évaporés, mais résidus secs + inertes restent $\approx$ 30-40\% masse initiale ? \textit{Attention : le rendement 35\% est massique entrée$\rightarrow$sortie compost. La perte de masse 65\% part en \ce{CO2}/\ce{H2O} (gaz) + résidus de tri (plastiques, verre, pierres dans organiques). En bilan de masse entrée/sortie : Entrée 45 t. Sorties solides : Compost (9,765) + Recyclables récupérés (9,72) + Résidus. Les gaz (\ce{CO2}, \ce{H2O}, \ce{CH4}) sortent du bilan solide.})
        \end{itemize}
        \textbf{Bilan de masse solide (Entrées = Sorties solides) :}
        \begin{align*}
            \text{Entrée totale} &= 45,00\,\text{t/j} \\
            \text{Sorties valorisées (solides)} &= \text{Compost} + \text{Recyclables récup.} = 9,765 + 9,72 = 19,485\,\text{t/j} \\
            \text{Résidus ultimes (enfouissement)} &= 45,00 - 19,485 = \boxed{25,515\,\text{t/j}} \quad (\approx 25,5\,\text{t/j})
        \end{align*}
        \textit{Vérification par flux résidus :} 
        Organiques non compostés (résidus tri + pertes solides) + Recyclables non récupérés (20\%) + Textiles (1,35) + Autres (3,6) 
        $= (27,9 - 9,765)_{\text{partie solide résiduelle}} + 2,43 + 1,35 + 3,6$. 
        Si on considère que les 65\% pertes compostage sont \ce{CO2}/\ce{H2O} (gaz), le résidu solide de la plateforme compostage est $\approx$ 10-15\% entrée. Mais l'énoncé dit « 65\% de pertes du compostage » (masse). On fait le bilan global simple : Entrée - Sorties valorisées = Résidus.
    \end{enumerate}

    \item \textbf{Bilan économique et DD (3 pts)}
    \begin{enumerate}
        \item \textbf{Chiffre d'affaires journalier (Compost + Plastiques seulement) :}
        \begin{align*}
            \text{Compost} &: 9,765\,\text{t} = 9\,765\,\text{kg} \times 150\,\text{FCFA/kg} = 1\,464\,750\,\text{FCFA} \\
            \text{Plastiques récup.} &: 5,4\,\text{t} \times 0,80 = 4,32\,\text{t} = 4\,320\,\text{kg} \times 120\,\text{FCFA/kg} = 518\,400\,\text{FCFA} \\
            \text{CA Total} &= \boxed{1\,983\,150\,\text{FCFA/j}} \quad (\approx 1,98\,\text{millions FCFA/j})
        \end{align*}
        \item \textbf{Économie enfouissement évité :}
        \[ \text{Masse diversionnée} = \text{Organiques traitées} + \text{Recyclables récup.} = 27,9 + 9,72 = 37,62\,\text{t/j} \]
        \[ \text{Économie} = 37,62 \times 5\,000 = \boxed{188\,100\,\text{FCFA/j}} \]
        \item \textbf{Pertinence CVD (3 piliers DD) :}
        \begin{enumerate}
            \item \textbf{Économique :} Revenus vente compost/recyclables ($\approx$ 2 M FCFA/j) + économie enfouissement ($\approx$ 0,19 M FCFA/j) $>$ Coûts exploitation (main d'œuvre, énergie, maintenance) $\rightarrow$ Autofinancement partiel/total, création emplois locaux (trieurs, composteurs, chauffeurs).
            \item \textbf{Social :} Amélioration salubrité publique (suppression décharge à ciel ouvert, vecteurs, odeurs, fumées) $\rightarrow$ santé riverains ; intégration trieurs informels (reconnaissance, EPI, revenus stables) ; compost abordable pour maraîchers locaux (sécurité alimentaire).
            \item \textbf{Environnemental :} Réduction drastique lixiviats (Mifi sauvée) ; réduction GES (évitement \ce{CH4} décharge, séquestration C dans sols via compost) ; économie ressources (pétrole, minerais, eau) par recyclage ; retour matière organique au sol (lutte désertification).
        \end{enumerate}
    \end{enumerate}

    \item \textbf{Sensibilisation et gouvernance (2 pts)}
    \begin{enumerate}
        \item \textbf{Action 1 (Porte-à-porte / Chef de quartier) :} \textit{Support :} Réunions quartier + affiches bilingues (Français/Ghomala') + kits « 2 bacs » (seau vert cuisine + sac jaune sec) distribués. \textit{Message clé :} « \textbf{Séparez à la cuisine : épluchures dans le vert, emballages dans le jaune = compost pour nos champs, argent pour la commune !} »
        \item \textbf{Action 2 (Marché / Radio locale) :} \textit{Support :} Spots radio « Radio Medumba » + démonstration publique au marché central (tri en direct par animateurs) + panneaux « Ici on trie ». \textit{Message clé :} « \textbf{Dschang propre, c'est mon affaire : je trie mes déchets, je protège la Mifi, je crée de l'emploi pour nos jeunes.} »
    \end{enumerate}
\end{enumerate}
\textbf{Points de vigilance Probatoire :} 
\begin{itemize}
    \item Unités : tonnes vs kg (attention conversions $\times 1000$ pour prix FCFA/kg).
    \item Rendement compostage : s'applique sur masse \textbf{entrante organique}, pas sur total déchets.
    \item Taux récupération 80\% : s'applique sur \textbf{quantité présente dans déchets mélangés} (gisement), pas sur total déchets.
    \item Bilan masse : Résidus = Entrée - (Compost + Recyclables récupérés). Les pertes gazeuses (\ce{CO2}, \ce{H2O}) sortent du système solide.
    \item Argumentaire DD : 1 point par pilier, phrase complète avec lien causal clair (projet $\rightarrow$ impact).
\end{itemize}
\end{corrige}

\begin{corrige}[Exercice 10 — Sujet type Probatoire : Pollution plastique Wouri (10 pts)]
\textbf{Barème indicatif :} Q1 (3 pts) + Q2 (3 pts) + Q3 (4 pts) = 10 pts.

\begin{enumerate}
    \item \textbf{Quantification et devenir (3 pts)}
    \begin{enumerate}
        \item \textbf{Nombre bouteilles PET et sachets PE :}
        \begin{align*}
            \text{Masse totale} &= 1\,200\,\text{kg} \\
            \text{PET (55\%)} &: 1\,200 \times 0,55 = 660\,\text{kg} = 660\,000\,\text{g} \\
            N_{\text{bouteilles}} &= \frac{660\,000}{35} = \boxed{18\,857\,\text{bouteilles}} \quad (\approx 18\,900) \\
            \text{PE sachets (20\%)} &: 1\,200 \times 0,20 = 240\,\text{kg} = 240\,000\,\text{g} \\
            N_{\text{sachets}} &= \frac{240\,000}{5} = \boxed{48\,000\,\text{sachets}}
        \end{align*}
        \item \textbf{Problème technique sachets PE (films fins) :}
        \begin{itemize}
            \item \textbf{Enroulement / Enchevêtrement :} Films souples, légers, grande surface $\rightarrow$ s'enroulent sur arbres de broyeurs, vis d'extrudeuse, tambours de séparation densimétrique $\rightarrow$ arrêts fréquents, maintenance coûteuse, risque opérateur.
            \item \textbf{Encrassement / Tri optique :} Se collent sur bandes transporteuses, bouchent cribles, perturbent soufflerie (trop légers) et tri optique NIR (transparents, sales, multicouches).
            \item \textbf{Contamination :} Souvent sales (restes alimentaires, sable) $\rightarrow$ qualité granulats médiocre, eau de lavage polluée.
            \item \textit{Contraste :} Bouteilles PET rigides, creuses, lourdes $\rightarrow$ flux régulier, triable par NIR, broyable sans enroulement, lavable efficacement.
        \end{itemize}
        \item \textbf{Masse pétrole économisée (80\% PET recyclé) :}
        \begin{align*}
            m_{\text{PET recyclé}} &= 660 \times 0,80 = 528\,\text{kg} \\
            \text{Pétrole économisé} &= 528 \times 2,5 = \boxed{1\,320\,\text{kg pétrole brut}}
        \end{align*}
    \end{enumerate}

    \item \textbf{Impacts écotoxicologiques (3 pts)}
    \begin{enumerate}
        \item \textbf{Fragmentation bouteille PET $\rightarrow$ microplastiques ($<5\,\text{mm}$) :}
        \begin{enumerate}
            \item \textbf{Physique :} Action mécanique (vagues, courant, abrasion sable/roches, choc thermique jour/nuit) $\rightarrow$ fissuration, éclatement en morceaux (mésoplastiques $5\text{--}25\,\text{mm}$).
            \item \textbf{Chimique :} Photodégradation UV (cassure chaînes polymères, oxydation, fragilisation) $\rightarrow$ poudre/flocons. Hydrolyse lente (esthers) à long terme.
            \item \textbf{Biologique :} Biofilm (bactéries, algues) $\rightarrow$ alourdissement (coulement), biodégradation de surface très lente (enzymes cutinases, PETases rares), fragmentation par ingestion/excrétion zooplancton/poissons.
        \end{enumerate}
        \item \textbf{Deux mécanismes d'impact sur physiologie poissons (\emph{Barbus}, \emph{Silure}) :}
        \begin{enumerate}
            \item \textbf{Fausse satiété / Obstruction :} Ingestion particules non nutritives $\rightarrow$ remplissage estomac/intestin $\rightarrow$ réduction ingestion proies réelles $\rightarrow$ perte condition, croissance ralentie, mortalité par occlusion.
            \item \textbf{Vecteur de polluants (Effet « Cheval de Troie ») :} Microplastiques adsorbent polluants hydrophobes (HAP, PCB, DDT, métaux) $\rightarrow$ ingestion $\rightarrow$ désorption dans tube digestif (surfactants biliaires, pH, température) $\rightarrow$ transfert contaminants vers tissus (foie, muscle, gonades) $\rightarrow$ toxicité hépatique, reprotoxique, immunodépression.
        \end{enumerate}
        \item \textbf{Risque santé humaine (Perturbateurs endocriniens + Bioamplification) :}
        \begin{itemize}
            \item Microplastiques + phtalates/BPA (additifs plastiques ou adsorbés) ingérés par proies (zooplancton, petits poissons).
            \item \textbf{Bioamplification :} Concentration croissante le long de la chaîne trophique (facteur biomagnification $>1$ pour molécules lipophiles persistantes).
            \item Homme (consommateur final poisson \emph{Silure}, \emph{Barbus}) $\rightarrow$ exposition chronique faibles doses $\rightarrow$ perturbation système endocrinien (œstrogénique, anti-androgénique) $\rightarrow$ infertilité, malformations génitales, puberté précoce, cancers hormone-dépendants, troubles métaboliques (diabète, obésité). Risque accru fœtus/enfants.
        \end{itemize}
    \end{enumerate}

    \item \textbf{Solutions intégrées et réglementation (4 pts)}
    \begin{enumerate}
        \item \textbf{Valorisation matière sachets PE (films fins) :} \textbf{Fabrication de « pavés plastiques » / « bois plastique » (ex: startup \textit{Kemit Ecology}, \textit{Plastic Waste to Wealth} à Douala/Yaoundé).}
        \textbf{Principe :} Collecte $\rightarrow$ Lavage/Séchage $\rightarrow$ Broyage $\rightarrow$ Mélange avec sable (60-70\% sable / 30-40\% PE fondu) $\rightarrow$ Extrusion/Compression en moules $\rightarrow$ Pavés de voirie, dalles, poteaux, bancs publics. 
        \textit{Avantages :} Pas de tri fin requis (tolère saletés, multicouches), pas de granulation fine, produit local durable (imperméable, imputrescible, moins cher que ciment), valorise sable local.
        \item \textbf{Trois mesures préventives (Autorités Locale/MINEPDED) :}
        \begin{enumerate}
            \item \textbf{Réglementaire + Contrôle :} \textbf{Application stricte Décret 2012} (interdiction sachets $<60\,\mu m$) : contrôles inopinés marchés/usines, saisie/destruction stocks illégaux, amendes lourdes (ex: 5 millions FCFA), fermeture boutiques récidivistes. \textit{Efficacité :} Coupe l'offre à la source, force substitution (papier, tissu, panier, sachets biodégradables certifiés $>60\,\mu m$).
            \item \textbf{Économique :} \textbf{Consigne obligatoire sur bouteilles PET/Canettes/Verre} (ex: 50 FCFA/bouteille) gérée par éco-organisme (REP - Responsabilité Élargie du Producteur). \textit{Efficacité :} Incitation financière forte retour consommateur $\rightarrow$ taux collecte $>80\%$, propreté rues/rivières, gisement propre pour recycleurs.
            \item \textbf{Éducative/Infrastructure :} \textbf{Déploiement bacs tri sélectif (Vert/Jaune/Bleu) + collecte séparée porte-à-porte/points d'apport} dans tous quartiers Douala, avec calendrier ramassage différencié (organiques 3x/sem, recyclables 1x/sem). Campagne continue « Mon quartier trie ». \textit{Efficacité :} Change habitudes, fournit gisement trié au CVD, réduit mélange à la source.
        \end{enumerate}
        \item \textbf{Affiche sensibilisation piroguiers/riverains Wouri :}
        \begin{center}
            \fbox{\parbox{0.9\linewidth}{\centering
            \textbf{\LARGE WOURI PROPRE = NOTRE SANTÉ, NOTRE FIERTÉ} \\[4pt]
            \textit{(Slogan : 7 mots)}
            }}
        \end{center}
        \textbf{Trois messages clés :}
        \begin{enumerate}
            \item \textbf{Bien commun :} « Le Wouri n'est pas une poubelle : chaque sachet jeté tue le poisson qui nourrit ta famille et bouche le chenal qui fait vivre ton commerce. »
            \item \textbf{Santé publique :} « Plastiques brûlés ou pourris = fumées toxiques et eau empoisonnée : cancers, asthma, diarrhées pour tes enfants. Protège-les, trie tes déchets. »
            \item \textbf{Action concrète :} « Ramène tes bouteilles vides au point de collecte \textit{[Nom lieu]} : 50 FCFA la bouteille, l'argent dans ta poche, le fleuve propre pour tous. »
        \end{enumerate}
    \end{enumerate}
\end{enumerate}
\textbf{Points de vigilance Probatoire :}
\begin{itemize}
    \item Calculs : attention unités (kg/g, tonnes/kg). Arrondis raisonnables.
    \item Fragmentation : distinguer facteurs physiques, chimiques, biologiques (3 mécanismes distincts).
    \item Microplastiques : citer « adsorption polluants » + « fausse satiété » + « transfert chaîne alimentaire ».
    \item Perturbateurs endocriniens : lier \textbf{adsorption sur MP} $\rightarrow$ \textbf{bioamplification} $\rightarrow$ \textbf{effet cocktail faibles doses} $\rightarrow$ \textbf{santé humaine}.
    \item Solutions locales : citer exemples réels camerounais (pavés plastiques, startups) = points bonus.
    \item Mesures préventives : 1 Réglementaire (coercitive), 1 Économique (incitative), 1 Éducative/Infrastructure (systémique). Justifier chacune.
    \item Slogan : court, percutant, positif, inclusif (« notre », « nous »).
\end{itemize}
\end{corrige}

\begin{corrige}[Exercice 11 — Sujet type Probatoire : Compostage et agriculture urbaine durable (8 pts)]
\textbf{Barème indicatif :} Q1 (2 pts) + Q2 (3 pts) + Q3 (3 pts) = 8 pts.

\begin{enumerate}
    \item \textbf{Équilibrage de la recette (2 pts)}
    \begin{enumerate}
        \item \textbf{Calcul C/N global initial (pondéré par masses de Matière Sèche) :}
        \textbf{Étape 1 : Masse MS pour 100 kg mélange initial (base 100 kg frais).}
        \begin{align*}
            \text{Déchets cuisine (60 kg frais)} &: MS = 60 \times (1 - 0,85) = 60 \times 0,15 = \mathbf{9\,kg\,MS} \\
            \text{Fientes poules (40 kg frais)} &: MS = 40 \times (1 - 0,60) = 40 \times 0,40 = \mathbf{16\,kg\,MS} \\
            \text{Total MS} &= 25\,kg
        \end{align*}
        \textbf{Étape 2 : Masse Carbone (C) et Azote (N) dans la MS (hypothèse C $\approx$ 45\% MS pour les deux, standard).}
        \begin{align*}
            \text{Cuisine : } C &= 9 \times 0,45 = 4,05\,kg \quad ; \quad N = \frac{C}{15} = \frac{4,05}{15} = \mathbf{0,27\,kg} \\
            \text{Fientes : } C &= 16 \times 0,45 = 7,20\,kg \quad ; \quad N = \frac{C}{10} = \frac{7,20}{10} = \mathbf{0,72\,kg} \\
            \text{Total C} &= 11,25\,kg \quad ; \quad \text{Total N} = 0,99\,kg
        \end{align*}
        \textbf{Étape 3 : C/N global.}
        \[ \text{C/N}_{\text{global}} = \frac{11,25}{0,99} = \boxed{11,4} \]
        \textit{(Méthode simplifiée pondérée par MS demandée : $\frac{9\times15 + 16\times10}{9+16} = \frac{295}{25} = 11,8$. Les deux donnent $\approx 11\text{--}12$.)}
        
        \textbf{Interprétation :} C/N $\approx 11\text{--}12$ \textbf{très inférieur à l'optimal (25--35)}. Mélange \textbf{trop riche en azote} (fientes 40\% apportent beaucoup de MS et N). Risque : perte d'ammoniac (\ce{NH3}) $\rightarrow$ odeurs, pollution air, perte valeur fertilisante, ralentissement dégradation carbone (manque C pour énergie micro-organismes). \textbf{Pas équilibré.}

        \item \textbf{Modification recette : ajout sciure (C/N=400, MS=90\%) pour C/N cible = 30.}
        \textbf{Notations :} $x$ = kg sciure fraîche ajoutée pour 100 kg mélange initial.
        \begin{itemize}
            \item Sciure : $MS_s = 0,90x$ ; $C_s \approx 0,45 \times MS_s = 0,405x$ ; $N_s = C_s / 400 = 0,0010125x$.
            \item Mélange initial (100 kg) : $C_i = 11,25\,kg$ ; $N_i = 0,99\,kg$ (valeurs ci-dessus).
        \end{itemize}
        \textbf{Équation C/N final = 30 :}
        \[ \frac{C_i + C_s}{N_i + N_s} = 30 \quad \Rightarrow \quad \frac{11,25 + 0,405x}{0,99 + 0,0010125x} = 30 \]
        \[ 11,25 + 0,405x = 29,7 + 0,030375x \]
        \[ 0,374625x = 18,45 \]
        \[ x = \frac{18,45}{0,374625} \approx \boxed{49,2\,kg\,\text{sciure fraîche}} \]
        \textbf{Pourcentage sciure / mélange initial total :}
        \[ \frac{49,2}{100} \times 100 = \boxed{49,2\%} \]
        \textit{(Nouveau mélange total = 149,2 kg ; sciure = 33\% du nouveau total).}
    \end{enumerate}

    \item \textbf{Suivi processus et qualité (3 pts)}
    \begin{enumerate}
        \item \textbf{Trois paramètres à surveiller + cibles optimales :}
        \begin{center}
            \begin{tabular}{|l|l|l|}
                \hline
                \textbf{Paramètre} & \textbf{Cible optimale} & \textbf{Rôle / Justification} \\ \hline
                \textbf{Température} (cœur andin) & \textbf{$55\text{--}65^\circ\text{C}$} (phase thermophile $\ge 3$ jours consécutifs) & Hygiénisation (pathogènes, graines adventices) ; activité microbienne max. $>70^\circ\text{C}$ = mort micro-organismes. \\ \hline
                \textbf{Humidité} & \textbf{50--60\%} (au toucher : goutte à peine entre doigts) & Milieu aqueux pour enzymes/bactéries. $<40\%$ = arrêt activité ; $>65\%$ = anaérobiose (odeurs, lixiviats). \\ \hline
                \textbf{Aération / O$_2$} (ou CO$_2$ émis) & \textbf{O$_2 > 10\%$ dans air interstitiel} (retournement si $<5\%$ ou T$>70^\circ\text{C}$) & Maintien conditions aérobies. Évite méthanisation, \ce{H2S}, \ce{NH3}. \\ \hline
                \textbf{pH} (optionnel 3ème) & \textbf{6,5--8,0} (légèrement acide $\rightarrow$ neutre) & Activité bactérienne/fongique. $<5,5$ = inhibition bactéries ; $>8,5$ = pertes \ce{NH3}. \\ \hline
            \end{tabular}
        \end{center}
        \item \textbf{Interprétation C/N final = 12 :}
        \begin{itemize}
            \item \textbf{Stabilité :} C/N $< 15\text{--}20$ indique une \textbf{bonne stabilité} (matière organique humifiée, peu de C facilement biodégradable résiduel). Le compost est mature.
            \item \textbf{Risque « faim d'azote » :} \textbf{Faible à nul.} C/N $= 12$ signifie que la matière organique est déjà riche en N relatif. Lors de l'épandage, la minéralisation libérera de l'azote minéral (\ce{NH4+}, \ce{NO3-}) \textbf{disponible pour les plantes} (effet fertilisant immédiat), sans immobilisation nette d'azote du sol par les micro-organismes (qui aurait lieu si C/N $> 25$).
        \end{itemize}
        \item \textbf{Conformité hygiénique (\emph{E. coli}) :}
        \begin{itemize}
            \item Analyse : $500\,\text{UFC/g}$ frai. Norme NF U 44-051 : $< 1000\,\text{UFC/g}$.
            \item \textbf{Conclusion :} \textbf{CONFORME} ($500 < 1000$).
            \item \textbf{Rôle phase thermophile ($>55^\circ\text{C}$) :} La chaleur détruit les pathogènes (\emph{E. coli}, \emph{Salmonella}, œufs vers, graines adventices) par dénaturation protéines/enzymes. La norme exige maintien $>55^\circ\text{C}$ pendant \textbf{au moins 3 jours consécutifs} (ou 7 jours selon référentiels) dans la masse de l'andin, garanti par retournements réguliers. C'est la \textbf{garantie d'hygiénisation} du processus.
        \end{itemize}
    \end{enumerate}

    \item \textbf{Modèle économique et impact agricole (3 pts)}
    \begin{enumerate}
        \item \textbf{Production mensuelle compost mature :}
        \begin{align*}
            \text{Entrée journalière} &= 2\,\text{t} = 2\,000\,\text{kg} \\
            \text{Perte masse totale} &= 55\% \quad \Rightarrow \quad \text{Rendement} = 45\% \\
            \text{Sortie journalière} &= 2\,000 \times 0,45 = 900\,\text{kg/j} = 0,9\,\text{t/j} \\
            \text{Production mensuelle (30 j)} &= 0,9 \times 30 = \boxed{27\,\text{tonnes}}
        \end{align*}
        \item \textbf{Résultat net mensuel :}
        \begin{align*}
            \text{Recettes} &= 27\,000\,\text{kg} \times 100\,\text{FCFA/kg} = 2\,700\,000\,\text{FCFA} \\
            \text{Charges} &= 1\,500\,000\,\text{FCFA} \\
            \text{Résultat net} &= 2\,700\,000 - 1\,500\,000 = \boxed{+1\,200\,000\,\text{FCFA (Bénéfice)}}
        \end{align*}
        \item \textbf{Apports épandage 10 t/ha (base fraîche) :}
        \begin{itemize}
            \item Compost frais : $10\,\text{t/ha} = 10\,000\,\text{kg/ha}$.
            \item MS ($60\%$) : $6\,000\,\text{kg MS/ha}$.
            \item \textbf{Matière Organique (MO) :} $45\%$ MS $\rightarrow 6\,000 \times 0,45 = \boxed{2\,700\,\text{kg MO/ha} = 2,7\,\text{t MO/ha}}$.
            \item \textbf{Azote total (N) :} $2\%$ MS $\rightarrow 6\,000 \times 0,02 = \boxed{120\,\text{kg N/ha}}$ (unités d'azote).
        \end{itemize}
        \textbf{Intérêt agronomique sols bas-fonds (sableux, pauvres humus) :}
        \begin{itemize}
            \item \textbf{MO (2,7 t/ha) :} Augmente \textbf{Capacité d'Échange Cationique (CEC)} (sable CEC $\approx 0$), rétention eau/éléments nutritifs, structure (agrégation), activité biologique (vers, microbes). Lutte contre lessivage.
            \item \textbf{N (120 kg/ha) :} Couverture besoins azotés majeure partie cycle culturale (ex: maïs $\approx 150\text{--}200$ kg N/ha ; légumes-feuilles $\approx 100\text{--}150$). Libération progressive (minéralisation) = engrais « à libération lente », moins lessivable que urée.
            \item \textbf{Synergie :} MO tamponne pH, complexe métaux lourds éventuels. Amendement + Engrais organique complet. Rentabilité pour maraîcher (économie engrais chimiques).
        \end{itemize}
    \end{enumerate}
\end{enumerate}
\textbf{Points de vigilance Probatoire :}
\begin{itemize}
    \item C/N initial : pondération par \textbf{masse MS} (pas masse fraîche). Hypothèse %C dans MS à expliciter (standard 45-50\%).
    \item Calcul sciure : bien distinguer masse fraîche / MS. Équation algébrique claire.
    \item Paramètres suivi : Température (hygiénisation), Humidité (vie microbienne), Aération (aérobie). Cibles chiffrées précises.
    \item C/N final 12 : « Stable » + « Pas faim azote » (au contraire, effet fertilisant N). Ne pas confondre avec C/N initial.
    \item Hygiénisation : Seuil norme ($1000$ UFC/g) vs mesure ($500$). Rôle T$>55^\circ\text{C}$ + durée (3-7 j) + retournements.
    \item Bilan économique : Unités cohérentes (t $\rightarrow$ kg pour prix FCFA/kg). Bénéfice = Recettes - Charges.
    \item Apports : Base fraîche $\rightarrow$ MS $\rightarrow$ MO et N sur MS. Unités : t MO/ha, kg N/ha (UF/ha).
    \item Discussion agronomique : Lier propriétés sable (faible CEC, lessivage) $\rightarrow$ apports MO/N du compost.
\end{itemize}
\end{corrige}

\newpage

\coursec{Fiche bilan}

\begin{popfigure}
\begin{tikzpicture}[mindmap, grow cyclic, every node/.style={concept, font=\small}, concept color=popblue, text=white,
  level 1/.append style={level distance=3.4cm, sibling angle=60, font=\small},
  level 2/.append style={level distance=2.1cm, sibling angle=45, font=\scriptsize}]
\node[concept, font=\bfseries] {Déchets : transformation et recyclage}
  child[concept color=poporange] { node {Typologie}
    child { node {organiques} }
    child { node {plastiques} }
    child { node {dangereux} }
  }
  child[concept color=poppink] { node {Impacts}
    child { node {pollution sol, eau, air} }
    child { node {santé : paludisme, choléra} }
  }
  child[concept color=poppurple] { node {Règle des 3R}
    child { node {Réduire} }
    child { node {Réutiliser} }
    child { node {Recycler} }
  }
  child[concept color=popgreen] { node {Tri}
    child { node {à la source} }
    child { node {selon la nature} }
  }
  child[concept color=popgold] { node {Recyclage}
    child { node {filières : verre, métal, plastique} }
    child { node {transformation} }
  }
  child[concept color=popblue!70!popgreen] { node {Valorisation locale}
    child { node {compostage} }
    child { node {biogaz (méthanisation)} }
  };
\end{tikzpicture}
\legende{Carte mentale de la leçon : les six idées forces autour de la gestion durable des déchets.}
\end{popfigure}

\begin{bilan}
\textbf{Définitions clés}
\begin{itemize}
  \item \cle{Déchet} : tout résidu d'une activité humaine (production, consommation) dont on se débarrasse.
  \item \cle{Déchet organique} : déchet biodégradable d'origine végétale ou animale (épluchures, restes de repas, fumier).
  \item \cle{Déchet dangereux} : déchet toxique, inflammable ou infectieux (piles, médicaments périmés, déchets hospitaliers, huiles usées).
  \item \cle{Tri} : séparation des déchets à la source selon leur nature pour orienter chacun vers la bonne filière.
  \item \cle{Recyclage} : transformation d'un déchet en matière première réutilisable.
  \item \cle{Compostage} : décomposition aérobie contrôlée des déchets organiques produisant du \cle{compost} (engrais naturel).
  \item \cle{Méthanisation} : décomposition anaérobie des déchets organiques produisant du \cle{biogaz} (méthane \ce{CH4}) et un digestat fertilisant.
\end{itemize}

\textbf{Connaissances à retenir par cœur}
\begin{center}
\begin{tabularx}{\linewidth}{|l|X|}
\hline
\rowcolor{popblueL} \textbf{Notion} & \textbf{À savoir} \\
\hline
Règle des 3R & \textbf{Réduire} à la source, \textbf{Réutiliser}, \textbf{Recycler} \\
\hline
Durées de dégradation & sac plastique : $\sim$ 450 ans ; mégot : 1 à 2 ans ; pelure de banane : quelques mois ; verre : quasi indégradable \\
\hline
Impacts majeurs & pollution des sols et des nappes phréatiques, émissions de gaz à effet de serre, gîtes larvaires (moustiques), maladies hydriques (choléra, typhoïde) \\
\hline
Compostage (aérobie) & déchets organiques + \ce{O2} $\rightarrow$ compost + \ce{CO2} + \ce{H2O} + chaleur \\
\hline
Méthanisation (anaérobie) & déchets organiques $\rightarrow$ biogaz (\ce{CH4} + \ce{CO2}) + digestat \\
\hline
Exemples camerounais & compost maraîcher (périurbain de Yaoundé et Douala), biodigesteurs en zone rurale, récupération informelle des plastiques et métaux \\
\hline
\end{tabularx}
\end{center}

\textbf{Méthodes express}
\begin{itemize}
  \item \emph{Trier un déchet} : identifier sa nature (organique / recyclable / dangereux / non valorisable) $\rightarrow$ l'orienter vers la bonne filière (compost, recyclage, filière spécifique, enfouissement).
  \item \emph{Proposer une valorisation} : déchet organique $\rightarrow$ compostage ou biogaz ; plastique, verre, métal, papier $\rightarrow$ recyclage ; déchet dangereux $\rightarrow$ collecte spécialisée, jamais en nature ni au feu.
  \item \emph{Sensibiliser} : message simple + exemple local + geste concret (poubelle de tri, composteur familial, refus des sachets plastiques).
\end{itemize}
\end{bilan}

\begin{aretenir}[Checklist avant le Probatoire]
\begin{itemize}
  \item[$\square$] Je sais définir un déchet et citer les trois grandes catégories (organiques, plastiques, dangereux) avec un exemple chacune.
  \item[$\square$] Je connais les durées de dégradation repères (plastique, verre, matière organique).
  \item[$\square$] Je sais expliquer deux impacts des déchets sur l'environnement (sols, eaux, air, climat).
  \item[$\square$] Je sais relier les dépôts sauvages de déchets à des maladies précises (paludisme, choléra, typhoïde).
  \item[$\square$] Je connais la règle des 3R et je sais donner un exemple concret pour chaque R.
  \item[$\square$] Je sais trier un déchet donné et justifier la filière choisie.
  \item[$\square$] Je distingue recyclage, réutilisation et valorisation.
  \item[$\square$] Je sais décrire le compostage : conditions (aérobie), étapes, produit obtenu.
  \item[$\square$] Je sais décrire la méthanisation : conditions (anaérobie), biogaz produit (\ce{CH4}), usages (cuisson, éclairage).
  \item[$\square$] Je connais au moins deux exemples locaux de valorisation des déchets au Cameroun.
  \item[$\square$] Je sais proposer des actions de sensibilisation adaptées à mon quartier ou mon village.
  \item[$\square$] Je rédige mes réponses avec le vocabulaire scientifique exact (biodégradable, anaérobie, digestat, filière).
\end{itemize}
\end{aretenir}

\findecours

\end{document}
